\documentclass[11pt]{article}
\usepackage[a4paper,margin=22mm]{geometry}
\usepackage[no-math]{fontspec}
\setmainfont{DejaVu Serif}
\usepackage{amsmath,amssymb,amsthm,mathtools,graphicx,xcolor,hyperref,xurl,xspace,ifthen}
\setlength{\emergencystretch}{4em}
\SetMathAlphabet{\mathbf}{normal}{OT1}{cmr}{bx}{n}
\SetMathAlphabet{\mathbf}{bold}{OT1}{cmr}{bx}{n}
\newtheorem{theo}{Theorem}[section]
\newtheorem{lemm}[theo]{Lemma}
\newtheorem{theorem}[theo]{Theorem}
\newtheorem*{remas*}{Remarks}
\newtheorem{remas}[theo]{Remarks}
\newtheorem*{exam*}{Example}
\newtheorem{ques}[theo]{Question}
\providecommand{\og}{«\nobreakspace}
\providecommand{\fg}{\nobreakspace»}
\newtheorem{lem}[theo]{Lemma}
\newtheorem{lemma}[theo]{Lemma}
\newtheorem{prop}[theo]{Proposition}
\newtheorem{coro}[theo]{Corollary}
\newtheorem{corol}[theo]{Corollary}
\newtheorem{conj}[theo]{Conjecture}
\newtheorem{Proposition}[theo]{Proposition}
\newtheorem{theoreme}[theo]{Theorem}
\theoremstyle{definition}
\newtheorem{defi}[theo]{Definition}
\newtheorem{exam}[theo]{Example}
\newtheorem{exem}[theo]{Example}
\newtheorem{remark}[theo]{Remark}
\newtheorem{example}[theo]{Example}
\newtheorem{rema}[theo]{Remark}
\newtheorem{nota}[theo]{Notation}
\newtheorem*{rema*}{Remark}
\newtheorem*{defi*}{Definition}
\newtheorem*{theo*}{Theorem}
\newtheorem*{lemm*}{Lemma}
\newtheorem*{prop*}{Proposition}
\newcommand{\enoncename}{}
\newtheorem{corpusenonce}[theo]{\enoncename}
\newenvironment{enonce}[2][]{\renewcommand{\enoncename}{#2}\begin{corpusenonce}}{\end{corpusenonce}}
\newenvironment{enonce*}[2][]{\par\medskip\noindent\textbf{#2.}\itshape}{\par\medskip}
\newenvironment{proofc}[1][\proofname]{\begin{proof}[#1]}{\end{proof}}
\newcommand{\Mc}{\mathcal{M}}
\newcommand{\pointir}{.\ ---\ }
\providecommand{\mathof}[1]{\mathrm{#1}}

\numberwithin{equation}{section}
\numberwithin{figure}{section}
\numberwithin{table}{section}
\def\endappendix{}
%~Mouliné par MaN_auto v.0.21.6 2019-08-22 13:20:55


\usepackage{mathtools}
\usepackage{amssymb}
%\usepackage{enumerate}
%\usepackage{psfrag}
\usepackage{mathrsfs}
% Original dutchcal mathcal font is unavailable; standard Computer Modern mathcal rendering, unchanged formula tokens.

\DeclareMathOperator{\supp}{Supp}
\DeclareMathOperator{\hyb}{hyb}
\DeclareMathOperator{\Gal}{Gal}
\DeclareMathOperator{\conv}{Conv}
\DeclareMathOperator{\tr}{tr}
\DeclareMathOperator{\id}{id}
\DeclareMathOperator{\red}{red}
\DeclareMathOperator{\length}{length}
\DeclareMathOperator{\ord}{ord}
\DeclareMathOperator{\Lm}{Lim}
\DeclareMathOperator{\sph}{sph}
\DeclareMathOperator{\diam}{diam}
\DeclareMathOperator{\SL}{SL}
\DeclareMathOperator{\PGL}{PGL}
\DeclareMathOperator{\Aff}{Aff}
\DeclareMathOperator{\Aut}{Aut}
\DeclareMathOperator{\res}{res}
\DeclareMathOperator{\diag}{diag}

\newcommand{\cc}{\mathbb{C}}
\newcommand{\re}{\mathbb{R}}
\newcommand{\dD}{\mathbb{D}}
\newcommand{\nn}{\mathbb{N}}
\newcommand{\pp}{\mathbb{P}}
\newcommand{\hh}{\mathbb{H}}
\newcommand{\prob}{\mathsf{P}}
\newcommand{\phyb}{{\mathbb{P}^1_{\mathrm{hyb}}}}
\newcommand{\Rep}{\mathrm{rep}}
\newcommand{\Att}{\mathrm{att}}

\newcommand{\cA}{\mathcal{A}}
\newcommand{\cD}{\mathcal{D}}
\newcommand{\sm}{\mathcal{m}}
\newcommand{\cO}{\mathcal{O}}

\newcommand{\cdt}{\cc(\!(t)\!)}

\newcommand{\mm}{\mathbb{M}}
\newcommand{\A}{{\mathbb A}}
\newcommand{\R}{{\mathbb R}}
\newcommand{\Z}{{\mathbb Z}}
\newcommand{\Q}{{\mathbb Q}}
\newcommand{\B}{\mathbb{B}}
\newcommand{\D}{\mathbb{D}}
\newcommand{\N}{{\mathbb N}}

\newcommand{\sE}{\mathsf{E}}

\newcommand{\an}{\mathrm{an}}

\newcommand{\e}{\varepsilon}
\newcommand{\cv}{\rightarrow}
\newcommand{\cvf}{\rightharpoonup}

\newcommand{\om}{\Omega}
\newcommand{\set}[1]{\left\{#1\right\}}
\newcommand{\norm}[1]{\left\Vert#1\right\Vert}
\newcommand{\abs}[1]{\left\vert#1\right\vert}
\newcommand{\cd}{{\cc^2}}
\newcommand{\pu}{{\mathbb{P}^1}}
\newcommand{\PP}{\mathbb{P}}

\newcommand{\berk}{{\mathbb{P}^{1, \mathrm{an}}_k}}
\newcommand{\berkl}{\mathbb P^{1,\an}_{\Laurent}}
\newcommand{\Laurent}{\mathbb{L}}

\newcommand{\rest}[1]{ \arrowvert_{#1}}
\newcommand{\m}{\textbf{M}}
\newcommand{\unsur}[1]{\frac{1}{#1}}

\newcommand{\lrpar}[1]{\left(#1\right)}
\newcommand{\la}{\lambda}

\newcommand{\att}{\mathrm{att}}
\newcommand{\rep}{\mathrm{rep}}
\newcommand{\na}{\mathsf{na}}

\newcommand{\SLC}{\SL(2,\mathbb C)}
\newcommand{\SLk}{\SL(2,k)}
\newcommand{\PGLk}{\PGL(2, k)}

\newcommand{\inv}{^{-1}}

\newtheorem{mthm}{Theorem}
\renewcommand{\themthm}{\Alph{mthm}}

\newcommand*{\dd}{\mathrm{d}}
\newcommand*{\du}{\dd u}
\newcommand*{\dm}{\dd m}
\newcommand*{\dmu}{\dd \mu}
\newcommand*{\dnu}{\dd \nu}
\newcommand*{\dsm}{\dd \sm}
\newcommand*{\dprob}{\dd \prob}

\newenvironment{step}{\begin{proof}}{\let\qed\relax \end{proof}}









\graphicspath{{./figures/}}

\newcommand*{\mk}{\mkern -1mu}
\newcommand*{\Mk}{\mkern -2mu}
\newcommand*{\mK}{\mkern 1mu}
\newcommand*{\MK}{\mkern 2mu}

\hypersetup{urlcolor=purple, linkcolor=blue, citecolor=red}

\newcommand*{\romanenumi}{\renewcommand*{\theenumi}{\roman{enumi}}}
\newcommand*{\Romanenumi}{\renewcommand*{\theenumi}{\Roman{enumi}}}
\newcommand*{\alphenumi}{\renewcommand*{\theenumi}{\alph{enumi}}}
\newcommand*{\Alphenumi}{\renewcommand*{\theenumi}{\Alph{enumi}}}
\let\oldforall\forall
\renewcommand*{\forall}{\mathrel{\oldforall}}


\begin{document}
\section*{\raggedright 2691: Unique atomless stationary Berkovich probability measures for finitely generated non-elementary random matrix groups over the complex Laurent-series field}
\noindent\textbf{Author:} Romain Dujardin; Charles Favre\par\noindent\textbf{Source:} Degenerations of SL(2, C) representations and Lyapunov exponents. Annales Henri Lebesgue 2 (2019), publisher version, DOI 10.5802/ahl.24. \url{https://ahl.centre-mersenne.org/item/AHL_2019__2__515_0/}
\par\noindent\textbf{License:} CC BY 4.0, \url{https://creativecommons.org/licenses/by/4.0/}. The original publisher PDF has the explicit license notice. See the saved source/license records.
\par\noindent\textbf{Editorial material note (not author text):} The selected problem is the stationary-measure part of Theorem 3.1 over the complex Laurent-series field with its t-adic norm, for countable support generating a finitely generated non-elementary group in the exact proximal/strongly-irreducible sense of Definition 2.17. No moment assumption is imposed for this selected conclusion. The complete author body retains all local projective/tree, compact-limit-set, martingale/Dirac-contraction, uniformly-small-ball, atomlessness and classical-point concentration arguments, including the deferred Lemma 3.8 proof. The Lyapunov, Furstenberg-formula, degeneration and hybrid-space conclusions are unselected context; named Berkovich/probability foundations remain references. All proof text and formulas below are editable author TeX. Mathematical wording, language and proof order are retained. Only the article class, title matter and bibliography presentation are adapted for independent compilation; no proof has been generated or repaired.
\par\noindent\textbf{Editorial difficulty:} H2: The author obtains boundary stationarity without local compactness by combining Berkovich compactification, hyperbolic contractions, a martingale limit argument and concentration on classical points. This non-Archimedean extension of random matrix theory exceeds ordinary qualifying probability and geometry.
\medskip\hrule\medskip
\noindent\textbf{Author statement, complete proof and local supporting text follow in source order.}
\section*{Introduction}

Let $G$ be a finitely generated group, endowed with a probability measure $\sm$, satisfying the following two conditions:
\begin{enumerate}\renewcommand{\theenumi}{A\arabic{enumi}}
\item \label{A1} $\supp(\sm)$ generates $G$;
\item \label{A2} $ \int \length(g) \dsm(g)<\infty$.
\end{enumerate}
In a few occasions we shall also require the following stronger moment condition:
\begin{enumerate}\renewcommand{\theenumi}{A\arabic{enumi}${}^+$}\setcounter{enumi}{1}
\item \label{A2+} there exists $\delta >0$ such that $ \int (\length(g))^{1+\delta} \dsm(g)<\infty$.
\end{enumerate}
In \eqref{A2} and~\eqref{A2+}, $\length({\,\cdot\,})$ denotes the word-length relative to some fixed, unspecified, finite symmetric set of generators of $G$. It depends of course on the choice of generators but the moment conditions do not.

%\medskip

If $\rho:G \to \SLC$ is any representation, the random walk on $G$ induced by $\sm$ gives rise through $\rho$ to a random product of matrices in $\SLC$. If $\norm{{\,\cdot\,}}$ denotes any matrix norm on $\SLC$, then under the moment condition \eqref{A2} we can define the \emph{Lyapunov exponent} $\chi = \chi(G, \sm, \rho)$ by the formula
\begin{equation}
\label{def:lyapunov exponent intro}
\!\chi \mk=\mk \lim_{n\to\infty} \unsur{n} \mk\int\mk \log\norm{\rho(g_n\cdots g_1)} \dsm(g_1) \cdots \dsm(g_n) \mk=\mk \lim_{n\to\infty} \unsur{n} \mk\int\mk \log\norm{\rho(g)} \dsm^n(g),
\end{equation}
where $\sm^n$ is the image of $\sm^{\otimes n}$ under the $n$-fold product map $(g_1, \ldots, g_n)\mapsto g_n\cdots g_1$. Observe that the limit exists because if we choose the matrix norm to be submultiplicative then the sequence of integrals is subadditive. The Lyapunov exponent is the most basic dynamical invariant associated to the random product of matrices, and its properties have been the object of intense research since the seminal work of Furstenberg~\cite{furstenberg} in the 1960's.

It is quite customary that $(\rho, \sm)$ depends on certain parameters, in which case the dependence of $\chi$ as a function of the parameters becomes an interesting problem. One famous instance of this problem, motivated by statistical physics, is the study of discrete Schr\"odinger operators, which involves random products of matrices of the form
\[
\begin{pmatrix} E-v & -1 \\ 1 & 0
\end{pmatrix},
\]
where $v$ is a real random variable and $E$ (the energy) is a real or complex parameter
(see e.g.~\cite{bougerol-lacroix}).

%\medskip

We are interested in the situation where the representation depends holomorphically on a complex parameter $t$, that is, we consider a family of representations $(\rho_t)$ such that for every $g\in G$, the function $t\mapsto \rho_t(g)$ is holomorphic. Then the Lyapunov exponent defines a function $\chi(t) = \chi(G, \sm, \rho_t)$ on the parameter space. Recall that a representation $\rho: G \to \SLC$ is said to be \emph{non-elementary}\footnote{Another common terminology is ``strongly irreducible and proximal''.} if there exist two elements $g_1, g_2 \in G$ such that $\rho(g_1)$ and $\rho(g_2)$ are both hyperbolic (i.e. their eigenvalues have modulus $\neq 1$) and have no common eigenvectors. A celebrated result due to Furstenberg~\cite{furstenberg} (see also~\cite{furstenberg-kifer}) asserts that if $t_0$ is such that $\rho_{t_0}$ is non-elementary, then under the assumptions \eqref{A1}--\eqref{A2}, $\chi(t)$ is positive and continuous at $t_0$. H\"older continuity can also be derived under stronger moment conditions (see~\cite{lepage-holder}). And it was recently proved by Bocker and Viana~\cite{bocker-viana} that if $\sm$ is finitely supported then $\chi$ is continuous at elementary representations as well.

It is a classical observation that $t\mapsto \chi(t)$ is subharmonic. In~\cite{kleinbif, bers1} Deroin and the first named author have studied the complex analytic properties of the Lyapunov exponent function in relation to the classical bifurcation/stability theory of Kleinian groups (devised by Bers, Maskit, Sullivan, etc.) and established that the harmonicity of $\chi$ over some domain is equivalent to the structural stability of the corresponding family of representations.

%\medskip

Our purpose in this paper is to study the asymptotic properties of $\chi(t)$ in non-compact families of representations with ``algebraic behavior'' at infinity. To be specific, we consider a holomorphic family $(\rho_t)_{t\in \dD^*}$ of representations of $G$ into $\SLC$, parameterized by the punctured unit disk, and such that the function $t\mapsto \rho_t(g)$ extends meromorphically through the origin for every $g$. An obvious but crucial observation is that this data is equivalent to that of a single representation with values in $\SL(2, \mm)$ where $\mm$ is the ring of holomorphic functions on the punctured unit disk with meromorphic extension through the origin.

We will show that the behavior of $\chi(t)$ at $t\to 0$ is controlled by a quantity which can be interpreted as a non-Archimedean Lyapunov exponent associated to the family $(\rho_t)_{t\in \dD^*}$. To make sense of this statement, observe first that $\mm$ may be viewed as a subring of the field of formal Laurent series\footnote{Beware that in some of the references cited in our bibliography, $\Laurent$ denotes the completion of the algebraic closure of $\cdt$.} $\Laurent:=\cdt$, which is a complete metrized field when endowed with the $t$-adic norm $\abs{f}_\na = \exp(-\ord_{t=0}(f))$. A representation $\rho : G \to \SL(2,\mm)$ thus canonically yields a representation $\rho_\na : G \to \SL(2,\Laurent)$ and exactly as in~\eqref{def:lyapunov exponent intro} we define $\chi_\na$ to be the Lyapunov exponent of the representation $\rho_\na$, where the matrix norm $\norm{\,\cdot\,}$ is now associated to the $t$-adic absolute value on $\Laurent$.


We are now in a position to state our main result.

\begin{mthm}\label{Thm:lyapunov na}
Let $(G, \sm)$ be a finitely generated group endowed with a measure satisfying~\eqref{A1} and \eqref{A2+}, and let $\rho: G\to \SL(2, \mm)$ be any representation. Then
\begin{equation}
\label{eq:lyapunov na}
\unsur{{\log\abs{t}\inv}}\chi(t) \longrightarrow \chi_\na \text{ as } t\cv 0.
\end{equation}
\end{mthm}


To illustrate the result in a simple case, consider a random product of Schr\"odinger matrices of the form
\begin{equation}
\label{eq:schrodinger matrix}
\begin{pmatrix} {\unsur{t}(E-v)} & -1 \\ 1 & 0
\end{pmatrix},
\end{equation}
where $v$ is a bounded random variable, $E$ is fixed, and $t\to 0$. Then it is easy to show in this case that $\chi(t)\sim \log\abs{t}\inv$ as $t\to 0$. One reason for this is that the pole structure of the $n$-fold product of such matrices is explicit and easy to describe: such a product will be of the form $\unsur{t^n}\lrpar{
\begin{smallmatrix} O(1) & O(t) \\ O(t) & O(t^2)
\end{smallmatrix}}$.
Avron, Craig and Simon~\cite{avron-craig-simon} gave in this situation a refined asymptotics of $\chi(t)$ at the order $o(1)$ (and a conjectural asymptotics at the order $O(t^{2})$).

For a general random product of matrices with meromorphic coefficients, the poles can add up or cancel in a rather subtle way, and the non-Archimedean formalism allows to deal efficiently with this issue. The idea of using non-Archimedean representations to describe the degenerations of $\SLC$ representations is now classical and was pioneered by Culler and Shalen~\cite{culler-shalen}. One main input of the present work is the incorporation of this technique into the theory of random matrix products.

%\medskip
Let us explain the strategy of the proof of Theorem~\ref{Thm:lyapunov na}. Since $\Laurent$ is a metrized field it makes sense to talk about hyperbolic elements in $\SL(2,\Laurent)$ so we can define natural notions of elementary and non-elementary subgroups. We refer to Section~\ref{sec:subgroups} for a thorough discussion of these concepts.

The proof of Theorem~\ref{Thm:lyapunov na} splits into two quite different parts according to the elementary or non-elementary nature of $\rho_\na$. The easier case is when $\rho_\na$ is elementary: then either the family is holomorphic at the origin (after conjugation by a suitable meromorphic family of M\"obius transformations and possibly taking a branched 2-cover of the base) or $\rho_t$ is elementary for all $t\in \D^*$. In the former case, $\chi_\na = 0$ and it follows from Furstenberg's theory that $\chi(t) = O(1)$ so we are done. In the latter case, up to meromorphic conjugacy, the image of $\rho_t$ in $\SLC$ is either upper triangular or lies in an index $2$ extension of the diagonal subgroup. The result then follows from a careful application of the law of large numbers. The details of the arguments are explained in Section~\ref{sec:elem}. Note that this is the only place where we use the stronger integrability condition~\eqref{A2+}.

%\medskip

Let us now assume that $\rho_\na$ is non-elementary. A first observation is that $\rho_t$ is then non-elementary for small enough $t$ (see Lemma~\ref{lem:specialization NE}; by rescaling we may assume that this holds for $\abs{t}<1$) so we may apply Furstenberg's theory to analyze the Lyapunov exponent $\chi(t)$. The main step of the proof of the continuity of $\chi$ in the classical setting is the study of the Markov chain on $\pu$ induced by the image measure $\rho_*\sm$ in $\SL(2,\cc)$. More specifically, the measure $\mu_t = (\rho_t)_* \sm$ acts by convolution on the set of probability measures on the Riemann sphere by $\nu \mapsto \mu_t \ast \nu$, where $\mu_t \ast \nu = \int \gamma_*\nu \; \dmu_t(\gamma)$. A \emph{stationary measure} is by definition a fixed point of this action. A fundamental result is that when $\rho_t$ is non-elementary, there is a \emph{unique} stationary probability measure $\nu_t$. Furthermore, the Lyapunov exponent $\chi(t)$ is positive and can be expressed by an explicit formula involving $\nu_t$:
\begin{equation}\label{eq:fur}
\chi(t) = \int \log \frac{\| \gamma \cdot v\|}{\|v \|} \, \dmu_t(\gamma)\, \dnu_t(v)~.
\end{equation}
The continuity of $\chi$ at non-elementary representations immediately follows: since $\nu_t$ is unique it varies continuously with $t$, and~\eqref{eq:fur} implies the result.

%\medskip

The first step of the proof consists in extending these results to random products in $\SL(2, \Laurent)$. We actually work over an arbitrary complete metrized field $k$ and show in Section~\ref{sec:furstenberg} how to generalize the above results to any non-elementary representation $\rho: G \to \SLk$. Of particular interest to us is the fact that the representation $\rho_\na$ admits a unique stationary measure $\nu_\na$ which lives on the Berkovich analytification $\PP^{1,\an}_{\Laurent}$ of the projective line and for which a non-Archimedean analog of~\eqref{eq:fur} holds. In particular we obtain the positivity of the non-Archimedean Lyapunov exponent.

Let us point out that this positivity also follows from the recent work of Maher and Tiozzo~\cite{maher-tiozzo} on random walks on groups of isometries of non-proper Gromov hyperbolic spaces. Maher and Tiozzo also discuss stationary measures, however, they work in a compactification which is a priori hard to relate to the Berkovich space.

%\medskip

From this point, two different paths lead to the main theorem. Both of them deal with the asymptotic properties of the stationary measure $\nu_t$ and can be seen as ways to imitate the Furstenberg argument for the continuity of the Lyapunov exponent.

%\medskip

The first method belongs to complex geometry and is described in Section~\ref{sec:non elem}. It relies on a correspondence between certain finite subsets of $\PP^{1,\an}_{\Laurent}$ and \emph{models} of $\D \times \PP^1_\cc $. Here by model we mean a complex surface $Y$ endowed with a birational map $\pi: Y \to \D \times \PP^1_\cc$ which is a biholomorphism over $\D^*\times \PP^1_\cc$.

For $t\neq 0$ we denote by $\nu_{Y_t}$ the pull-back of $\nu_t$ on $Y$, which should be understood as ``the measure $\nu_t$ viewed on the model $Y$''. We obtain the following result.

\begin{mthm}\label{Thm:atomic}
Let $(G, \sm)$ be a finitely generated group endowed with a measure satisfying (A1) and let $\rho: G\to \SL(2, \mm)$ be a non-elementary representation.

Then for every model $Y\cv \dD\times \PP^1_\cc$, the canonical family of stationary probability measures $\nu_{Y_t}$ converges as $t\to 0$ to a purely atomic measure $\nu_{Y}$.

Furthermore if $\nu_\na$ denotes the unique stationary probability measure on $\PP^{1,\an}_{\Laurent}$, then $\nu_Y = (\res_Y)_*\nu_\na$ is the residual measure of $\nu_\na$ on $Y$.
\end{mthm}

Here the residue map $\res_Y$ is a canonical anti-continuous\footnote{i.e. the pull back of an open set (resp. a closed set) set is closed (resp. open).} map from $\PP^{1,\an}_{\Laurent}$ to the $\cc$-scheme $\pi^{-1}(\{ 0 \} \times \PP^1_\cc)$. In particular the push-forward of a \emph{non-atomic} measure on $\PP^{1,\an}_{\Laurent}$ is an \emph{atomic} measure on $Y$. We refer to Section~\ref{subs:reduction} for a detailed discussion on this map. The asymptotics~\eqref{eq:lyapunov na} of the Lyapunov exponent in Theorem~\ref{Thm:lyapunov na} then follows from an analysis of~\eqref{eq:fur} as $t\cv 0$ in a carefully chosen family of models.

%\medskip

The second approach to Theorem~\ref{Thm:lyapunov na} relies on the notion of \emph{hybrid (Berkovich) space}. This is a topological space which allows to give a precise meaning to the intuitive idea that $\nu_t$ ``converges'' to $\nu_\na$ as $t\to 0$, and derive the asymptotics~\eqref{eq:lyapunov na} directly from this weak convergence.

This space was first constructed by Berkovich in~\cite{berko}. It has been recently realized by Boucksom and Jonsson in~\cite{boucksom-jonsson} and by the second named author in~\cite{degeneration} that it is well-adapted to the description of the limiting behavior of families of measures such as the $(\nu_t)_{t\in \dD^*}$. Concretely, $\phyb$ is a compact topological space endowed with a continuous surjective map\footnote{The choice of the value $1/e$ for the radius is convenient, but any other would do.} $p_{\hyb}: \phyb \to \overline \dD_{1/e}$ such that $p_{\hyb}^{-1}(0)$ can be identified with the non-Archimedean analytic space $\PP^{1,\an}_{\Laurent}$ while $p_{\hyb}$ is a trivial topological fibration over $\overline \dD_{1/e}^*$ with $\PP^1_\cc$ fibers. More precisely there exists a canonical homeomorphism $\psi\colon \overline\dD_{1/e}^* \times \PP^1_\cc \rightarrow p_{\hyb}^{-1}(\overline\dD_{1/e}^*) $ such that $p_{\hyb} \circ \psi$ is the projection onto the first factor. Likewise we denote by $\psi_\na$ the canonical identification $\PP^{1,\an}_{\Laurent} \to p_{\hyb}\inv(\set{0})$.

A key point in the construction of $\phyb$ is that its topology is designed so that for every $g\in G$ the function
\[
(t,v) \longmapsto \frac1{\log|t|^{-1}} \, \log \frac{\|\rho_t(g)\cdot v\|}{\| v \|}
\]
extends \emph{continuously} to the hybrid space for $t=0$. Theorem~\ref{Thm:lyapunov na} hence follows immediately from:

\begin{mthm}\label{Thm:hybrid}
Let $(G, \sm)$ be a finitely generated group endowed with a measure satisfying~\eqref{A1} and $\rho: G\to \SL(2, \mm)$ be a non-elementary representation.

Then in $\phyb$, we have that $ (\psi_{t})_*(\nu_t) \longrightarrow (\psi_\na)_*\nu_{\na}$ as $t\cv 0$ in the weak topology of measures, where $\nu_t$ is the unique stationary probability measure under $\mu_t$, and $\psi_t({\,\cdot\,}) =\psi(t,\cdot\,)$.
\end{mthm}

We discuss the construction of the hybrid space and prove Theorem~\ref{Thm:hybrid} in Section~\ref{sec:hybrid}.

%\medskip
%\begin{center}$\diamond$
%\end{center}
%\medskip


As for~\cite{kleinbif, bers1}, this work was prompted by analogous results in the context of iteration of rational mappings, in accordance with the celebrated \emph{Sullivan dictionary}. Consider a holomorphic family of rational maps $(f_t)_{t\in \D^*}$ of degree $d\ge2$ that extends meromorphically through $0$, and denote by $\mu_{f_t}$ their measure of maximal entropy. In this context, the analogue of~\eqref{eq:fur} is a formula for the blow-up of the Lyapunov exponent of $\mu_{f_t}$ which follows from the work of DeMarco~\cite{demarco03,demarco16}, Theorem~\ref{Thm:lyapunov na} is due to~\cite{degeneration} in arbitrary dimension, and Theorem~\ref{Thm:atomic} was proven by DeMarco and Faber~\cite{demarco-faber}. It is particularly interesting to note that proving the convergence of the measures $\mu_{f_t}$ relies on pluripotential theory and the interpretation of $\mu_{f_t}$ as the Monge--Amp\`ere measure of a suitable metrization on an ample line bundle, whereas in our case it follows from the uniqueness of the stationary measure.

%\medskip

Our work raises several natural open questions.
\begin{enumerate}
\item Is it possible to estimate the error term $\chi(t) - \chi_\na \log\abs{t}^{-1}$? The answer is easy when $\chi(t)$ is harmonic in a punctured neighborhood of the origin, in which case one obtains an expansion of the form
\[
\chi(t) = \chi_\na \log\abs{t}^{-1} + C +o(1)
\]
for some constant $C$ (this situation happens e.g. in~\eqref{eq:schrodinger matrix}).
In the general case, continuity holds under appropriate moment assumptions if $\rho_\na$ is elementary (see Section~\ref{subs:continuity elementary}). However in the nonelementary case our method seems to produce errors of magnitude $\e\log\abs{t}^{-1}$ (see Section~\ref{subs:asymptotics}) so new ideas have to be developed.

\item Can our results be extended to higher dimensions? Random matrix products in arbitrary dimension over an Archimedean or a local field are well-understood. Some of our arguments should carry over to the study of the extremal Lyapunov exponents, even if we use the hyperbolic structure of $\PP^{1, \an}_{\Laurent}$ at some places. Note however that even the Oseledets theorem does not seem to have received much attention over arbitrary metrized fields.
\end{enumerate}
%\medskip
%\begin{center}$\diamond$
%\end{center}
%\medskip

The plan of the paper is as follows. In Section~\ref{sec:prelim} we recall some basics on Berkovich theory. In Section~\ref{sec:subgroups} we classify subgroups of $\PGL(2,k)$ for an arbitrary complete metrized field $k$. In particular we define the notion of non-elementary subgroup and classify elementary ones. Part of this material follows from the classical theory of groups acting on trees. In Section~\ref{sec:furstenberg} we develop the non-Archimedean Furstenberg theory. The complex geometric proof of Theorem~\ref{Thm:lyapunov na} is given in Sections~\ref{sec:non elem} and~\ref{sec:elem}. In Section~\ref{sec:non elem} we prove Theorem~\ref{Thm:atomic} and deal with the non-elementary case, while the elementary case is treated in Section~\ref{sec:elem}. The hybrid formalism and Theorem~\ref{Thm:hybrid} are explained in Section~\ref{sec:hybrid}.

\subsection*{Acknowledgement}

We are grateful to Bertrand Deroin for many useful conversations.


\section{The Berkovich projective line}\label{sec:prelim}

In this section, we collect some basic facts on the Berkovich analytification of the projective line over a complete non-trivially metrized field $(k, |{\,\cdot\,}|)$. We write $k^* = k \setminus \{ 0 \}$ and let $|k^*|$ be the value group of $k$. Observe that we do not assume $k$ to be algebraically closed (since we apply these results to $k = \cdt $ later on) which leads to a few subtleties. The reader is referred to~\cite{Ber,temkin} for a general discussion on Berkovich spaces, and to~\cite{baker-rumely,jonsson} for a detailed description of the Berkovich projective line.

\subsection{Analytification of the projective line}

We denote by $\PP^1_k$ the projective line over a field $k$, viewed as an algebraic variety, endowed with its Zariski topology, and by $\PP^1(k)$ its set of $k$-points which is in bijection with $k \cup \{ \infty\}$. When $k = \cc$, we often simply denote by $\pu = \pu(\cc)$ the Riemann sphere, that is, the complex projective line with its usual structure of compact complex manifold\footnote{Note that formally $\pu$ can be viewed as the analytification $\PP^{1,\an}_\cc$ of the variety $\PP^1_\cc$.}.

In the remainder of this section, we suppose that $(k,|{\,\cdot\,}|)$ is a complete metrized \emph{non-Archimedean} field. We also assume that the norm on $k$ is non-trivial, so in particular $k$ is infinite. We denote by $\PP^{1,\an}_k$ the Berkovich analytification of $\PP^1_k$ which is a compact topological space endowed with a structural sheaf of analytic functions. Only its topological structure will be used in this paper, and we refer the interested reader to~\cite{berko} for the description of the structural sheaf.

The Berkovich space $\PP^{1,\an}_k$ is defined as follows. The analytification of the affine line $\A^{1,\an}_k$ is the space of all multiplicative semi-norms on $k[Z]$ whose restriction to $k$ coincides with $|{\,\cdot\,}|$, endowed with the topology of pointwise convergence. Given a point $x\in \A^{1,\an}_k$ and a polynomial $P\in k[Z]$, the value of the semi-norm defined by $x$ on $P$ is usually denoted by $\mathopen{|}P\mathclose{|}_x \in \R_+$. It is also customary to denote it by $|P(x)|$, the reason for this notation should be clear from the classification of semi-norms below. The Gau{\ss} norm $\sum_i a_i Z^i \mapsto \max |a_i|$ defines a point denoted by $x_{\mathrm{g}}$, and referred to as the Gau{\ss} point.

The Berkovich projective line can be defined as a topological space to be the one-point compactification of $\A^{1,\an}_k$ so that we write $\PP^{1,\an}_k = \A^{1,\an}_k \cup \{\infty\}$. More formally it is obtained by gluing two copies of $\A^{1,\an}_k$ in a standard way using the transition map $z\mapsto z^{-1}$ on the punctured affine line $(\A^1)^{*,\an}_k$.

A \emph{rigid} point in $\PP^{1,\an}_k$ is a point defined by a multiplicative semi-norm having a non-trivial kernel. For any point $z$ lying in a finite extension of $k$, the semi-norm $\abs{{\,\cdot\,}}_z$ defined by $P \mapsto |P(z)|$ is a rigid point in $\A^{1,\an}_k$. The induced map yields a canonical bijection between closed points of the $k$-scheme $\PP^1_k$ and rigid points in $\PP^{1,\an}_k$. In particular $\pp^1(k)$ naturally embeds as a set of rigid points in $\PP^{1,\an}_k$, and in the following we simply view $\pp^1(k)$ as a subset of $\PP^{1,\an}_k$.

The Berkovich projective line $\PP^{1,\an}_k$ is an $\R$-tree in the sense that it is uniquely pathwise connected, see~\cite[\S 2]{jonsson} for precise definitions. In particular for any pair of points $x,y\in \PP^{1,\an}_k$ there is a well-defined segment $[x,y]$. Recall that the convex hull of a subset $F$ in an $\R$-tree is the smallest connected set $\conv (F)$ which contains $F$, that is the union of all segments $[x,y]$ with $x, y \in F$. Any point in $\PP^{1,\an}_k$ admits a well defined projection to a closed convex subset.

In this paper, by \emph{measure} on $\PP^{1,\an}_k$ we mean a Radon measure, that is a Borel measure which is inner regular, or equivalently a bounded linear functional on the vector space of continuous functions on $\PP^{1,\an}_k$ endowed with the sup norm.

Using the tree structure, one can show the following result (see~\cite[Lem.~7.15]{valtree}).

\begin{lemm}\label{lem:support}
The support of any Radon measure in $\PP^{1,\an}_k$ is compact and metrizable.
\end{lemm}

\subsection{Balls in the projective line and semi-norms}\label{subsec:balls}

We still assume that the norm on $k$ is non-Archimedean and non-trivial. Closed and open balls in $k$ (of radius $R \in \R_+$) are defined as usual:
\[
\overline{B}(z_0,R)= \{ z\in k, \, | z -z_0| \le R\} \quad\text{ and}\quad {B}(z_0,R)= \{ z\in k, \, | z -z_0| < R\}\,.
\]
By definition, a ball in $\PP^1(k)$ is either a ball in $k$ or the complement of a ball in $k$.

Any closed (or open) ball $B\subsetneq\PP^1(k)$ determines a point $x_B\in \PP^{1,\an}_k$. When $B$ or its complement is a singleton $\{ z\}$, this point $x_B $ is the rigid point attached to $z$. When $B$ is a (open or closed) ball of finite radius in $k$, we let $x_B$ be the point in $\A^{1,\an}_k$ corresponding to the semi-norm $|P(x_B)| := \sup_B |P|$. Otherwise the complement of $B$ is a ball of finite radius in $k$, and we set $|P(x_B)| := \sup_{k\setminus B} |P|$. Observe that a closed ball $x_B$ is rigid iff its diameter is zero, and that the Gau{\ss} point is equal to $x_{\overline{B}(0,1)}$.

\begin{rema}
When $|k^*|$ is dense in $\R_+$, we have $x_{\overline{B}(z_0,R)}=x_{B(z_0,R)}$ for all $z_0\in k$ and $R\in \R_+$. Otherwise $k$ is discretely valued, $|k^*| = r^\Z$ for some $r>1$, and we have $x_{B(z_0,r^n)}= x_{\overline{B}(z_0,r^{n-1})}$.
\end{rema}

\subsection{The spherical metric}

Let $(k,|{\,\cdot\,}|)$ be any non-Archimedean complete metrized field. We can endow $\PP^1(k)$ with the spherical metric
\begin{equation}\label{eq:dsph}
d_{\sph} ([z_0:z_1], [w_0:w_1]) = \frac{|z_0w_1 - z_1w_0|}{\max\{ |z_0|, |z_1|\} \, \max\{ |w_0|, |w_1|\}}~.
\end{equation}
which satisfies the ultrametric inequality. The spherical diameter of $\PP^1(k)$ is equal to $1$. For any $z\in \PP^1(k)$ and $r\le 1$ we define closed and open \emph{spherical} balls
\[
\overline{B}^{\sph}(z,r) = \{ d_{\sph} ({\,\cdot\,}, z) \le r\} \quad\text{ and}\quad B^{\sph}(z,r) = \{ d_{\sph} ({\,\cdot\,}, z) < r\}\,.
\]
Observe that for all $r\ge1$, $\overline{B}^{\sph}(z,r) = \PP^1(k)$. A spherical ball is either $\PP^1(k)$ or a ball in $\PP^1(k)$ in the sense of the previous section. Conversely a ball in $\PP^1(k)$ is either a spherical ball or the complement of a spherical ball.

\subsection{The hyperbolic space (\texorpdfstring{$k$}{k} algebraically closed)} \label{subs:hyperbolic alg closed}

Let $(k,|{\,\cdot\,}|)$ be any \emph{algebraically closed} non-Archimedean complete metrized field, and let us describe the geometry of $\PP^{1,\an}_k$ under this assumption. First observe that rigid points are in bijection with balls of zero diameter. Following the Berkovich terminology we say that points corresponding to balls of diameter $\diam(B) \in |k^*|$ (resp. $\diam(B) \notin |k^*|$) are of type $2$ (resp. of type $3$). Rigid points are said to be of type $1$.

More generally, points in $\PP^{1,\an}_k$ are in bijection with (equivalence classes of) decreasing sequences of balls, see~\cite[Thm.~1.2]{baker-rumely}. When $k$ is spherically complete, that is, every decreasing intersection of balls is non-empty, then $\PP^{1,\an}_k$ consists of type $1$, $2$ or $3$ points. In general, there may exist a fourth type of points, associated with decreasing sequences of balls with empty intersection. All these types of points (whenever non-empty) yield dense subsets of $\PP^{1,\an}_k$.

Types of points relate with the tree structure as follows. Type $1$ and $4$ points are precisely the ones at which the $\R$-tree $\PP^{1,\an}_k$ has only one branch, that is, they are endpoints of the tree. Type $2$ points are branching points (i.e. $\PP^{1,\an}_k$ has at least three branches at these points) and type $3$ points are regular points (i.e. $\PP^{1,\an}_k$ has exactly two branches at these points).

The hyperbolic space $\hh_k$ is by definition the complement of the set of rigid points in $\PP^{1,\an}_k$, i.e. $\hh_k = \PP^{1,\an}_k \setminus \PP^{1}(k)$. It is a proper subtree of $\PP^{1,\an}_k \setminus \PP^1(k)$ which contains no rigid point and is neither open nor closed.

To describe the structure of $\hh_k$, for any $r\in \R_+$ introduce the semi-norm $x_r\in \A^{1,\an}_k$ defined by
\[
|P(x_r)| = \max \left\{ |a_n| r^n, \, a_n \neq 0 \right\}
\]
where $P(Z) = \sum a_n Z^n$. In particular $x_r = x_{\bar B(0,r)}$. Let $\hh_k^\circ$ be the orbit under $\PGL(2,k)$ of the ray $\{ x_r, \, r\in \R_+^*\}$. Then $\hh_k^\circ$ is a dense subtree of $\hh_k$, and $\hh_k \setminus \hh_k^\circ$ coincides with the set of type 4 points in $\PP^{1,\an}_k$.

By~\cite[Prop.~2.30]{baker-rumely}, one can endow $\hh_k^\circ$ with a unique $\PGL(2,k)$-invariant metric such that
\[
d_{\hh} (x_{r_1}, x_{r_2}) = \log \left|\frac{r_1}{r_2}\right|
\]
for any $r_1\ge r_2 >0$. Observe that
\[
d_{\hh} (x_{B_1}, x_{B_2}) = \log \left(\frac{\diam(B_2)}{\diam(B_1)}\right),
\]
for any closed two balls $B_1 \subset B_2 \subset k$ (the diameter is relative to the metric induced by the norm $|{\,\cdot\,}|$).

A proof of the next result can be found in~\cite[Prop.~2.29]{baker-rumely}.

\begin{lemm}
The metric defined above on $\hh_k^\circ$ extends to a distance on $\hh_k$, which makes $(\hh_k, d_{\hh})$ a complete metric $\R$-tree upon which $\PGLk$ acts by isometries.
\end{lemm}

Recall that by metric $\R$-tree we mean that for any pair of distinct points $x,y$ there exists a unique isometric embedding $\gamma : [0,d_{\hh}(x,y)] \to \hh_k$ such that $\gamma(0) = x$, $\gamma(1) =y$, and $d_{\hh}(\gamma(t), \gamma(t')) = |t -t'|$.

\subsection{Field extensions}\label{subs:field extensions}

We no longer assume $k$ is algebraically closed. Let $K/k$ be any complete field extension. The inclusion $k[Z] \subset K[Z]$ yields by restriction a canonical surjective and continuous map $\pi_{K/k}\colon \PP^{1,\an}_K \to \PP^{1,\an}_k$, and the Galois group $\Gal(K/k)$ acts continuously on $\PP^{1,\an}_K$.

Let $\bar{k}^a$ be the completion of an algebraic closure of $k$. Then $\PP^{1,\an}_{k}$ is homeomorphic to the quotient of $\PP^{1,\an}_{\bar{k}^a}$ by $\Gal(\bar{k}^a/k)$, see~\cite[Cor.~1.3.6]{berko}. The group $\Gal(\bar{k}^a/k)$ preserves the types of points in $\PP^{1,\an}_{\bar{k}^a}$, so that we may define the type of a point $x \in \PP^{1,\an}_k$ as the type of any of its preimage by $\pi_{\bar{k}^a/k}$ in $\PP^{1,\an}_{\bar{k}^a}$. Note that since the field extension $\bar{k}^a/k$ is not algebraic in general, it may happen that some type 1 points in $\PP^{1,\an}_k$ have trivial kernel hence are not rigid (this phenomenon occurs when $k$ is the field of Laurent series over any field).

By~\cite[Prop.~2.15]{baker-rumely}, the natural action of $\PGL(2, k)$ on $\PP^{1,\an}_{\bar{k}^a}$ preserves the types of points so the same holds for its action on $\PP^{1,\an}_k$.

\begin{prop}\label{prop:orbit gauss}
The following assertions are equivalent.
\begin{enumerate}
\item \label{1.4.1} The point $x\in \PP^{1,\an}_k$ belongs to the orbit of the Gau{\ss} point $x_{\mathrm{g}}$ under the action of $\PGL(2,k)$ (in particular it is of type 2).
\item \label{1.4.2} There exists $z\in k$ and $r\in |k^*|$ such that $x= x_{\bar{B}(z,r)}$.
\end{enumerate}
\end{prop}

\begin{proof}
Pick $z\in k$ and assume $r\in \abs{k^*}$ so that there exists $y\in k^*$ with $r= |y|$. The image of $\bar{B}(0,1)$ by the affine map $Z \mapsto yZ + z$ is equal to $\bar{B}(z,r)$ hence $\eqref{1.4.2}\Rightarrow\eqref{1.4.1}$. Conversely any element in $\PGL(2,k)$ can be decomposed as a product of affine maps and the inversion $\Phi(Z) := 1/Z$. Thus we conclude that $\eqref{1.4.1}\Rightarrow\eqref{1.4.2}$ by observing that $\Phi(x_{\bar{B}(z,r)})= x_{\bar{B}(z^{-1},r/|z|^2)}$ if $r \le |z|$, $\Phi(x_{\bar{B}(0,r)})= x_{\bar{B}(0,r^{-1})}$ and $A(x_{\bar{B}(z,r)}) = x_{\bar{B}(az+b, |a|r)}$ if $A(Z) = aZ +b$.
\end{proof}

\begin{prop}\label{prop:projection}
Suppose $x,y, z$ belong to the orbit of $x_{\mathrm{g}}$ under $\PGL(2,k)$. Then the projection of $z$ on $[x,y]$ also belongs to the orbit of $x_{\mathrm{g}}$ under $\PGL(2,k)$.
\end{prop}

\begin{proof}
We can normalize the situation so that $x = x_{\mathrm{g}} = x_{\bar B(0,1)}$ and $y = x_{\bar B(0,r)}$ for some $1<r\in |k^*|$. Let $z = x_{\bar B(a, r')}$. If $\bar B(a, r')$ is disjoint from $\bar B(0,1)$ or contains it then the projection is the Gau{\ss} point and we are done. Otherwise $\bar B(a, r') \subset \bar B(0,1)$ with $\abs{a}\leq 1$ and $r'<1$. If $\bar B(a, r') \subset \bar B(0,r)$ the projection equals $\bar B(0,r)$ and again we are done. The remaining case is when $\abs{a} >r $, in which case the projection is $\bar B(0,\abs{a})$ and we conclude by Proposition~\ref{prop:orbit gauss}.
\end{proof}

\subsection{The hyperbolic space (\texorpdfstring{$k$}{k} arbitrary)} \label{subs:hyperbolic arbitrary}

The Galois group $\Gal(\bar{k}^a/k)$ acts on $\bar{k}^a$ by isometries, so the diameter of balls is $\Gal(\bar{k}^a/k)$-invariant. As a consequence the action of $\Gal(\bar{k}^a/k)$ on $(\hh_{\bar{k}^a},d_\hh)$ is also isometric. Let $\widetilde{\hh}_k\subset\hh_{\bar{k}^a}$ be the set of fixed points of this action and recall that $\conv({\,\cdot\,})$ denotes the convex hull.

\begin{lemm}\label{lem:fixx}
The set of fixed points of the action of $\Gal(\bar{k}^a/k)$ on $\PP^{1,\an}_{\bar{k}^a}$ is $\overline{\conv(\PP^1(k))}$.
\end{lemm}

\begin{proof}
Denoting by $\mathcal{F}$ this set of fixed points, it is clear that $\mathcal{F}$ contains $\PP^1(k)$. Since the Galois action preserves the tree structure we infer that $\mathcal{F}\supset \mathrm{Conv}(\PP^1(k))$ and since it is weakly continuous we finally deduce that $\mathcal{F}$ contains $\overline{\mathrm{Conv}(\PP^1(k))}$.

Suppose by contradiction that there exists a point $x\in \mathcal{F}$ which does not belong to $\overline{\mathrm{Conv}(\PP^1(k))}$. Since type $2$ points are dense on any ray in the tree $\PP^{1,\an}_{\bar{k}^a}$, we may suppose $x =x_{\bar{B}(z,r)}$ for some $z\in \bar{k}^a$ and $r \in |k^*|^{\Q}$. It is enough to show that $\bar{B}(z,r)$ contains a point of $k$. Indeed in this case we get that $x \in\mathrm{Conv}(\PP^1(k))$ which contradicts our assumption.

To show this, first note that algebraic points over $k$ are dense in $\bar{k}^a$, hence we may assume that $z$ is algebraic over $k$. Let $P$ be its minimal polynomial, and suppose first that its degree $d$ is prime to the characteristic of $k$. The point $x$ is fixed by $\Gal(\bar{k}^a/k)$ hence so does $\bar{B}(z,r)$, so this ball contains all the roots $z=z_1, \ldots, z_d$ of $P$ (repeated according to their multiplicity if needed). In particular letting $z^* = \frac1{d} (z_1+ \cdots + z_d)$ we have that $|z^* - z| \le r$ and $z^*\in k$ and we are done.

When $d$ is not prime to the characteristic, we modify this argument as follows. Fix $a\in k^*$ such that $ |a| \cdot |z|^{d+1} \ll \abs{P(x)}$ (recall that $\abs{P(x)} = \sup_B\abs{P}$) and consider the polynomial $\widetilde{P}(X)= a X^{d+1} + P(X)$. By definition of $P$ we have $|\widetilde{P}(z)| = |a| \cdot |z|^{d+1}$, and on the other hand $|\widetilde{P}(x)| \ge |P(x)| $. This classically implies the existence of a root of $\widetilde{P}$ in $\bar{B}(z,r)$, so $\widetilde P$ is a polynomial in $k[X]$ of degree $d+1$ with a root in $\bar{B}(z,r)$ and we can apply the previous argument to $\widetilde P$.
\end{proof}

By the previous lemma we have that in $\PP^{1,\an}_{\bar{k}^a}$, $\widetilde{\hh}_k:= \overline{\conv(\PP^1(k))}\setminus \PP^1(k)$. In particular $(\widetilde{\hh}_{k},d_\hh)$ is a complete metric $\R$-tree by Section~\ref{subs:hyperbolic alg closed}.

Since $\PP^{1,\an}_{k}$ is homeomorphic to the quotient of $\PP^{1,\an}_{\bar{k}^a}$ by $\Gal(\bar{k}^a/k)$, the restriction map $\pi_{\bar{k}^a/k}$ induces a homeomorphism from $\widetilde{\hh}_k$ onto its image, which we denote by $\hh_k$ and call the hyperbolic space\footnote{When $k$ is a $p$-adic field, $\hh_k$ is \emph{not} the Drinfeld upper half-plane which is equal as a set to $\PP^{1,\an}_k \setminus \PP^1(k)$. } over $k$. We endow it with the metric $d_\hh$ making $\pi_{\bar{k}^a/k}\colon (\widetilde{\hh}_k,d_\hh)\to (\hh_k,d_\hh)$ an isometry.

The following proposition summarizes the properties of $\hh_k$ obtained so far.

\begin{prop}
The hyperbolic space $\hh_k$ is the complement of $\PP^1(k)$ in the closure of the convex hull of $\PP^1(k)$ in $\PP^{1,\an}_{k}$, that is $\hh_k:= \overline{\conv(\PP^1(k))}\setminus \PP^1(k)\subset \PP^{1,\an}_k$. Endowed with the metric $d_\hh$, it is a complete metric $\R$-tree upon which $\PGLk$ acts by isometries.
\end{prop}

Let $K/k$ be any complete field extension, then there is a canonical $\PGL(2,k)$-equivariant continuous map $\sigma_{K/k} : \overline{\conv(\PP^1(k))} \to \PP^{1,\an}_K$ such that $\pi_{K/k} \circ \sigma_{K/k} = \id$ and which sends the point $x_{\bar{B}(0,r)}\in\PP^{1,\an}_k$ to the corresponding point $x_{\bar{B}(0,r)}\in \PP^{1,\an}_K$ for all $r\in \R_+$. A detailed discussion of this map can be found in~\cite{poineau}. The map $\sigma_{K/k}$ is injective hence induces a homeomorphism from $\overline{\conv(\PP^1(k))}$ in $ \PP^{1,\an}_k$ onto its image which is the closure of the convex hull of $ \PP^1(k)$ in $ \PP^{1,\an}_K$.

\begin{prop}\label{prop:quadratic}
For any pair $x, y$ of points in $ \hh_k$ lying in the orbit of $x_{\mathrm{g}}$ under $\PGL(2,k)$, there exists a quadratic extension $K/k$ and $g\in \PGL(2,K)$ such that $g\cdot x_{\mathrm{g}}\in \PP^{1,\an}_K$ is the middle point of the segment $[\sigma_{K/k}(x),\sigma_{K/k}(y)]$.
\end{prop}

\begin{proof}
Applying a suitable M\"obius transformation, we may assume that $x= x_{\mathrm{g}} =x_{\bar{B}(0,1)}$ and $y = x_{\bar{B}(0,r)}$ with $r\in |k^*|$. Fix $z\in k^*$ such that $|z| = r$, and pick any square root $z'$ of $z$. The middle point of the segment $[x,y]:= \{x_{B(0,t)}, \, t \in [1,r]\}$ is the type 2 point $x_{\bar B(0, \sqrt{r})}$ which lies in the orbit of $x_{\mathrm{g}}$ by $\PGL(2,k(z'))$. The assertion is proved with $K=k(z')$.
\end{proof}

\subsection{Balls and simple domains in \texorpdfstring{$\PP^{1, \an}_k$}{PP1, ank}}

For any $a\in k$ and any $r\geq 0$, we set
\begin{equation}\label{eq:ball1}
\bar{B}^\an(a, r) = \set{x\in \A^{1, \an}_k, \ \abs{Z-a}_x \leq r}
\end{equation}
and
\begin{equation}\label{eq:ball2}
B^\an(a, r) = \set{x\in \A^{1, \an}_k, \ \abs{Z-a}_x < r}.
\end{equation}
When no confusion can arise, we drop the ``an'' superscript.

A closed ball in the Berkovich projective line is a set of the form $\bar B^\an(a, r)$ or the complement of a set of the form $B^\an(a, r) $ in $\PP^{1, \an}_k$. One similarly defines open balls. Observe that any ball in $\PP^{1, \an}_k$ not containing $\infty$ is of the form~\eqref{eq:ball1} or~\eqref{eq:ball2}. A closed (resp. open) ball is a closed (resp. open) subset of $\PP^{1, \an}_k$ with one single boundary point. When the ball is $\bar B^\an(a, r) $ or $B^\an(a, r)$, or their complements, this boundary point is $x_{\bar{B}(a,r)}$.

When $k$ is algebraically closed, a closed (resp. open) ball in $\PP^{1, \an}_k$ is the convex hull of a closed (resp. open) ball in $\PP^1(k)$.


Following the terminology of~\cite{baker-rumely}, by a \emph{simple domain} we mean any open set $U\subsetneq \berk$ whose boundary is a (non-empty) finite set of type 2 points. Open balls are simple domains, and it follows from~\cite[Thm.~4.2.1]{berko} that simple domains form a basis for the topology of $\berk$.

The next result will play an important role in our approach to Theorem~\ref{Thm:lyapunov na}.

\begin{prop}\label{prop:subtree Hk}
Let $\nu$ be any probability measure on $\berk$ having no atom. Then for every $\e>0$ there exists a finite set $S$ of type 2 points such that every connected component $U$ of $\berk\setminus S$ satisfies $\nu(U)<\e$.
\end{prop}

\begin{proof}
Pick any $\e>0$. Since $\nu$ is a Radon measure, any point $x$ is included in a simple domain $U_x$ such that $\nu(U_x) \le \e$. By compactness, we may cover the support of $\nu$ by finitely many of these domains $U_{x_1}, \ldots, U_{x_n}$. Let $S$ be the union of all boundary points of $U_{x_i}$ for $i =1, \ldots,n$. Any connected component $U$ of $\berk\setminus S$ intersects one of the open sets say $U_{x_1}$. Since $U \cap \partial U_{x_1} = U\cap S = \emptyset$, it follows that $U\subset U_{x_1}$ and $\nu(U) \le \e$ as claimed.
\end{proof}


\section{Subgroups of \texorpdfstring{$\PGLk$}{PGLk}}\label{sec:subgroups}
%\ici

In this section $(k, \abs{{\,\cdot\,}})$ is an arbitrary non-trivially valued field that is complete and non-Archimedean and by $k^\circ$ its valuation ring. The case of most interest to us is $\Laurent:=\cdt$ which is a complete metrized field when endowed with the $t$-adic norm $\abs{f}_\na = \exp(-\ord_{t=0}(f))$.

We consider a subgroup $\Gamma\leq \PGLk = \Aut(\pp^1_k)$ and study the geometric properties of its action on the projective and Berkovich spaces. Much of this material is a reformulation in our context of classical results on groups acting on trees (see e.g.~\cite{culler-morgan, kapovich,otal}, and also~\cite{yang-wang} for related material).

Over the complex numbers the corresponding results are well-known (see~\cite{beardon}).

\subsection{Basics}

As in the complex setting, there is a morphism $\SLk\to\PGLk$, defined by associating a M\"obius transformation to a 2-by-2 matrix by the usual formula
\[
\left(
\begin{matrix} a & b \\ c & d
\end{matrix}\right) \longmapsto \lrpar{z\mapsto \frac{az+b}{cz+d}},
\]
whose kernel is $\set{\pm \id}$.

Beware that in general this morphism is not surjective. It is so when the field $k$ is algebraically closed, but not for instance in the case $k = \cdt$. The trouble is that for a general M\"obius transformation $\frac{az+b}{cz+d}$, the determinant $ad-bc$ need not be a square in $k$. Thus, if we denote by $\e$ the generator of the Galois group of the quadratic field extension $\cc(\!(t^{1/2})\!) / \cdt$ (i.e. $\e \cdot t^{1/2} = - t^{1/2}$), we have a surjective morphism from the subset of matrices $M \in \SL(2, \cc(\!(t^{1/2})\!))$ for which $ \e \cdot M = \pm M$ onto $\PGL(2,\cdt)$ whose kernel is again $\set{\pm \id}$. In other words, after a base change we can always lift a meromorphic family of M\"obius transformations to a family of matrices in $\SL(2, \mm)$. The same phenomenon happens for triangular matrices over $k$ and the affine group $\Aff_k$.

Working with matrices is often more convenient for calculations, and when no confusion can arise, we simply identify $\gamma\in \SLk$ with the corresponding M\"obius transformation, denoted by $z\mapsto \gamma(z)$.

%\medskip
An element of $\PGLk$ induces an automorphism of $\berk$ preserving $\pu(k)$ and $\hh_k$. Recall that it preserves the types of points and acts by isometries on $(\hh_k,d_\hh)$.
%\smallskip

For $A = \lrpar{\begin{smallmatrix} a & b \\ c & d
\end{smallmatrix}} \in \SLk$ we denote by
$\norm{A} = \max\lrpar{\abs{a}, \abs{b}, \abs{c}, \abs{d}}$, which in the ultrametric case is the matrix norm associated to the sup norm on $ k^2$.

When $\gamma \in \PGL(2,k)$ as explained above there exists a quadratic extension $K/k$ and $A \in \SL(2,K)$ inducing $\gamma$ on $\PP^1(k)$, and we set $\| \gamma \| := \| A\|$. This is well-defined since $K$ carries a unique complete norm whose restriction to $k$ is $|{\,\cdot\,}|$ and $A$ is defined up to multiplication by $\pm \id$. Likewise, we define $\abs{\tr(\gamma)} := \abs{\tr(A)}$.

\begin{prop}
For all $\gamma$ and $\gamma' $ in $\SLk$, we have that:
\begin{enumerate}%\romanenumi%[{(i)}]
\item \label{2.1.1} $\norm{\gamma}\geq 1$,
\item \label{2.1.2} $\norm{\gamma \gamma'}\leq \norm{\gamma}\cdot\norm{\gamma'}$,
\item \label{2.1.3} $\norm{\gamma} = \norm{\gamma\inv}$,
\item \label{2.1.4} if furthermore $\gamma\in \SL(2,k^\circ)$, then $\gamma$ induces an isometry of $(\PP^1(k),d_{\sph})$.
\end{enumerate}
\end{prop}

The proof is left to the reader (note that~\eqref{2.1.1} follows from the ultrametric property of the absolute value).


\subsection{Classification of elements in \texorpdfstring{$\PGL(2,k)$}{PGL(2,k)}}

\begin{prop}\label{prop:claselt}
Let $\gamma\in \PGL(2,k)$, $\gamma\neq \id$. Then exactly one of the following holds:
\begin{itemize}
\item $\mathopen{|}\tr(\gamma)| >1$: then $\gamma$ is diagonalizable over $k$ and has one attracting (resp. repelling) fixed point $x_\att \in \PP^1(k)$ (resp. $x_\rep \in \PP^1(k)$). Furthermore for every $x\neq x_\rep$ in $\PP^{1,\an}_k$, the sequence $\gamma^n\cdot x $ converges to $ x_\att$ when $n\to\infty$.
\item $\mathopen{|}\tr(\gamma)| \leq 1$: then $\gamma$ admits a fixed point in $\hh_k$ and more precisely if:
\begin{itemize}
\item $\tr^2(\gamma) =4$: then $\gamma$ is not diagonalizable, and is conjugate in $\PGL(2,k)$ to $z\mapsto z+1$; thus it fixes a segment $[x,y] \in \PP^{1,\an}_k$ where $x$ (resp. $y$) is a type 2 (resp. type 1) point belonging to the $\PGL(2,k)$ orbit of the Gau\ss\ point.
\item $\tr^2(\gamma) \neq 4$: then in some at most quadratic extension $K/k$ the matrix $\gamma$ is diagonalizable; in addition it is conjugate to an element in $\PGL(2,K^\circ)$ and fixes a type $2$ point in $\hh_k$.
\end{itemize}
\end{itemize}
\end{prop}

In accordance with the terminology of group actions on trees, when $\mathopen{|}\tr(\gamma)| >1$ we say that $\gamma$ is \emph{hyperbolic}, otherwise it is called \emph{elliptic}. When required we can be more precise: if $\gamma$ is elliptic and $\gamma\neq \id$, we say that $\gamma$ is \emph{parabolic} when $\tr^2(\gamma) =4$ and \emph{strictly elliptic} otherwise.

%\medskip

One cannot say much more on the action of $\gamma$ on $\PP^1_k$ in the elliptic case. It depends heavily on the residue field $\widetilde{k}$. When the characteristic of $\widetilde{k}$ is $p>0$, then the closure of the subgroup generated by $\gamma$ is isomorphic to $\Z_p$ and $\gamma^{p^n} \to \id$ when $n\to\infty$.

%\medskip

Following standard terminology, we say that $\gamma$ has \emph{good reduction} if it fixes the Gau{\ss} point (i.e. belongs to $\PGL(2,k^\circ)$), and \emph{potential good reduction} over $K/k$ if it is conjugate in $\SL(2,K)$ to a map having good reduction. With notation as in Section~\ref{subs:hyperbolic arbitrary} this is equivalent to saying that $\gamma$ fixes a type 2 point $x\in \PP^{1,\an}_k$ such that $\sigma_{K/k}(x)$ lies in the $\PGL(2,K)$-orbit of the Gau\ss\ point.



\begin{proof}
The diagonalizability of $\gamma$ depends on the roots of its characteristic polynomial.
If $\mathrm{char}(k)\neq 2$, this can be read off the discriminant $\tr^2(\gamma)-4$.

Suppose that $\mathopen|\tr(\gamma)\mathclose| >1$. Then $\gamma$ is diagonalizable in a quadratic extension $K$ of $k$ and its eigenvalues have respective norms larger and smaller than 1. This implies that there exists a global attracting point $\PP^1(K)$ which in particular attracts all elements of $k$. Hence by completeness this fixed point belongs to $\pu(k)$. Applying the same reasoning to the inverse, we conclude that $\gamma$ is diagonalizable over $k$.


Suppose now that $\mathrm{char}(k)\neq 2$ and $\tr^2(\gamma) =4$. Then (up to sign) 1 is an eigenvalue of multiplicity 2, hence since $\gamma$ is not the identity it is conjugate in $\PGLk$ to $\lrpar{
\begin{smallmatrix} 1& 1 \\ 0& 1
\end{smallmatrix}}$ and $\gamma(z)$ is conjugate to a translation.
If $\mathrm{char}(k)=2$ and $\tr(\gamma) = 0$, then the characteristic polynomial is $X^2+1 = (X-1)^2$ and the same discussion applies.

Suppose finally that $\mathopen|\tr(\gamma)\mathclose|\leq 1$ and $\tr(\gamma) \neq 2$. Then $\gamma$ is diagonalizable over an extension $K$ of $k$ of degree at most 2, and both eigenvalues belong to $K^\circ$. If $K=k$, then the existence of the announced fixed point in $\hh_k$ is clear. Otherwise consider the geodesic joining the two fixed points in $\PP^{1,\an}_K$. The Galois group $K/k$ acts on this geodesic and permutes these two points. Thus it admits a fixed point which is of type $2$ and lies in $\hh_k$ by Lemma~\ref{lem:fixx}.
\end{proof}

Here is a noteworthy consequence of this classification.

\begin{coro}
For $\gamma\in \PGLk$, $\norm{\gamma^n} \to \infty$ as $n\cv\infty$ if and only if $\gamma$ is hyperbolic.
\end{coro}

The next result is an analogue of the Cartan KAK decomposition in the non-Archimedean setting. It will play an important role in the following.

\begin{prop}\label{prop:KAK}
Any element $\gamma\in \SL(2,k)$ can be decomposed as a product $\gamma= m\cdot a \cdot n$ with $m,n \in \SL(2,k^\circ)$ and $a = \diag (\lambda, \lambda^{-1})$ with $\lambda\in k$, $|\lambda|\ge 1$. Furthermore $\norm{\gamma} = \norm{a} = \abs{\lambda}$.
\end{prop}

\begin{proof}
Let $x_{\mathrm{g}}$ be the Gau{\ss} point. Pick an element $m\in \SL(2,k^\circ)$ such that $m^{-1} \gamma\cdot x_{\mathrm{g}}$ belongs to the segment $[x_{\mathrm{g}}, \infty]$. Likewise, choose $n\in \SL(2,k^\circ)$ such that $n^{-1} \gamma^{-1} \cdot x_{\mathrm{g}}$ belongs to the segment $[x_{\mathrm{g}}, 0]$. Then $\gamma' = m^{-1} \gamma n$ either fixes $x_{\mathrm{g}}$ or maps $x_{\mathrm{g}}$ into $(x_{\mathrm{g}}, \infty)$ and its inverse into $(x_{\mathrm{g}},0)$.

In the former case $\gamma'$ belongs to $\SL(2,k^\circ)$ hence $\gamma$ too and we can choose $a = \id$, $m = \gamma$, $n=\id$.

In the latter case, $\gamma'$ is hyperbolic with two fixed points $|c^+| >1$ and $|c^-|<1$. We claim that in this case we can conjugate it by an element in $\SL(2,k^\circ)$ so that it becomes diagonal. Indeed we first conjugate by $\lrpar{
\begin{smallmatrix}
1 & -c^- \\ 0 & 1
\end{smallmatrix}}$ (i.e. by the translation $z\mapsto z-c^-$), which belongs to $\SL(2,k^\circ)$, to send $c^-$ to 0. This maps $c^+$ to $ \tilde c^+ = c^+ - c^-$ which has the same norm. Then we use the element $\lrpar{
\begin{smallmatrix}
1 & 0\\ -1/\tilde c^+ & 1
\end{smallmatrix}}\in \SL(2,k^\circ)$ to send $ \tilde c^+$ to $\infty$, and we are done.

To prove the identity on $\norm{\gamma}$ simply observe that $\norm{\gamma} = \norm{man} \le \norm{a}$ and $ \norm{a} = \norm{m^{-1} \gamma n^{-1}} \le \norm{\gamma}$.
\end{proof}

\subsection{Norms of elements in \texorpdfstring{$\SLk$}{SLk}}

\begin{lemm}
For $\gamma\in \PGLk$, one has the identity $d_{\hh_k}(x_{\mathrm{g}}, \gamma\cdot x_{\mathrm{g}}) = \log \norm{\gamma}$.
\end{lemm}

\begin{proof}
Any element in $\PGL(2,k^\circ)$ has norm $1$ and also fixes the Gau{\ss} point so the formula is clear in this case. In the general case we use the KAK decomposition and write $\gamma = man$. Then $d_{\hh_k}(x_{\mathrm{g}}, \gamma\cdot x_{\mathrm{g}}) = d_{\hh_k}(x_{\mathrm{g}}, a \cdot x_{\mathrm{g}}) = \norm{a}=\norm{\gamma}$.
\end{proof}

A similar argument shows:

\begin{lemm}\label{lem:lipschitz}
For $\gamma\in \PGLk$ and $x, y\in \PP^1(k)$, then
\[
\norm{\gamma}^{-2} d_{\sph}(x, y)\leq d_{\sph}(\gamma x, \gamma y)\leq
\norm{\gamma}^{2} d_\mathrm{sph}(x, y)
\]
and this bound is optimal.
\end{lemm}

The following geometric consequence of Proposition~\ref{prop:KAK} will be very useful.

\begin{prop}\label{prop:B att rep}
For every $\gamma\in \PGLk$, there exist two closed balls $B_\att(\gamma)$ and $B_\rep(\gamma)$ in $\PP^{1,\an}_k$ of spherical radius $\norm{\gamma}\inv$ and such that $\gamma(\PP^{1,\an}_k\setminus B_\rep(\gamma))\subset B_\att(\gamma)$.
\end{prop}

\begin{proof}
This is a straightforward consequence of Proposition~\ref{prop:KAK}. Indeed, with notation as in Proposition~\ref{prop:KAK}, the result is obvious for $a$, with $B_\rep(a) = \overline{B}^{an}(0, \abs{\la}\inv)$ and $B_\att(a) = \overline{B}^{an}(\infty, \abs{\la}\inv) = \pp^{1,\an}_k\setminus B(0, \abs{\la})$.

In the general case, writing $\gamma =man$, it suffices to put $B_{\rep}(\gamma) = n\inv (\overline{B}^{an}(0, \abs{\lambda}\inv))$ and $B_{\att}(\gamma) = m (\overline{B}^{an}(\infty, \abs{\lambda}\inv))$.
\end{proof}

For $\gamma = \lrpar{
\begin{smallmatrix} a & b \\ c & d
\end{smallmatrix}} \in \SLk$ and $v\in \pp^1(k)$, let
\begin{equation}\label{eq:sigma1}
\sigma(\gamma,v) = \log\frac{\norm{\gamma V}}{\norm{V}} =\log \lrpar{\max (\abs{ax+by}, \abs{cx+dy}) }\,,
\end{equation}
where $V\in k^2$ is any representative of $v$, and $(x,y)\in k^2$ is a representative of $v$ with $\max (\abs{x}, \abs{y})=1$. The second equality shows that $\sigma(\gamma,{\cdot\,})$ extends continuously to $\berk$, and we have the cocycle relation
\begin{equation}\label{eq:cocycle}
\sigma(\gamma_1 \gamma_2,v) = \sigma(\gamma_1, \gamma_2 \cdot v) + \sigma(\gamma_2, v),
\end{equation}
for all $\gamma_1, \gamma_2 \in \SL(2,k)$ and for any $v\in \PP^{1,\an}_k$.


\begin{lemm}\label{lem:expand vector 1}
For any $\gamma\in \SL(2,k)$ we have that
\[
- \log \|\gamma\| \le \sigma(\gamma,v) \le \log\|\gamma\|,
\]
and furthermore if $v\notin B_\rep(\gamma)$, then
\[
\sigma(\gamma,v) = \log\lrpar{\|\gamma\| d_{\mathrm{sph}}(v,B_\rep(\gamma))}.
\]
\end{lemm}

\begin{proof}
The first assertion is obvious from~\eqref{eq:sigma1}. For the second one, observe first that $\sigma(\gamma,v) = 0$ for all $v$ when $\gamma\in \SL(2,k^\circ)$. Otherwise, writing, $ \gamma = m a n$ as in Proposition~\ref{prop:KAK} and using the cocycle relation we infer that $\sigma(\gamma, v) = \sigma(a, nv)$, where $a = \diag(\lambda, \lambda^{-1})$, and $\abs{\lambda}>1$.

From Proposition~\ref{prop:B att rep} we have that $B_{\rep}(a) = B(0, \abs{\la}^{-1})$ and $B_{\rep}(\gamma) = n^{-1}B_{\rep}(a)$. Thus $v\notin B_{\rep}(\gamma)$ if and only if $nv\notin B_{\rep}(a)$, that is $d_{\mathrm{sph}}(nv, 0) > \abs{\lambda}\inv$, and by the ultrametric property
\[
d_{\mathrm{sph}}(v, B_{\rep}(\gamma)) =d_{\mathrm{sph}}(nv, B_{\rep}(a))= d_{\mathrm{sph}}(nv, 0).
\]
Finally let $nv = [x:y]$ with $\max(\abs{x}, \abs{y}) = 1$, so that $d_{\mathrm{sph}}(nv, 0) = \abs{x}$. If $\abs{x}\leq 1$ we get $\sigma(a,nv) = \log \abs{\lambda x} $ and if $\abs{y}\leq 1$ we get $\sigma(a,nv) = \log \abs{\lambda}$. The proof is complete.
\end{proof}

\begin{lemm}\label{lem:expand vector 2}
For any \emph{hyperbolic} element $\gamma\in \SL(2,k)$ the balls $B_\att(\gamma)$ and $B_\rep(\gamma)$ are disjoint and we have
\begin{equation}\label{eq:expand}
\|\gamma\| = \frac{\mathopen|\tr(\gamma)\mathclose|}{\delta}
\end{equation}
where $\delta = d_{\sph} (B_\att(\gamma), B_\rep(\gamma)) =\inf_{x\in B_\att(\gamma)} \inf_{y\in B_\rep(\gamma)} d_{\sph}(x,y)$.
\end{lemm}

\begin{proof}
If $\gamma$ is hyperbolic the disjointness of the two balls $B_\att(\gamma)$ and $B_\rep(\gamma)$ follows from their construction and the ultrametric property implies that $\delta = d_{\sph} (x,y)$ for any pair $(x,y)\in B_\att(\gamma) \times B_\rep(\gamma)$.

We follow the reasoning of~\cite[Lem.~2.1]{kleinbif}. Conjugate $\gamma$ by some element in $\SL(2,k^\circ)$ to send the attracting fixed point to $\infty$. This does not affect neither $\norm{\gamma}$ nor $\tr(\gamma)$, and after this
conjugacy we have
\[
\gamma =
\begin{pmatrix} a & b \\ 0 & 1/a
\end{pmatrix}
\]
so that as a M\"obius transformation
$\gamma(z) = a^2z+ ab$ for some $\abs{a}>1$. The repelling fixed point is $ab/(1-a^2)$, and its distance to $\infty$ is equal to $\delta = \min \{ 1, |(1-a^2)/ab|\} = \min \{ 1, {\abs{a}}/{\abs{b}}\}$. We then have
\[
\| \gamma \| = \max\{ |a|, |b|\} = \frac{|a|}{\min \{ 1, {\abs{a}}/{\abs{b}}\}} = \frac{|\tr(\gamma)|}{\delta}~,
\]
as was to be shown.
\end{proof}

Let us point out a kind of converse to the previous lemma.

\begin{lemm}\label{lem:distinct balls}
Let $\gamma \in \PGL(2,k)$ be such that there exist two disjoint balls $B_\mathrm{a}$ and $B_\mathrm{r}$ of radius $<1$ such that $\gamma(B_\mathrm{r}^c)\subset B_\mathrm{a}$. Then $\gamma$ is hyperbolic with attracting and repelling fixed points respectively contained in $B_\mathrm{a}$ and $B_\mathrm{r}$.
\end{lemm}

\begin{proof}
Since the complement of a ball is a ball, the existence of an attracting (resp. repelling) type 1 fixed point in $B_\mathrm{a}$ (resp. $B_\mathrm{r}$) follows from~\cite[Thm.~10.69]{baker-rumely}. The result follows.
\end{proof}



\begin{lemm}\label{lem:two points}
For any pair of distinct points $z_1, z_2 \in \PP^1(k)$, there exists a constant $C = C(z_1, z_2)>0$ such that
\[
\left| \max \{ \sigma(\gamma, z_1), \sigma(\gamma, z_2)\} - \log \| \gamma \| \right|
\le C
\]
for all $\gamma\in \SL(2,k)$.
\end{lemm}
\begin{proof}
Since $ \sigma(\gamma, z) \le \log \| \gamma\|$, we only need to prove the lower bound
\[
\max \{ \sigma(\gamma, z_1), \sigma(\gamma, z_2)\} \ge \log \| \gamma \| -C.
\]
Pick $g\in \SLk$ sending $z_1$ to $0$ and $z_2$ to $\infty$. We have
\begin{align*}
\max \{ \sigma(\gamma, z_1), \sigma(\gamma, z_2)\} &=
\max \left\{ \sigma(\gamma g^{-1}, 0) + \sigma (g,z_1), \sigma(\gamma g^{-1}, \infty) + \sigma (g,z_2)\right\}\\
&\ge \log \| \gamma g^{-1}\| -\log \| g \| \ge \log \| \gamma\| - 2\log \| g \|~.
\end{align*}
This concludes the proof.
\end{proof}



\subsection{Elementary and non-elementary subgroups} \label{subs:elementary}

A subgroup $\Gamma\leq \SLk$ (resp. $\Gamma \leq \PGLk$) is said \emph{reducible} if its action on $\PP^1(k)$ fixes a point, and \emph{irreducible} otherwise. It is \emph{strongly irreducible} if it does not admit a finite orbit in $\PP^1(k)$.

We say that $\Gamma$ has \emph{good reduction} if it takes values in $\SL(2,k^\circ)$, or equivalently, fixes the Gau\ss~point. It has \emph{potential good reduction} if there exists a finite field extension $K/k$ such that $\Gamma$ is conjugate in $\SL(2,K)$ to a subgroup of $\SL(2,K^\circ)$. Finally, $\Gamma$ is \emph{proximal } if it contains at least one hyperbolic element.

\begin{prop}\label{prop:fixed point}
A finitely generated subgroup $\Gamma$ of $\PGLk$ is either proximal or has potential good reduction. If $\Gamma$ has potential good reduction and in addition $k$ is discretely valued, then $\Gamma$ is conjugate to a subgroup of $\SL(2,K^\circ)$ in some quadratic extension $K$ of $k$.
\end{prop}

Observe that the group of translations is not proximal but has not potential good reduction when the norm on $k$ is non-trivial, so that the assumption that $\Gamma$ is finitely generated is necessary in the previous statement.

%\smallskip

This proposition is essentially a formulation in our language of the well-known fact that a group acting on a tree with only elliptic elements has a global fixed point. We sketch the proof for convenience.

The key is the following lemma (see~\cite[Lem.~10.4]{kapovich} or~\cite[Lem.~40]{otal}).

\begin{lemm}\label{lem:semi group}
Any finitely generated semi-group of $\SL(2,k)$ which does not contain any hyperbolic element fixes a type 2 point lying in $\hh_k$.
\end{lemm}

\begin{proof}
Let $S$ be a finitely generated semi-group which does not contain any hyperbolic element. We shall prove by induction on the number of generators the existence of a type 2 point in $\hh_k$ fixed by $S$.

When $S$ is generated by a single element this is a direct consequence of Proposition~\ref{prop:claselt}. If $S$ is generated by two elements $g$ and $h$, pick $x,y\in \hh_k$ two type 2 points fixed by $g$ and $h$ respectively. Let $x'$ be the unique point satisfying $[x,g(y)] \cap [x,y] = [x,x']$. By Proposition~\ref{prop:projection} this is a type 2 point. Similarly define $y'$ to be the unique type 2 point satisfying $[y,h(x))] \cap [y,x] = [y,y']$. If the segment $[x',y']$ is degenerate, the segment $[y',h(y')]$ is a fundamental domain for the action of $gh$ which is therefore hyperbolic. Otherwise $gh$ fixes pointwise $[x',y']$. This proves the result when $S$ is generated by two elements.

Now suppose $S$ is generated by $g_1, \ldots, g_l$ with $l\ge 3$, and that the result is known for semi-groups generated by $l-1$ elements. For $i =1,2,3$, let $S_i$ be generated by $\set{ g_1, \ldots, g_l}\setminus \set{g_i}$. By the induction hypothesis $S_i$ admits a type 2 fixed point $x_i \in \hh_k$. Then the projection of $x_3$ on $[x_1, x_2]$ is a type 2 point in $\hh_k$ fixed by $S$ and we are done.
\end{proof}



\begin{proof}[Proof of Proposition~\ref{prop:fixed point}]
It follows from Lemma~\ref{lem:semi group} that if $\Gamma$ is not proximal then it fixes a type 2 point $x_\star$ in $\hh_k$. Using the notation of Section~\ref{subs:hyperbolic arbitrary} this means that $\sigma_{\bar{k}^a/k}(x_\star)$ lies in the $\PGL(2,\bar{k}^a)$-orbit of the Gau\ss{} point. Since algebraic points over $k$ are dense in $\bar{k}^a$, the ball corresponding to $\sigma_{\bar{k}^a/k}(x_\star)$ contains a point of $k^a$ so we get that $\sigma_{\bar{k}^a/k}(x_\star)$ lies in fact in the $\PGL(2,k^a)$-orbit of the Gau\ss{} point. In other words, we can find a finite field extension $K/k$ and conjugate $\Gamma$ by a matrix in $\PGL(2,K)$ so that it fixes the Gau\ss{} point.

Assume now that $k$ is discretely valued so that $\hh_k$ is a simplicial tree. The point $x_\star$ is either in the $\PGL(2,k)$-orbit of the Gau\ss{} point or it belongs to a unique segment $[x_0,x_1]$ of $\hh_k$ whose extremities lie in the $\PGL(2,k)$-orbit of the Gau\ss{} point. Any element fixing $x_\star$ either fixes pointwise $[x_0,x_1]$ or acts upon it as an involution switching the two extremities. It follows that the middle point of $[x_0,x_1]$ is fixed by $\Gamma$. We conclude using Proposition~\ref{prop:quadratic}.
\end{proof}


There is a simple classification of subgroups that are not strongly irreducible, analogous to the Archimedean case.

\begin{prop}\label{prop:strongly irreducible}
Let $\Gamma\leq \PGLk$ be a finitely generated subgroup that is not strongly irreducible. Then one of the following situations occurs:
\begin{enumerate}
\item $\Gamma$ has potential good reduction;
\item $\Gamma$ is conjugate to a subgroup of the affine group $\{z\mapsto a z + b, \, a \in k^\times, \, b\in k\}$;
\item $\Gamma$ is conjugate to a subgroup of $\{z\mapsto \lambda z^{\pm 1},\, \lambda \in k\}$.
\end{enumerate}
\end{prop}

\begin{proof}
By assumption there exists a finite $\Gamma$-orbit $x_1, \ldots, x_n$ on $\pp^1(k)$. If $n=1$ then $\Gamma$ is conjugate to a subgroup of the affine group. If $n=2$, we may assume that $x_1 =0$ and $x_2 =\infty$ and it follows that any element $\Gamma$ is conjugate to $\lambda /z$ or $\lambda z$ for some $\lambda\in k^*$.

Assume now that $\Gamma$ leaves invariant a set of $n\ge3$ distinct points $E= \{ x_1, \ldots, x_n \}$ in $\pp^1(k)$. The first observation is that for every $\gamma\in \Gamma$, some iterate $\gamma^m$ fixes $E$ pointwise, therefore $\gamma^m = \id$. All elements of $\Gamma$ are thus elliptic and the previous proposition shows that $\Gamma$ has potential good reduction.
\end{proof}

\begin{rema}
If $\mathrm{char}(k) = 0$, then by the Selberg lemma (see e.g.~\cite{alperin}) the existence of a finite orbit of cardinality $n\geq 3$ implies that $\Gamma$ is finite.
\end{rema}



\begin{prop}\label{prop:non elementary}
Let $\Gamma\leq \PGLk$ be a finitely generated subgroup. If $\Gamma$ is proximal and strongly irreducible then it contains two hyperbolic elements with disjoint sets of fixed points.
\end{prop}

\begin{proof}
Since $\Gamma$ is proximal, it contains a hyperbolic element $g$. Denote by $x_{\att/\rep}$ its fixed points. We claim that there exists an element $h \in \Gamma$ such that $h(\{ x_\att, x_\rep\}) \cap\{ x_\att, x_\rep\}= \emptyset$. Indeed since $\Gamma$ is strongly irreducible, $\{ x_\att, x_\rep\}$ is not a $\Gamma$-orbit so that there exists $h\in \Gamma$ satisfying $h(x_\att) \notin \{ x_\att, x_\rep\}$. There are 3 possibilities:
\begin{itemize}
\item $h(x_\rep)\notin \{ x_\att, x_\rep\}$;
\item $h(x_\rep) = x_\att$: then either $hgh$ or $h^2$ sends $\{ x_\att, x_\rep\}$ to a disjoint pair;
\item $h(x_\rep) = x_\rep$: then there exists $j\in \Gamma$ such that $j(x_\rep)\notin \{ x_\att, x_\rep\}$ and $hg^nj\inv$ is convenient for large $n$ (use Proposition~\ref{prop:B att rep}).
\end{itemize}
In any case there exists $\tilde{h}$ such that $\tilde{h}(\{ x_\att, x_\rep\})\cap \{ x_\att, x_\rep\}=\emptyset$, thus $\tilde{h}^{- 1} \gamma \tilde{h}$ is a hyperbolic element whose fixed points are disjoint from $\{ x_\att, x_\rep\}$.
\end{proof}

Proposition~\ref{prop:non elementary} motivates the following definition.

\begin{defi}
A finitely generated subgroup $\Gamma$ of $\PGL(2,k)$ is \emph{non-elemen\-tary} if it is proximal and strongly irreducible.

A finitely generated subgroup $\Gamma$ of $\SL(2,k)$ is non-elementary if its image in $\PGL(2,k)$ is non-elementary.
\end{defi}

Propositions~\ref{prop:fixed point} and~\ref{prop:strongly irreducible} imply the following characterization of non-elementary subgroups. The details are left to the reader.

\begin{prop}
Let $\Gamma\leq \PGLk$ be a finitely generated subgroup. The following assertions are equivalent:
\begin{enumerate}
\item $\Gamma$ is non-elementary;
\item $\Gamma$ does not admit a finite orbit on $\PP^{1, \mathrm{an}}_k$;
\item for every $z\in \PP^{1, \mathrm{an}}_k$, $\# \Gamma \cdot z\geq 3$.
\end{enumerate}
\end{prop}

Let us note for further reference the following variation on Proposition~\ref{prop:non elementary}.

\begin{lemm}\label{lem:proximal semi group}
Let $\Gamma$ be a non-elementary finitely generated subgroup of\linebreak $\PGLk$. Then for every set $S$ of generators of $\Gamma$, the semi-group generated by $S$ contains two hyperbolic elements with distinct attracting fixed points.
\end{lemm}

\begin{proof}
First note that since $\Gamma$ is finitely generated, there is a finite subset $S'\subset S$ such that $\langle S'\rangle$ contains a finite set of generators of $\Gamma$, hence $\langle S'\rangle = \Gamma$, so, replacing $S$ by $S'$ we may assume that $S$ is finite. Denote by $G_0$ the semi group generated by $S$. Since $\langle S\rangle = \Gamma$ the elements of $S$ do not admit a common fixed point, hence by Lemma~\ref{lem:semi group} there exists a hyperbolic element $g\in G_0$. Thus, letting $\rho = \norm{g}>1$ we infer that for $n\geq 0$, $g^n$ maps $B(x_\rep(g), \rho^n)^c$ into $B(x_\att(g), \rho^n)$ (all the balls here are in $\PP^{1, an}_k$). Since $S$ has no fixed point there exists $h\in S$ such that $h(x_\att(g))\neq x_\att(g)$. Then for every $n\geq 1$, $hg^n$ belongs to $G_0$ and maps $B(x_\rep(g), \rho^n)^c$ into $B(h(x_\att(g)), C\rho^n)$ for some $C = C(h)$. If $h(x_\att(g))\neq x_\rep(g)$, from Lemma~\ref{lem:distinct balls} we infer that $hg^n$ is hyperbolic for sufficiently large $n$, and its attracting fixed point is close to $h(x_\att(g))$, hence distinct from $x_\att(g)$.

If $h(x_\att(g))= x_\rep(g)$, then $h\inv(x_\rep(g)) =x_\att(g)$, and we consider $hg^n h$ instead of $hg^n$. Indeed for some $C$ we have that
\[
B(x_\att(g), C\inv\rho^n)^c\overset{h}\longrightarrow B(x_\rep(g), \rho^n)^c\overset{g^n}\longrightarrow B(x_\att(g), \rho^n)
\overset{h}\longrightarrow B(x_\rep(g), C \rho^n),
\]
so again we see that $hg^nh$ is hyperbolic for large $n$, and its attracting fixed point is distinct from $x_\att(g)$.
\end{proof}

\subsection{The limit set}

\begin{theo}\label{thm:struct limit}
Let $\Gamma$ be a finitely generated and non-elementary subgroup of $\PGLk$. Then the following sets coincide:
\begin{itemize}
\item the closure in $\PP^{1,\an}_k$ of the set of fixed points of all hyperbolic elements of $\Gamma$;
\item the smallest non-empty $\Gamma$-invariant closed subset of $\PP^{1,\an}_k$;
\item for any given $x\in \PP^{1,\an}_k$, the set of points $y$ such that there exists a sequence $(g_n)\in \Gamma^\nn$ such that $\norm{g_n}\to\infty$ and $g_n \cdot x \to y$.
\end{itemize}
This set is compact, metrizable, and included in $\conv(\PP^1(k)) =\hh_k \cup \PP^1(k)$. It is by definition the limit set $\Lm(\Gamma)$ of $\Gamma$.
\end{theo}

\begin{proof}
Denote by $\Lambda_0$ the set of fixed points of all hyperbolic elements, and by $\Lambda$ the smallest closed $\Gamma$-invariant subset of $\PP^{1,\an}_k$. If $x$ is fixed by some hyperbolic element $g$, then for every $h\in \Gamma$, $h(x)$ is fixed by $hgh\inv$. We infer that $\overline{\Lambda}_0$ is a closed $\Gamma$-invariant set thus $\Lambda\subset \overline{\Lambda}_0$.

Conversely, pick any $x\in \Lambda$. By Proposition~\ref{prop:non elementary} there exists a hyperbolic element $g\in \Gamma$ whose fixed point set $\{ x_\att, x_\rep\}$ is disjoint from $x$. Since $g^{\pm n}(x) \to x_{\att/\rep} $ it follows that $\{x_\att, x_\rep\} \subset \Lambda$ so $\Lambda$ admits at least three points. Therefore, for an arbitrary hyperbolic element $g'\in\Gamma$ with fixed point set $\{ x'_\att, x'_\rep\}$, there exists $y\in \Lambda\setminus \{ x'_\att, x'_\rep\}$. Then $(g')^{\pm n}(y) \to x'_{\att/\rep}$ as $n\to\infty$, from which we infer that $\Lambda_0 \subset \Lambda$. We conclude that $\Lambda = \overline{\Lambda}_0 $.

%\medskip

Fix now any point $x\in \PP^{1,\an}_k$ and denote by $\Lambda_1$ the set of all $y$ for which there exists a sequence $(g_n)$ with $\norm{g_n}\to\infty$ and $g_n\cdot x \to y$. Observe that $\Lambda_1$ is $\Gamma$-invariant and non-empty since $\Gamma$ contains a hyperbolic element. We claim that it is also closed. Indeed by a theorem of Poineau~\cite[Thm.~5.3]{poineau}, for any $y'$ in the closure of $\Lambda_1$ there exists a sequence $y_n \in \Lambda_1$ such that $y_n \to y'$. For each $n$, pick a sequence with $\norm{g_{m,n}}
\ge m+n$ such that $g_{m,n}\cdot x \to y_n$. The set $\{ g_{m,n}\cdot x\}$ contains $\{y_n\}$ in its closure hence $y'$ too. Again by Poineau's theorem, there exists a subsequence $g_{m_j,n_j}\cdot x \to y$. This shows that $\Lambda_1\supset\Lambda$.

%\smallskip

Now suppose $\norm{g_n}\to\infty$ and $g_n\cdot x \to y\in\Lambda_1$. We want to show that $y$ belongs to $\Lambda$. If $x$ belongs to $\Lambda$ then the closure of $\Gamma\cdot x$ is contained in $\Lambda$ and the result follows. So suppose that $x$ does not belong to $\Lambda$.

Recall from Proposition~\ref{prop:B att rep} that we can associate to every $g\in \PGLk$ two closed balls $B_{\att}(g)$ and $B_{\rep}(g)$ in $\PP^{1,\an}_k$ such that $g (\berk\setminus B_{\rep}(g)) \subset B_{\att}(g)$. We claim that for large enough $n$, $B_{\att}(g_n)$ intersects $\Lambda$. Indeed pick any 2 distinct points in $\Lambda$. Then since $\norm{g_n}\to \infty$, for large $n$ one of these points does not belong to $B_{\rep}(g_n)$, hence its image under $g_n$ belongs to $B_{\att}(g_n)$, and also to $\Lambda$ by invariance, so we get that $B_{\att}(g_n)\cap \Lambda\neq \emptyset$. Similarly, $B_{\rep}(g_n)\cap \Lambda\neq \emptyset$.

In particular we see that for large $n$, $x\notin B_{\rep}(g_n)$. Indeed otherwise since the diameter of $B_{\rep}(g_n)$ tends to zero we would infer that $x$ belongs to $\Lambda$, which is not the case.

Thus we conclude that $g_n\cdot x\in B_{\att}(g_n)$ for large $n$, so every neighborhood of $y$ intersects $\Lambda$ and it follows that $\Lambda_1\subset\Lambda$.

%\medskip

The limit set is a closed subset of $\PP^{1,\an}_k$ which is compact, hence it is also compact. Since $\Gamma$ is countable hence $\Lambda_0$ is countable too. It follows that $\Lambda$ is included in the closure of the convex hull of a countable set. Such a set is always metrizable (see e.g. the proof of~\cite[Lem.~7.15]{valtree}, or~\cite[Lem.~5.7]{baker-rumely}). Finally $\Lambda_0$ is a subset of $\PP^1(k)$, hence $\Lm(\Gamma)$ is included in its closure which is contained in $\hh_k \cup \PP^1(k)$.
\end{proof}


\section{Random products of matrices in \texorpdfstring{$\SLk$}{SLk}}\label{sec:furstenberg}

In this section we work on an arbitrary complete non-trivially valued field $(k, \abs{{\,\cdot\,}})$ (shortly to be assumed non-Archimedean). We keep the notation of the previous section. We consider a measure $\mu$ with at most countable support in $\SLk$, and make the following assumptions:
\begin{enumerate}\renewcommand{\theenumi}{B\arabic{enumi}}
\item \label{B1} $\Gamma = \langle \supp(\mu) \rangle$ is non-elementary.
\item \label{B2} $\mu$ has finite first moment $ \int \log \norm{\gamma} \dmu(\gamma)<\infty$.
\end{enumerate}

The measure $\mu$ acts by convolution on the set of probability measures on $\berk$ by $\nu \mapsto \mu \ast \nu$. The measures invariant under this action are called \emph{stationary}. We use the probabilistic notation $(\om, \prob) = (\SLk^{\nn^*}, \mu^{\nn^*})$, and for $\omega = (\gamma_n)_{n\geq 1}\in \om$ we respectively let $r_n(\omega) = \gamma_1\cdots \gamma_n$ and $\ell_n(\omega) = \gamma_n\cdots \gamma_1$ be the $n^{\rm th}$ step of the right and left random walk on $\Gamma$ with transition probabilities given by $\mu$. We denote by $\mu^n$ the $n^{\rm th}$ convolution power of $\mu$, that is the image of $\mu^ {\otimes n}$ under the map $(g_1, \ldots,g_n)\mapsto g_1\cdots g_n$. It is also the distribution of both $r_n(\omega)$ and $\ell_n(\omega)$.

The \emph{Lyapunov exponent} of $\mu$ is defined to be the following non-negative real number:
\begin{equation}\label{def:lyapunov exponent}
\chi (\mu) : = \lim _{n\rightarrow \infty} \frac{1}{n}\int
\log \norm{\gamma} \dmu^{n}(\gamma) = \lim _{n\rightarrow \infty} \frac{1}{n}\int \log \norm{ \gamma_1\cdots
\gamma_n} \dmu(\gamma_1)\cdots \dmu(\gamma_n)~.
\end{equation}
It follows from Kingman's sub-additive ergodic theorem that the limit of $ \frac1n \log \norm {\ell_n(\omega)}$ as $n\to\infty$ exists for $\prob$-a.e. $\omega$ and equals $\chi(\mu)$.

%\medskip

The main result of this section is the following theorem. Recall the notation $\sigma(\gamma, v) = \log\norm{\gamma V}$, where $V = (V_1, V_2)$ is a lift of norm 1 of $v$ and $\norm{V} = \max(\abs{V_1}, \abs{V_2})$.

\begin{theo}\label{thm:positive}
Let $\mu$ be a probability measure with countable support in $\SLk$, satisfying~\eqref{B1}. Then there exists a unique stationary probability measure $\nu$ on $\berk$, that is stationary under the action of $\mu$. This measure has no atoms, it is supported on the limit set of $\Gamma$, and when $k$ is non-Archimedean it gives full mass to $\PP^1(k)$ so that $\nu (\hh_k)=0$.

If furthermore~\eqref{B2} holds, then the Lyapunov exponent $\chi(\mu)$ of the associated random product of matrices is positive and satisfies the following formula
\begin{align}\label{eq:furstenberg}
\chi(\mu) & = \int \sigma(\gamma, v) \, \dmu(\gamma)\, \dnu(v).
\end{align}

Finally for $\prob$-a.e. $\omega$ and for $\nu$-a.e. $v$, we have:
\begin{align}
\chi(\mu) = \lim_{n\to\infty} \frac1n \log \sigma(\ell_n(\omega), v) \label{eq301}
.
\end{align}
\end{theo}

\begin{rema}
If $\Gamma$ is generated by $\supp(\mu)$ as a semi-group then $\supp(\nu) = \mathrm{Lim}(\Gamma)$. Indeed $\supp(\nu)$ is contained in $\mathrm{Lim}(\Gamma)$, closed, and $\supp(\mu)$-invariant, hence $\Gamma$-invariant. In the general case, however, the inclusion $\supp(\nu)\subset \mathrm{Lim}(\Gamma)$ can be strict.
\end{rema}

There are many statements of this kind in the literature, and when $k$ is archimedean this statement is precisely Furstenberg's theorem on random products of matrices~\cite{furstenberg} (see~\cite[Chap.~II]{bougerol-lacroix} for a simple exposition). In the non-Archimedean setting, it was observed by several authors (see in particular~\cite{guivarch}) that Furstenberg's theory can be adapted without much harm to local fields. The novelty here is that $k$ is arbitrary, in particular may not be locally compact. This leads us to resort to Berkovich theory, and also prevents us from taking cluster limits of sequences of elements in $\SLk$, which is commonplace in the classical presentation of the topic.

On the other hand it was recently proved by Maher and Tiozzo~\cite{maher-tiozzo} that non-elementary random walks on non-necessarily proper Gromov hyperbolic spaces have positive drift, from which the first conclusion of the theorem follows.

Nevertheless we provide a complete proof of the theorem for at least two reasons: first our algebraic setting allows us to provide a relatively short and self-contained proof, and also the conclusion on the stationary measure does not straightforwardly follow from~\cite{maher-tiozzo} since Maher and Tiozzo work in a horofunction compactification that is not directly related to the Berkovich projective line.

%\medskip

As in the classical case, the uniqueness of the stationary measure and the positivity of the Lyapunov exponent follow from a contraction statement, which asserts that if $\nu$ is any stationary measure, then for $\prob$-a.e. $\omega$, $\norm{r_n(\omega)}\to \infty$ and $(r_n(\omega))_*\nu$ converges to a Dirac mass at a point ${e(\omega)}$ which does not depend on $\nu$ (see below Lemma~\ref{lem:contraction}). To prove this result, we adapt the arguments of Guivarc'h and Raugi~\cite{guivarch-raugi}.

In the remaining of this section we assume that the norm on $k$ is \emph{non-Archimedean}.

\subsection{Uniqueness of the stationary measure} \label{subs:abstract2}

As a first step towards Theorem~\ref{thm:positive} in this section we prove the following result.

\begin{theo}\label{thm:partial unique}
Let $\mu$ be a probability measure with countable support in $\SLk$, satisfying~\eqref{B1}, and suppose $\nu$ is a $\mu$-stationary measure having no atoms and such that $\nu(\PP^1(k)) = 1$.
\begin{enumerate}
\item There exists a measurable map $e: \Omega \to \PP^1(k)$ such for a.e. $\omega$, $r_n(\omega)_* \nu \to \delta_{e(\omega)}$. Moreover for a.e. $\omega$ we have $\|r_n(\omega)\| \to \infty$, and $d_{\mathrm{sph}}(e(\omega), B_\att(r_n(\omega))) \le 2 \|r_n(\omega)\|^{-1}$.
\item The identity $\nu = \int \delta_{e(\omega)} \, \dprob(\omega)$ holds, and $\nu$ is the unique $\mu$-stationary measure which has no atoms and gives full mass to $\PP^1(k)$.
\end{enumerate}
\end{theo}

We start with the following classical lemma.

\begin{lemm}\label{lem:guivarch raugi}
Under the assumptions of Theorem~\ref{thm:partial unique}, for $\prob$-a.e. $\omega$ there exists a probability measure $\nu_\omega$ such that for every $\gamma$ belonging to the semi-group $G$ generated by $\supp(\mu)$, the sequence $r_n(\omega)_*\gamma_*\nu$ converges weakly to $\nu_\omega$. This measure a.s. gives full mass on $\PP^1(k)$ and $\nu = \int \nu_\omega \; \dprob(\omega)$.
\end{lemm}

\begin{proof}
The support of $\nu$ is a compact metrizable space by Lemma~\ref{lem:support}, so that we may apply~\cite[Lem.~2.1, p.~19]{bougerol-lacroix}. We obtain for $\prob$-a.e. $\omega$ the existence of a probability measure $\nu_\omega$ such that for $\mu^\infty$-a.e. $\gamma$, $r_n(\omega)_*\gamma_*\nu$ converges weakly to $\nu_\omega$, where $\mu^\infty = \sum_{n = 0}^\infty 2^{-n-1}\mu^n$. Since $\Gamma$ is countable the measure $\mu$ is purely atomic and its support is precisely the semi-group generated by $\supp(\mu)$. Thus for every $\gamma$ in this semi-group, we obtain $r_n(\omega)_*\gamma_*\nu \cvf \nu_\omega$.

The stationarity property implies that for every $n$, $\nu = \int r_n(\omega)_* \nu\, \dprob(\omega)$, and the dominated convergence theorem implies $\nu = \int \nu_\omega \; \dprob(\omega)$. In particular $\nu_\omega(\PP^1(k)) =1$ almost surely.
\end{proof}

Next we prove the divergence of the norms of generic random products.

\begin{lemm}\label{lem:contraction}
For $\prob$-a.e. $\omega$ we have that $\| r_n(\omega)\| \cv\infty$.
\end{lemm}

\begin{proof}
Let $\omega \in G^\nn$ be a sequence satisfying the conclusion of Lemma~\ref{lem:guivarch raugi}, and let us show that $\| r_n(\omega)\| \cv\infty$. We proceed by contradiction, so assume there exists a subsequence $(n_j)$ such that $\|r_{n_j}(\omega))\|\leq C$ for some $C\geq 1$. Since $\Gamma$ is non elementary, by Lemma~\ref{lem:proximal semi group} the semi-group generated by $\supp(\mu)$ contains two hyperbolic elements $\gamma_1$, $\gamma_2$ with distinct attracting fixed points, that we fix from now on. Denoting by $x_{\att}(\gamma_i)$ the respective attractive fixed points, we fix $r$ small enough so that $d(x_\att(\gamma_1),x_\att(\gamma_2)) > 2r (C^2+1)$.

The measure $\nu_\omega$ charges $\PP^1(k)$, so one can find a point $z\in \PP^1(k)$ such that $m := \nu_\omega(B(z,r))>0$. Since the measure $\nu$ has no atoms, for $i=1, 2$ $ \nu(B(x_\rep(\gamma_i),\rho))$ tends to $0$ when $\rho\to0$, so we can fix $N$ (large) such that
\[
((\gamma_1^N)_*\nu) (B(x_\att(\gamma_1),r)) \geq 1- \frac{m}{4} \quad\text{ and}\quad ((\gamma_2^{N})_*\nu)(B(x_\att(\gamma_2),r)) \geq 1- \frac{m}{4}.
\]

Since on the other hand $B(z,r)$ is open, from the choice of $\omega$, there exists $j$ such that
\[
\lrpar{ r_{n_j}(\omega)_* (\gamma_1^N)_*\nu} (B(z,r)) \geq \frac{m}{2}
\quad\text{ and}\quad \lrpar{ r_{n_j}(\omega)_*(\gamma_2^{N})_*} (B(z,r)) \geq \frac{m}{2}.
\]

Applying Lemma~\ref{lem:lipschitz} it follows that for $i=1,2$,
\[
(\gamma_i^N)_*\nu(B(r_{n_j}(\omega)^{-1}z, rC^2))\geq m/2,
\]
hence
\[
B(r_{n_j}(\omega)^{-1}z, rC^2)\cap B(x_\att(\gamma_i),r) \neq \emptyset.
\]
which implies that $d(x_\att(\gamma_1), x_\att(\gamma_2)) \le 2rC^2 + 2r$, a contradiction.
\end{proof}

The proof of Theorem~\ref{thm:partial unique} will be complete if we prove that

\begin{lemm}\label{lem:Dirac}
For $\prob$-a.e. $\omega$, the measure $\nu_\omega$ is a Dirac mass at a point $e(\omega)\in \PP^1(k)$ which does not depend on $\nu$, and satisfies $d_{\sph}(e(\omega), B_\att(r_n(\omega))) \le 2\|r_n(\omega)\|^{-1}$.
\end{lemm}

The proof of this lemma relies on the following elementary fact which asserts that the measure of small balls is uniformly small.

\begin{lemm}\label{lem:uniform}
Let $\nu$ be an atomless Borel probability measure on a complete metric space $(X,d)$. There exists a function $\eta : \R_+ \to \R_+$ such that $\eta(r)\to 0$ as $r\to 0$ such that for every $x\in X$ we have that $\nu(B(x,r))\leq \eta(r)$.
\end{lemm}

\goodbreak
\begin{proof}
Assume by way of contradiction that there exists $\eta>0$, a sequence $(x_n)\in X^\nn$ and a sequence of radii $r_n\cv 0$ such that $\nu(B(x_n, r_n))\geq \eta$. Extracting a subsequence we may assume that $\sum r_n$ converges. We will show that $(x_n)$ admits a Cauchy subsequence $(x_{n_j})$, thus converging to some $x$. Since for any $r>0$, we have that $B(x_{n_j}, r_{n_j}) \subset B(x, r)$ for large $j$, it follows that $\nu (B(x,r))\geq \eta$ for every $r$, contradicting the assumption on $\nu$.

To show that $(x_n)$ admits a Cauchy subsequence, we define the set
\[
A = \set{n\in \nn,\ \forall p>n, \ B(x_p, r_p) \cap B(x_n, r_n) = \emptyset}.
\]
Write $A$ as an increasing sequence $A = \set{a_1, a_2, \ldots}$. By construction, for every $k>l$, $B(x_{a_k}, r_{a_k}) \cap B(x_{a_l}, r_{a_l}) = \emptyset$. Since each of these balls has mass at least $ \eta$, there are at most $1/\eta$ of them, hence $A$ is finite.

Now if $n_1> \max A$, by assumption there exists $n_2>n_1$ such that $B(x_{n_2}, r_{n_2})\cap B(x_{n_1}, r_{n_1})\neq \emptyset$. Repeating this process we construct a subsequence $(n_j)_{j\geq 1}$. Since the series $\sum r_n$ converges, the sequence $(x_{n_j})$ is Cauchy, and the lemma follows.
\end{proof}

\begin{proof}[Proof of Lemma~\ref{lem:Dirac}]
Recall the definition of $B_\att(\gamma)$ and $B_\rep(\gamma)$ from Proposition~\ref{prop:B att rep}: these are two balls in $\PP^{1,\an}_k$ of spherical diameter $\|\gamma\|^{-1}$ such that $\gamma (\PP^{1,\an}_k \setminus B_\rep(\gamma)) \subset B_\att(\gamma)$. Observe that they necessarily intersect $\PP^1(k)$.

From Lemmas~\ref{lem:guivarch raugi} and~\ref{lem:contraction}, for almost every $\omega$ we have that $(r_n(\omega))_*\nu\cvf\nu_\omega$ and $\| r_n(\omega)\| \to \infty$. With $\eta$ as in Lemma~\ref{lem:uniform}, we have that $\nu (B_{\rep}(r_n(\omega)))\leq \eta (\| r_n(\omega)\|^{-1})$, hence
\begin{equation}\label{eq:lw meas}
\nu (B_{\att}(r_n(\omega)))\geq 1- \eta (\| r_n(\omega)\|^{-1}) \ge \frac34
\end{equation}
for $n$ large enough. For each $n$ choose any $x_n \in B_{\att}(r_n(\omega)) \cap \PP^1(k)$. Then~\eqref{eq:lw meas} implies $d_{\sph}(x_n,x_m) \le 2\|r_n(\omega)\|^{-1}$ for all $m\ge n$ hence $(x_n)$ forms a Cauchy sequence. Observe that the limit of this sequence belongs to $\PP^1(k)$, is at distance at most $2\|r_n(\omega)\|^{-1}$ from $B_\att(r_n(\omega))$, and does depend neither on the choice of the sequence $(x_n)$, nor on the stationary measure. We may thus denote it by $e(\omega)$.

Finally, for each $r>0$, we have that $B_{\att}(r_n(\omega)) \subset B(x_\infty,r)$ for large $n$ since $\|r_n(\omega)\|^{-1} \to0$. We conclude that $\nu_\omega(B(e(\omega),r)) \ge \liminf_n \nu_\omega(B_{\att}(r_n(\omega))) =1$, so that $\nu_\omega = \delta_{e(\omega)}$.
\end{proof}

\subsection{Proof of Theorem~\ref{thm:positive}}

Recall that $\mu$ is a measure with countable support on $\SL(2,k)$ satisfying condition~\eqref{B1}. We first show the existence and uniqueness of a $\mu$-stationary measure, and then prove that this measure has no atoms and gives full mass on $\PP^1(k)$. Then assuming~\eqref{B2} we establish the analog of Furstenberg's formula that expresses the Lyapunov exponent $\chi(\mu)$ in terms of the stationary measure. The positivity of $\chi(\mu)$ and the identities~\eqref{eq:furstenberg} and~\eqref{eq301} follow from this formula.
%\medskip

Recall first from Theorem~\ref{thm:struct limit} that the limit set of $\Gamma$ is a compact metrizable space. For any $x\in \Lm(\Gamma)$, one can thus extract a converging subsequence from $\frac1{n} \sum_{i=0}^{n-1}\mu^ i \ast \delta_{x}$, and the limit measure $\nu$ is $\mu$-stationary. The uniqueness of the stationary measure then follows from Theorem~\ref{thm:partial unique} together with the next lemma whose proof will be given at the end of this section.

\begin{lemm}\label{lem:Hk}
Let $\mu$ be a probability measure with countable support in $\SLk$, satisfying~\eqref{B1}. If $\nu$ is any $\mu$-stationary probability measure on $\berk$, then $\nu$ has no atoms and gives full mass to $ \PP^1(k)$.
\end{lemm}

By Lemma~\ref{lem:guivarch raugi} the support of $\nu$ is contained in the limit set. It remains to prove~\eqref{eq:furstenberg},~\eqref{eq301} and the positivity of the Lyapunov exponent. From now on we assume that the moment condition~\eqref{B2} holds, and proceed in several steps.

\begin{step}[Step 1]
We first claim that for $(P\times \nu)$ a.e. $(\omega,v)$, one has $\sigma(\ell_n(\omega), v)\cv\infty$.

To see this we introduce the reversed random walk on $\Gamma$, that is the random walk associated to $\check \mu$, the image of $\mu$ under the involution $\gamma\mapsto \gamma^{-1}$. It satisfies the assumptions~\eqref{B1}--\eqref{B2} so from what precedes we know that it admits a unique stationary measure $\check\nu$ on $\berk$ having no atoms and putting full mass on $\PP^1(k)$. Define an involution $G^\nn\to G^\nn$ by $\omega = (\gamma_n)_{n\geq 1 }\mapsto \check \omega = (\gamma_n^{-1})_{n\geq 1}$, so that $\ell_n(\omega) = r_n(\check \omega)^{-1}$, and $B_{\rep}(\ell_n(\omega)) = B_{\att}(r_n(\check \omega))$.

By Theorem~\ref{thm:partial unique}, we have $\| \ell_n(\omega)\| \to \infty$ for a.e. $\omega$, and there exists a measurable map $\check{e}: \Omega \to \PP^1(k)$ such that $\check \nu = \int \delta_{\check e(\omega)} \dd\check\mu(\omega)$, and $d_{\sph}(\check e(\omega), B_\att(\ell_n(\omega))) \le 2\| \ell_n(\omega)\|^{-1}$.

If $v\in \pp^1(k)$ is different from $\check e(\omega)$, then for $n$ large enough, $d_{\mathrm{sph}}(v, B_{\rep}(\ell_n(\omega)))\geq \delta$, where $\delta = d_{\mathrm{sph}}(v, \check e(\omega))/2$, hence from Lemma~\ref{lem:expand vector 1} we get that $\sigma(\ell_n(\omega), v)\geq \log (\norm {\ell_n(\omega)}\delta)\to \infty$. Since $\nu$ gives no mass to points, this property holds for $\nu$-a.e. $v$. Thus by Fubini's theorem we conclude that $\sigma(\ell_n(\omega), v)\cv\infty$ holds $(P\times \nu)$ a.s., as claimed.

Denote by $\theta : \Omega \to \Omega$ the shift map $\omega = (\gamma_n)_{n\ge 1} \mapsto \theta(\omega) = (\gamma_{n+1})_{n\ge 1}$, and observe that the skew product map $\Theta(\omega,v) = (\theta(\omega), l_1(\omega) \cdot v)$ preserves the measure $m = \prob\times \nu$.
\end{step}

\begin{step}[Step 2]
We next show that $\prob\times \nu$ is ergodic under $\Theta$.

From (Kifer's version of) the Kakutani random ergodic theorem (see~\cite[Thm.~3.1]{furman}), it is sufficient to prove that for any Borel set $E$ such that $\nu(E\Delta \gamma E)=0$ for $\mu$-a.e. $\gamma$, then we have $\nu(E)=0$ or $1$.

Pick such a set $E$. First note that by stationarity and by the countability of $\supp(\mu)$, for $\mu$-a.e. $\gamma\in \Gamma$, we have $\gamma_*\nu\ll\nu$. In particular, we get $\gamma_*\nu(E\Delta \gamma E)=0$ so that
\[
\mu \ast (\nu|_E) = \int \gamma_*\nu \, \mathbf{1}_{\gamma (E)} \, \dmu(\gamma) = \int \gamma_*\nu \, \mathbf{1}_{E} \, \dmu(\gamma) =\nu|_E~.
\]
By the uniqueness of the $\mu$-stationary measure, we conclude that $\nu(E) =0$ or $1$, as required.
\end{step}

\begin{step}[Step 3]
Proof of~\eqref{eq:furstenberg} and~\eqref{eq301}.

The Birkhoff ergodic theorem applied to $\Theta$ and the function $\sigma(l_1(\omega),v)$, together with the cocycle relation~\eqref{eq:cocycle} now yield for $(\prob,\nu)$-a.e. $(\omega,v)$
\begin{equation}\label{eq001}
\lim_{n\to\infty} \frac1n \sigma(l_n(\omega),v) = \int \sigma (l_1(\omega),v) \, \dprob(\omega)\, \dnu(v) = \int \sigma(\gamma, v) \dmu(\gamma)\dnu(v).
\end{equation}
By Fubini's theorem, for $\prob$-a.e. $\omega$ we have that for $\nu$-a.e. $v$, the limit in~\eqref{eq001} exists. Since $\nu$ gives no mass to points, this holds for at least two distinct points so by Lemma~\ref{lem:two points} we get that for such $\omega$
\begin{equation}\label{eq002}
\lim_{n\to\infty} \frac1n \log \| l_n(\omega)\| = \lim_{n\to\infty} \frac1n \sigma(l_n(\omega),v)~.
\end{equation}
Thus, combining~\eqref{eq001} and~\eqref{eq002} we get~\eqref{eq:furstenberg} and~\eqref{eq301}.
\end{step}

\begin{step}[Step 4]
Positivity of the Lyapunov exponent.

To that end we rely on the following general lemma, a proof of which can be found in~\cite[Lem.~2.3, p.~22]{bougerol-lacroix}. For simplicity we write $F^+ = \max \{ 0, F\}$.

\begin{lemm}\label{lem:birkhoff}
Let $\Theta$ be a measure-preserving map on a probability space $(X, m)$. If $F:X\to \re$ is a measurable function such that $\int F^+ \dm <\infty$ and $\lim_{n\cv\infty} \sum_{i=0}^{n-1} F\circ \Theta^i = +\infty $ a.s., then $F\in L^1(X,m)$ and $\int F \,\dm >0$.
\end{lemm}

We apply the previous lemma to the function $F: (\omega, v) \mapsto \sigma(\ell_1(\omega), v)$ on $X=\om\times \berk$, and to the skew product map $\Theta(\omega,v) = (\theta(\omega), l_1(\omega) \cdot v)$. By Step 1 for $\prob\times \nu$-a.e. $(\omega,v)$ we have that
\[
\lim_{n\cv\infty} \sum_{i=0}^{n-1} F\circ \Theta^i (\omega,v) = \lim_{n\cv\infty} \sigma(l_n(\omega),v) = +\infty.
\]
Therefore Lemma~\ref{lem:birkhoff} yields
\begin{equation}\label{eqL}
\int \sigma(\gamma, v)\, \dmu(\gamma)\dnu(v) >0~,
\end{equation}
and the proof of Theorem~\ref{thm:positive} is complete.

\begin{proof}[Proof of Lemma~\ref{lem:Hk}]
Assume by contradiction that $\nu$ charges a point in $\PP^{1,\an}_k$. Then for every $\alpha>0$, the set $\set{x, \nu(\set{x})>\alpha}$ is finite, thus there exists an atom of maximal mass $\alpha_0$. It follows that the set $\set {x, \nu(\set{x})=\alpha_0}$ is finite and invariant under every element of $\supp(\mu)$, hence $\Gamma$-invariant, which is impossible since $\Gamma$ is non-elementary.

For the second assertion, we start by proving that $\nu(\hh_k) =0$. Observe first that any closed ball for the hyperbolic metric is also closed in $\PP^{1,\an}_k$ hence is a Borel set. Denote by $\B(x,R) = \{ d_{\hh}({\,\cdot\,}, x) < R\}$ (resp. $\overline{\B}(x,R) = \{ d_{\hh}({\,\cdot\,}, x) \le R\}$) the open (resp. closed) ball of radius $R$ relative to the hyperbolic metric. It is thus enough to prove that for every $x\in \hh_k$ and every $R>0$, $\nu(\overline{\B}(x,R)) = 0$.

By Lemma~\ref{lem:proximal semi group}, the semi-group generated by $\supp(\mu)$ contains a hyperbolic element $h$. Replacing $\mu$ by some $\mu^k$ if necessary (this does not affect the stationarity of $\nu$) we may assume $h\in \supp(\mu)$. Put $s= \sup_{x\in \hh_k} \nu(\B(x,R))$ and suppose by way of contradiction that $s>0$.

Since the series $\sum_{n\geq 0} \nu \lrpar{\overline{\B}(x_{\mathrm{g}}, (n+1)R)\setminus \B(x_{\mathrm{g}}, R)}$ converges, there exists $B$ such that for $d_\hh(x,x_{\mathrm{g}})> B$, $\nu(\B(x,R))\leq s/2$. Since $h$ is hyperbolic and is an isometry for $d_\hh$, we have $d_\hh(h^m(x_{\mathrm{g}}), x_{\mathrm{g}}) \to \infty$, hence there exists an integer $m$ such that $d_\hh(x,x_{\mathrm{g}}) \le B$ implies $d_\hh(h^m(x), x_{\mathrm{g}})\ge 2B$ hence $\nu (\B(h^m(x),R)) \le s/2$.

Put $\e = \mu^m(\set{h^m})$ and pick $x\in \hh_k$ such that $\nu(\B(x,R))\geq s(1-\e/3)$. In particular $d_{\hh_k}(x,x_{\mathrm{g}})\le B$. The invariance relation $\mu^m\ast \nu = \nu$ yields
\begin{align*}
s\lrpar{1-\frac{\e}{3}} \leq \nu(\B(x,R)) &= \sum_{\gamma \in \Gamma} \mu^m(\set{\gamma}) \nu(\gamma (\B(x,R))) \\
&\leq \sum_{\gamma \in \Gamma\setminus{h^m} } \mu^m(\set{\gamma}) \nu(\B(\gamma x,R)) +
\mu^m(\{h^m\}) \nu(\B(h^m(x), R))\\
&\leq (1- \mu^m(\{h^m\})) s + \mu^m(\{h^m\}) \frac{s}{2} = s \lrpar{1-\frac{\e}{2}}.
\end{align*}
From this contradiction we conclude that $\nu(\hh_k)=0$.

To prove that $\nu$ gives full mass to $\PP^1(k)$, consider the projection $\pi$ on the closed subtree $\overline{\mathrm{Conv}(\PP^1(k))} = \hh_k \cup \PP^1(k)$. Since $\hh_k \cup \PP^1(k)$ is $\Gamma$-invariant, for any $x\in \PP^{1, \an}_k$ and any $\gamma\in \Gamma$ we have that $\pi(\gamma(x)) = \gamma(\pi(x))$. It follows that $\pi_*\nu$ is stationary. By the first part of the proof, $\pi_*\nu$ gives full mass to $\PP^1(k)$ (which is a Borel set). Hence $\nu$ gives full mass to $\pi\inv(\PP^1(k)) = \PP^1(k)$. This completes the proof.
\end{proof}
\end{step}

\subsection{Distribution of attracting and repelling fixed points}

In the course of the proof of our main result we shall need the following interpretation of the stationary measures as the distribution of fixed points of hyperbolic elements.

Recall that the dual measure $\check{\mu}$ is defined as the image of $\mu$ under the involution $\gamma\mapsto \gamma^{-1}$, and that the measure $\mu$ satisfies~\eqref{B1} (resp.~\eqref{B2}) iff $\check{\mu}$ does. We also set $\check{\omega} = (\gamma_n^{-1})$ when $\omega = (\gamma_n)$ so that $l_n(\omega) = r_n(\check{\omega})^{-1}$. Observe that $\|l_n(\omega)\| = \|r_n(\check{\omega})^{-1}\| = \|r_n(\check{\omega})\|$ hence $\chi(\check{\mu}) = \chi(\mu)$.

\begin{theo}\label{thm:loxodromic}
Let $\mu$ be a probability measure with countable support in $\SLk$, satisfying~\eqref{B1}. Then
\begin{equation}\label{eq:proba hyperbolic}
\prob\lrpar{\left\{ \omega, \, r_n(\omega)
\text{ is hyperbolic}\, \right\}}
\underset {n\cv\infty}\longrightarrow 1.
\end{equation}
In addition, the asymptotic distribution for the weak-$\ast$ topology on $\berk$ of the attracting (resp. repelling) fixed point of $r_n(\omega)$ is given by the unique $\mu$-stationary (resp. $\check{\mu}$-stationary) probability measure $\nu$ (resp. $\check{\nu}$).

If furthermore $\mu$ satisfies~\eqref{B2}, then for every $\e>0$
\begin{equation}\label{eq:trace}
\prob\lrpar{\set{
\abs{ \unsur{n}\log\abs{\tr (r_n(\omega))} - \chi(\mu)} <\e }} \underset {n\cv\infty}\longrightarrow 1.
\end{equation}
\end{theo}

\begin{rema}\label{rem:301}
The meaning of the statement on the distribution of periodic points is the following. For each $n$ let $\Omega_n$ be the set of $\omega\in \Omega$ such that $r_n(\omega)$ is hyperbolic. One can then define the measurable map $\Att^n, \Rep^n: \Omega_n \to \PP^1(k)$ by sending $\omega$ to the attracting and repelling fixed points of $r_n(\omega)$. The theorem asserts that $\prob(\Omega_n) \to1$ when $n\to\infty$, and $\Att^n_* \prob \to \nu$, and $\Rep^n_* \prob \to \check{\nu}$ as $n\to\infty$.
\end{rema}

\begin{rema}\label{rem:302}
Under stronger moment assumptions on $\mu$, \eqref{eq:trace} can be turned
into an almost sure limit. For instance it is shown in~\cite[Thm.~1.4]{maher-tiozzo} that if $\mu$ has bounded support the probability in~\eqref{eq:proba hyperbolic} is exponentially close to $1$, thus by Borel--Cantelli, $\unsur{n}\log\abs{\tr (r_n(\omega))}$ converges a.s. to $\chi(\mu)$.
\end{rema}

\begin{proof}
Consider the probability measure $\nu\times\check{\nu}$ on the product space $\PP^1(k) \times \PP^1(k)$. This measure is the image of $\prob$ under the measurable map $\omega \mapsto (e(\omega), e(\check{\omega}))$.

By Fubini, and since the measures $\nu$ and $\check{\nu}$ do not have any atom, we have $\nu\times\check{\nu} (\Delta) = 0$ where $\Delta$ denotes the diagonal in $\PP^1(k) \times \PP^1(k)$. It follows that
\begin{equation}\label{eq333}
(\nu\times\check\nu)\set{(x,y)\in \pu(k)^2, \ d_{\sph}(x,y)\leq \delta} \underset {\delta\cv0}\longrightarrow 0~.
\end{equation}
By Theorem~\ref{thm:partial unique} $\| r_n(\omega)\|\to\infty$ a.s. and the asymptotic distribution of the $B_\att(r_n(\omega))$ is given by $\nu$. Similarly the asymptotic distribution of the $B_\rep(r_n(\omega))$ is given by $\check{\nu}$, so that for each $\delta>0$, we infer that
\begin{multline*}
\limsup_{n\cv\infty} \prob\left\{d_{\sph}(B_\att(r_n(\omega)), B_\rep(r_n(\omega)))\leq \delta\right\} \\
\leq (\nu\times\check\nu)\set{(x,y)\in \pu(k)^2, \ d_{\sph}(x,y)\leq \delta},
\end{multline*}
which can be made as small as we wish.

Fix any real number $\varepsilon>0$, and choose $\delta>0$ such that the left hand side in~\eqref{eq333} is at most $\e/2$. Then there exist $N = N(\e)$ and a set $\Omega_\varepsilon$ with $\prob(\Omega_\e)\geq 1-\e/2$, such that if $\omega \in \Omega_\e$ and $n\ge N(\omega)$, we have
\[
d_{\sph}(B_\att(r_n(\omega)), B_\rep(r_n(\omega))) > \delta~.
\]
In addition $\| r_n(\omega)\|\to\infty$ a.s, so increasing $N$ and discarding a set of probability $\e/2$ we may further assume that
\[
\| r_n(\omega) \|\ge 2\delta\inv \text{ on } \Omega_\varepsilon~.
\]
Now pick $\omega \in \Omega_\varepsilon$ and $n\ge N$. The two closed balls $B_\att(r_n(\omega)), B_\rep(r_n(\omega))$ have diameter at most $ \delta/2$ hence are disjoint. Since $r_n(\omega)$ maps $\PP^{1,\an}_k\setminus B_\rep(r_n(\omega))$ into $B_\att(r_n(\omega))$, we conclude by Lemma~\ref{lem:distinct balls}.

If furthermore~\eqref{B2} holds then from Lemma~\ref{lem:expand vector 2}, we conclude that
\[
\left| \log |\tr(r_n(\omega)| - \log \| r_n(\omega)\| \right|
\le - \log \delta~,
\]
and~\eqref{eq:trace} follows since $\lim_{n\to\infty} \frac1n \log\|r_n(\omega)\| = \chi(\check{\mu})= \chi(\mu)$.
\end{proof}


\section{Degenerations: non elementary representations}\label{sec:non elem}

In this section we fix a finitely generated group $G$, endowed with some probability measure $\sm$. Recall the notation $\Laurent = \cdt$ and that $\mm$ is the ring of holomorphic functions in $\dD^*$ with meromorphic extension at the origin. We fix a representation $\rho: G\to \SL(2, \mm)$, that is a family of representations $\rho_t : G \to \SLC$ for $t\in \D^*$ such that for any $g\in G$ $ t\mapsto \rho_t(g)$ is holomorphic on $\D^*$ and extends meromorphically through the origin. We suppose that
\begin{enumerate}\renewcommand{\theenumi}{A\arabic{enumi}}
\item \label{A1b} $\supp(\sm)$ generates $G$;
\item \label{A2b} $\sm$ has finite first moment $\int \length(g) \dsm(g)<\infty$.
\end{enumerate}

We also assume that if $\rho_\na$ denotes the induced representation $\rho_\na: G \to \SL(2,\Laurent)$, then $\rho_\na(G)$ is non-elementary. Our aim is to prove Theorem~\ref{Thm:atomic} and to infer the non-elementary case of Theorem~\ref{Thm:lyapunov na}.

\subsection{Basic remarks on meromorphic families of representations}

Since we have to deal with both Archimedean and non-Archimedean objects, we slightly change notation: $\norm{{\,\cdot\,}}$ denotes the operator norm in $\SLC$ associated to any norm on $\cd$, say $\norm{(x,y)} =\max({\abs{x}, \abs{y}})$, and $\norm{{\,\cdot\,}}_{\na}$ denotes the non-Archimedean norm on $\SL(2, \Laurent)$ associated to the $t$-adic norm in $\mm$ given by $\abs{f}_\na = \exp(-\ord_{t=0}(f))$. More generally we use the subscript $\na$ to label non-Archimedean objects.

For any $t\in \D^*$, we set $ \Gamma_t =\rho_t(G)$ and let $\mu_t$ be the push-forward of $\sm$ under $\rho_t$. Analogously, denoting by $\rho_\na:G\to \SL(2, \Laurent)$ the non-Archimedean representation naturally associated to $\rho$, we let $\Gamma_\na = \rho_\na(G)$ and $\mu_\na = (\rho_\na)_*\sm$.

Observe that for every $g\in G$, $\log \norm{\rho(g)}_\na\leq C \length (g)$ for some uniform constant $C>0$, and likewise for $\log\norm{\rho_t(g)}$. In particular we have

\begin{lemm}\label{lem:AtoB}
The condition~\eqref{A2b} implies the moment condition~\eqref{B2} for the measures $\mu_t$ and $\mu_\na$.
\end{lemm}

We will need some uniformity on the control of $\norm{\rho_t(g)}$.

\begin{lemm}\label{lem:uniform disk}
For every homomorphism $\rho: G\to \SL(2, \mm)$ there exists $C = C(\rho)>0$ such that for every $g\in G$ one can write $\rho_t(g) = t^{-\log\|\rho(g) \|_\na}\cdot \widetilde{\gamma}(t)$, with $\widetilde{\gamma}\in \SL(2,\cO(\D))$, and
\begin{equation}\label{eq:uniform disk}
\abs{\log\norm{\widetilde \gamma}_{L^{\infty}(\overline{\D}(0,1/2))} }\leq C\, \mathrm{length}(g).
\end{equation}
\end{lemm}

\begin{proof}
let $(s_i)$ be a finite symmetric set of generators of $G$ and write $g$ as a reduced word in $G$, $g = s_{i_1}\cdots s_{i_n}$, $n = \mathrm{length(g)}$. Let $\sigma_i = \rho(s_i)$ and write $\sigma_i = t^{-\alpha_i} \widetilde \sigma_i$ with $\alpha_i = \log\norm{\sigma_i}_\na$ and $\widetilde \sigma_i$ holomorphic and non vanishing at 0,
so that
\[
\rho_t(g)= t^{-\sum_{j=1}^n\alpha_{i_j}} \; \widetilde \sigma_{i_1}\cdots \widetilde \sigma_{i_n}~.
\]
Set $A:= \max \alpha_i$, and write $\rho_t(g) = t^{-\log\|\rho(g) \|_\na}\cdot \widetilde{\gamma}(t)$ with $\widetilde{\gamma} \in \SL(2, \cO(\D))$. Then we get that $\widetilde{\gamma}(t) = t^{-\alpha} \cdot \widetilde \sigma_{i_1}\cdots \widetilde \sigma_{i_n}(t)$ with $0\le \alpha = \sum_{j=1}^n\alpha_{i_j} - \log\|\rho(g) \|_\na \le An$. By using the maximum principle we can estimate
\begin{multline*}
\sup_{|t|\le1/2}|\widetilde{\gamma}(t)| \le \sup_{|t|=1/2} |\widetilde{\gamma}(t)| \le (1/2)^{-\alpha} \prod_{j=1}^n \lrpar {\sup_{|t|=1/2} \abs{\widetilde \sigma_{i_j}(t) }}\\
\le 2^{An} \, \left(\max_{i}\norm{ \widetilde \sigma_i}_{L^\infty(D(0, 1/2))} \right)^n\le D^n
\end{multline*}
for some $D\geq1$. Likewise we have that
\begin{align*}
\sup_{|t|\le1/2}|\widetilde{\gamma}(t)| &\ge \sup_{|t|=1/2} |\widetilde{\gamma}(t)|
\ge (1/2)^{-\alpha} \sup_{|t|=1/2} \abs{\widetilde \sigma_{i_1}\cdots \widetilde \sigma_{i_n}(t)}\\
&\ge (1/2)^{-\alpha} \abs{\widetilde \sigma_{i_1}\cdots \widetilde \sigma_{i_n}(0)} \ge E^n
\end{align*}
for some $E>0$ and we are done.
\end{proof}

Let us also note the following basic but crucial observation.

\begin{lemm}\label{lem:specialization NE}
If $\Gamma_\na$ is non-elementary, then for small enough $t\in \dD^*$, $\Gamma_t\leq \SLC$ is non-elementary.
\end{lemm}

\begin{proof}
By definition $\Gamma$ contains two hyperbolic elements $\gamma_1$ and $\gamma_2$ with disjoint fixed points on $\PP^{1,\an}_{\Laurent}$. Since for $i=1,2$ $\abs{\tr (\gamma_i)}_\na>1$, it follows that $\abs{\tr(\gamma_{i,t})} \to \infty$ as $t\cv 0$. Thus $\gamma_{1,t}$ and $\gamma_{2,t}$ are loxodromic for small $t$, and they have well-defined attracting and repelling fixed points $\att(\gamma_{i,t})$, $\rep(\gamma_{i,t})$. Saying that $\gamma_1$ and $\gamma_2$ have disjoint fixed points sets on $\berkl$ implies that the formal expansions of the curves $t\mapsto \att(\gamma_{i,t})$ and $t\mapsto \rep(\gamma_{i,t})$ are all distinct, thus the corresponding points in $\pu$ must be disjoint in some punctured neighborhood of the origin.
\end{proof}

\subsection{Models and \texorpdfstring{$\PP^{1,\an}_{\Laurent}$}{PP1, an Laurent}} \label{subs:models}

In this paragraph we review the notion of model. Set $X = \dD\times \PP^1_\cc$. A (bimeromorphic) model of $X$ is a surface $Y$ together with a bimeromorphic holomorphic map $\pi_Y:Y\to X$ that is biholomorphic above $X\setminus (\{0\}\times \pu)$. The fiber $\pi_Y\inv(\set{t}\times \pu)$ will be denoted by $Y_t$. We will only consider the case where $Y$ is smooth, in which case $\pi_Y$ is simply a composition of point blow-ups above the central fiber and $Y_0$ is a divisor with simple normal crossings.

We say that a model $Y'$ dominates a model $Y$ if the birational map $\pi_{Y'}:Y'\cv X$ factors through $Y$. The set of models is then directed in the sense that given two models $Y$ and $Y'$, there exists a third one $Y''$ dominating both.
%\medskip

We now explain the basic correspondence between models and finite subsets of $\berkl$. Recall as a set, $\PP^{1,\an}_{\Laurent}$ is the one point compactification of the space of multiplicative semi-norms on $\Laurent[z]$ whose restriction to $\Laurent$ is the $t$-adic norm. Let $Y$ be any model, and pick any irreducible component $E$ of $Y_0$: we will define a type 2 point $\zeta_E \in \PP^{1,\an}_{\Laurent}$.

Observe that any element $f\in \cc(t)[z]$ defines a rational function on $\D \times \PP^1_\cc$ so that we can define
\[
\abs{f}_{\zeta_E} = |f(\zeta_E)| = \exp \left(- \frac1{b_E}\, \ord_E(f \circ \pi_Y)\right)
\]
where $b_E = \ord_E (t \circ \pi_Y)$. Dividing by $b_E$ guarantees that $\zeta_E\rest{\Laurent} = \abs{{\,\cdot\,}}$ and also that the definition is model-independent in the sense that if $Y'$ dominates $Y$ and $E'$ is the strict transform of $E$ then $\zeta_E=\zeta_{E'}$. Since the field $\cc(t)$ is dense in $\Laurent$, it follows that $\zeta_E$ extends uniquely to a semi-norm on $\Laurent[z]$ hence defines a point in $\PP^{1,\an}_{\Laurent}$. This point is of type 2 and is defined over the field extension $\cc(\!(t^{1/{b_E}})\!)$. Conversely, any type 2 point of $\PP^{1,\an}_{\Laurent}$ is equal to $\zeta_E$ for some irreducible component $E$ of the central fiber of some model $Y$ over $X$, see~\cite[Lem.~7.16]{NL}.

We denote by $ S(Y)$ the set of all type 2 points $\zeta_E \in \PP^{1,\an}_{\Laurent}$ where $E$ ranges over the set of irreducible components of $Y_0$. It is an elementary fact that $S(Y') \supset S(Y)$ if and only if $Y'$ dominates $Y$.

\begin{prop}[{see~\cite[\S 4.1]{demarco-faber-2}}]
Let $S$ be any finite set of type 2 points in $\PP^{1,\an}_{\Laurent}$. Then there exists a (smooth) model $Y$ such that $S \subset S(Y)$.
\end{prop}

\begin{rema}
In fact there is an isomorphism of partially ordered sets between finite sets of type 2 points endowed with the inclusion and proper bimeromorphic maps $\pi: Y \to \D \times \PP^1_\cc$ with $Y$ a normal complex analytic variety, see~\cite[Thm.~7.18]{NL}.
\end{rema}

The previous proposition together with Proposition~\ref{prop:subtree Hk} yield the following corollary.

\begin{coro}\label{cor:cut to small}
Let $\nu$ be any probability measure on $\berkl$ which is non atomic and gives zero mass to $\hh_\Laurent$. Then for every $\e>0$ there exists a model $Y$ such that every connected component $U$ of $\berkl\setminus S(Y)$ satisfies $\nu(U)<\e$.
\end{coro}

\begin{rema}\label{rem:useless}
Suppose that $\nu$ gives full mass to the set $\PP^{1,\an}(\Laurent)$ of points of type $1$ defined over $\Laurent$. Then one can actually choose $\pi: Y \to \D\times \PP^{1}_\cc $ to be a composition of blow-ups at \emph{free} points, i.e. lying in the regular locus of the central fiber. In other words, one can choose the central divisor $\pi^* (\PP^{1}_\cc \times \{ 0 \})$ to be \emph{reduced}.
\end{rema}

\subsection{Reduction map and action of \texorpdfstring{$ \SL(2,\mm)$}{SL(2, mm)}}\label{subs:reduction}

Fix any model $Y$, and let $\zeta$ be any type 2 point in $\PP^{1,\an}_{\Laurent}$. Then we can find a model $Y'$ dominating $Y$ and an irreducible component $E$ of $Y'_0$ such that $\zeta = \zeta_E$. If the natural map $\pi_{Y',Y}: Y' \to Y$ contracts $E$ to a point, we set $\red_Y(\zeta_E) = \pi_{Y',Y}(E)$. Otherwise we let $\red_Y(\zeta_E)$ be the generic point of the curve $\pi_{Y',Y}(E)$ which is a (non-closed) point in the $\cc$-scheme $Y_0$. The mapping $\red_Y$ is called the \emph{reduction map}. It extends canonically to an anti-continuous map $\red_Y: \PP^{1,\an}_{\Laurent}\to Y_0$ (i.e. the preimage of a closed set in the Zariski topology is open for the Berkovich topology), see e.g.~\cite[\S 5.2.4]{temkin}, or~\cite[\S 4.2]{demarco-faber-2}.

It can be shown that if $\eta_E$ is the generic point of an irreducible component $E$ of $Y_0$, then $\red_Y^{-1}(\eta_E) = \{ \zeta_E\}$. If $p\in Y_0$ is a closed point then $\red_Y^{-1}(p)$ is a connected component of $\berkl\setminus {S}(Y)$ whose boundary consists of the points $\zeta_E$ where $E$ ranges over all irreducible components of $Y_0$ containing $p$. In particular the boundary of $\red_Y^{-1}(p)$ consists of one or two points.

The reduction map behaves well under proper modifications.

\begin{lemm}\label{lem:reduction modification}
If $\pi: Y'\to Y$ is a birational morphism, then $\red_Y = \pi\circ \red_{Y'}$.
\end{lemm}

\begin{proof}
If $\zeta$ is any type 2 point, there exists a model $Y''$ dominating $Y'$ such that $\zeta = \zeta_E$ for some component $E$ of $Y''_0$. Then it follows immediately from the definitions that $\red_Y (\zeta)= \pi\circ \red_{Y'}(\zeta)$. This identity then extends to $\berkl$ because $\red_Y$ and $\red_{Y'}$ admit unique anti-continuous extensions to $\berkl$ and $\pi$ is continuous for the Zariski topology.
\end{proof}

Let us now pick $\gamma \in \SL(2, \mm)$, and denote by $\gamma_\na$ its natural image in $\SL(2, \Laurent)$. Observe that $\gamma$ induces a biholomorphism from $\dD^*\times \PP^1_\cc$ to itself commuting with the first projection, and extending meromorphically to $\dD \times \PP^1_\cc$. More generally given any two models $Y,Y'$ this biholomorphism extends to a bimeromorphic map $\gamma_{Y,Y'}: Y \dashrightarrow Y'$. Its properties can be described from $\gamma_\na$ and the reduction map as follows.

\begin{prop}\label{prop:val-geom}
Let $Y,Y'$ be two models over $X$, and pick $\gamma\in \SL(2,\mm)$.
\begin{enumerate}
\item The induced bimeromorphic map $\gamma_{Y,Y'}: Y \dashrightarrow Y'$ has an indeterminacy point at $p\in Y_0$ iff there exists a type 2 point $\zeta \in \red^{-1}_{Y}(p)$ such that $\gamma_\na(\zeta) \in S(Y')$.
\item Suppose $\gamma_{Y,Y'}$ is holomorphic at $p\in Y_0$. Then for any $\zeta \in \red_{Y}^{-1}(p)$, we have $\gamma_{Y,Y'}(p) = \red_{Y'} (\gamma_\na(\zeta))$.
\end{enumerate}
\end{prop}

\begin{proof}
The proposition follows easily from the basic properties of the reduction map together with the following lemma.
\end{proof}

\begin{lemm}
Let $\zeta\in\berkl$ be a type 2 point. Fix a model $Y_1$ dominating $Y$ and a component $E$ of $(Y_1)_0$ such that $\zeta = \zeta_E$. Let $\zeta' = \gamma_\na(\zeta)$, and fix a model $Y'_1$ dominating $Y'$ with a component $E'$ of $(Y'_1)_0$ such that $\zeta' = \zeta_{E'}$. Then $\gamma_{Y_1, Y'_1}(\eta_E) = \eta_{E'}$, where $\eta_E$ (resp. $\eta_{E'}$) is the generic point of $E$ (resp. $E'$).

Conversely, given any two models $Y_1$, $Y_1'$ and respective irreducible components $E \subset (Y_1)_0$ and $E' \subset (Y'_1)_0$, if $\gamma_{Y_1, Y'_1}(\eta_E) = \eta_{E'}$ then $\zeta_{E'} = \gamma_\na(\zeta_{E})$.
\end{lemm}

\begin{proof}
The pull-back by $\gamma$ of any rational function on $Y'_1$ vanishing at $\eta_{E'}$ necessarily vanishes at $\eta_E$ since $\zeta_{E'} = \gamma_\na(\zeta_E)$. Since for any point $p$ not lying on $E'$ there exists a rational function on $Y'_1$ which is non-zero at $p$, and zero on $E'$, the first claim follows.

For the second claim, pick $f$ a rational function on $Y'_1$, and choose points $p\in E$, $p'\in E'$ such that $\gamma_{Y_1, Y'_1}$ is regular at $p$, $p' = \gamma_{Y_1, Y'_1}(p)$, and $E$ (resp. $E'$) is regular at $p$ (resp. at $p'$). Then we may choose coordinates $(x,y)$ at $p$ and $(x',y')$ at $p'$ such that $E = \{ x=0\}$, $E' = \{ x'=0\}$ and $\gamma_{Y_1, Y'_1} (x,y) = (x^k, \star)$. It follows that
\[
-\log|f(\gamma_\na(\zeta_{E}))| = \frac1{b_E}\ord_{E}(f \circ \gamma_{Y_1, Y'_1}) = \frac{k}{b_E}\ord_{E'}(f)
\]
which implies $ b_{E'} = b_E/k$ and $\zeta_{E'} = \gamma_\na(\zeta_{E})$.
\end{proof}

Using the reduction map, one can see that for large $\norm{\gamma}_\na$, the meromorphic map $\gamma_Y$ acts on the central fiber $Y_0$ like a (brutal) North-South transformation.

\begin{prop}\label{prop:nordsud}
Let $Y$ be any model.

There exists a constant $C=C(Y)$ depending only on $Y$ such that for any $\gamma \in \SL(2,\mm)$ such that $\norm{\gamma}_\na\geq C$, there exist two points $\mathrm{att}(\gamma_Y)$ and $\mathrm{rep}(\gamma_Y)$ in $Y_0$ (not necessarily distinct) such that the induced bimeromorphism $\gamma_Y: Y \dashrightarrow Y$ is holomorphic on $Y_0\setminus \mathrm{rep}(\gamma_Y)$ and $\gamma_Y(Y_0\setminus \mathrm{rep}(\gamma_Y)) = \mathrm{att}(\gamma_Y)$.
\end{prop}
\begin{proof}
Choose $C > \max_{\zeta \in S(Y)} \exp(d_\hh (\zeta, x_{\mathrm{g}}))$, and pick $\gamma$ such that $\norm{\gamma}_\na\geq C$.

By Proposition~\ref{prop:B att rep} there exist two disjoint closed Berkovich disks $B_{\att}(\gamma)$ and $B_{\rep}(\gamma)$ of spherical diameter $\norm{\gamma}^{-1}_\na$ such that $\gamma (\PP^{1,\an}_{\Laurent}\setminus B_{\rep}(\gamma)) \subset B_{\att}(\gamma)$.

Observe that $ B_{\rep}(\gamma)$ cannot contain any point $\zeta \in S(Y)$ since otherwise we would get that $d_\hh (\zeta, x_{\mathrm{g}}) \ge \log \norm{\gamma}_\na^{-1}$, thus contradicting the choice of $C$. Thus, $ B_{\rep}(\gamma) \cap S(Y) = \emptyset$, and since $B_{\rep}(\gamma)$ is connected, it is contained in a connected component of $\PP^{1,\an}_{\Laurent}\setminus S(Y)$. It follows that there exists a unique closed point $\mathrm{rep}(\gamma_Y)\in Y$ such that $B_{\rep}(\gamma) \subset \red_Y^{-1}(\mathrm{rep}(\gamma_Y))$. Similarly we have $B_{\att}(\gamma) \cap S(Y) = \emptyset$, and $B_{\att}(\gamma)\subset \red_Y^{-1}(\mathrm{att}(\gamma_Y))$ for some unique $\mathrm{att}(\gamma_Y)\in Y$.

Now pick any point $p \in Y_0$ different from $\mathrm{rep}(\gamma_Y)$. Then $\red_Y^{-1}(p)$ is disjoint from $B_{\rep}(\gamma)$ so it is mapped into $B_{\att}(\gamma)$ by $\gamma_\na$. The first item of Proposition~\ref{prop:val-geom} then asserts that $\gamma_{Y}$ is holomorphic at $p$ and the second one that $\gamma_{Y}(p) =\red_Y(B_{\att}(\gamma)) = \mathrm{att}(\gamma_Y)$. The proof is complete.
\end{proof}

\begin{rema}\label{rem:more precise}
The proof shows that $\mathrm{rep}(\gamma_Y) = \red_Y (\zeta)$ for any $\zeta \in B_{\rep}(\gamma)$. In particular if $\gamma$ is hyperbolic and $\zeta_{\mathrm{rep}}$ is its repelling fixed point, then $\mathrm{rep}(\gamma_Y) = \red_Y (\zeta_{\mathrm{rep}})$. Also if $S(Y)$ contains the Gau{\ss} point then taking $C > \exp (\mathrm{diam}_\mathbb{H} (S(Y)))$ is enough.
\end{rema}

Let $\nu$ be any Radon measure on the Berkovich projective line $\PP^{1,\an}_{\Laurent}$. Then by definition the \emph{residual measure} $(\red_Y)_* \nu$ is the atomic measure on $Y_0$ satisfying
\[
{(\red_Y)_* \nu} (\{ p \}) = \nu (\red_Y^{-1}(p))~,
\]
for any closed point $p\in Y_0$. Note that if $\nu$ gives no mass to $S(Y)$, then $(\red_Y)_* \nu$ is the push forward of $\nu$ in the usual sense. Since the union of the open sets $\red_Y^{-1}(p)$ as $p$ ranges through closed points of $Y_0$ is equal to the complement of $S(Y)$ in $\PP^{1,\an}_{\Laurent}$, it follows that the total mass of $(\red_Y)_* \nu$ is equal to the mass of the restriction of $\nu$ to $\PP^{1,\an}_{\Laurent}\setminus S(Y)$.

\begin{lemm}\label{lem:cvg mass}
If $\nu_n$ converges weakly to $\nu$ and $\nu(S(Y)) =0$, then $(\red_Y)_* \nu_n$ converges in mass to $(\red_Y)_* \nu$.
\end{lemm}

\begin{proof}
Pick any closed point $p\in Y_0$. Recall that Radon measures are regular. Since $\nu_n \to \nu$ and $\red_Y^{-1}(p)$ is open we deduce that
\[
\liminf_{n\to\infty} (\red_Y)_* \nu_n \{ p \} \ge (\red_Y)_* \nu \{ p \}~.
\]
On the other hand, since the boundary of $\red_Y^{-1} \{ p \}$ is included in $S(Y)$, the assumption $\nu(S(Y)) =0$ implies that
\[
\limsup_{n\to\infty} (\red_Y)_* \nu_n (\{ p \}) \le \nu \left(\overline{ \red_Y^{-1} \{ p \}} \right) = (\red_Y)_* \nu (\{ p \})~.
\]
Therefore we conclude that for every $p\in Y_0$, $(\red_Y)_* \nu_n(\{p \}) \to (\red_Y)_* \nu (\{p\})$. Since all these measures are atomic and of uniformly bounded mass, the result follows.
\end{proof}

\subsection{Proof of Theorem~\ref{Thm:atomic}}\label{subs:atomic}

Let $(G, \sm)$ satisfying~\eqref{A1}, and a representation $\rho : G \to \SL(2, \mm)$, and suppose that the induced representation $\rho_\na: G \to \SL(2, \Laurent)$ is non-elementary. Lemma~\ref{lem:specialization NE} implies that $\rho_t: G \to \SLC$ is also non-elementary for small $t$ so by Theorem~\ref{thm:positive} (in the complex case) it makes sense to talk about the unique probability measure $\nu_t$ on $\PP^1(\cc)$ that is stationary under $\mu_t = (\rho_t)_* \sm$. We also denote by $\nu_\na$ the unique probability measure on $\PP^{1,\an}_{\Laurent}$ that is stationary under $\mu_\na = (\rho_\na)_* \sm$.

Fix a model $Y$. Since $\pi_Y: Y \to \D \times \PP^1_\cc$ is a biholomorphism outside the central fiber, we can view $\nu_t$ as a probability measure on $Y_t$, which we denote by $\nu_{Y_t}$. Our aim is to prove that $\nu_{Y_t}$ converges as $t\to 0$ to an atomic measure on the central fiber given by $(\res_Y)_* \nu_\na$.

Fix $\e>0$, and choose a model\footnote{Observe that Remark~\ref{rem:useless} applies here so that we can further assume the central fiber to be reduced.} $Y'$ such that the conclusion of Corollary~\ref{cor:cut to small} holds for the unique stationary probability measure $\check\nu_{\sf na}$ associated to the reversed random walk. We can assume that $Y'$ dominates $Y$, and write $\nu_{Y'_t}$ for the stationary measure on $Y'_t$.

Consider the Markov operator $P_t = \int (\gamma_t)_* \dmu_t(\gamma)$ acting on the space of probability measures on $Y'_t$, whose $n$-fold iterate is $P_t^n = \int (\gamma_t)_* \dmu_t^n(\gamma)$. Observe that for every $n\geq 1$ we have that $P_t^n\nu_{Y'_t}= \nu_{Y'_t}$. To analyze this identity as $t$ tends to 0, we extract a sequence $t_j\to0$ such that $\nu_{Y_{t_j}}$ and $\nu_{Y'_{t_j}}$ converge to respective probability measures $\nu_{Y_0}$ on $Y_0$, and $\nu_{Y'_0}$ on $Y_0'$. Note that by construction the push-forward of $\nu_{Y'_0}$ under the canonical projection map $Y' \cv Y$ is $\nu_{Y_0}$. We will show that $\nu_{Y_0} = (\res_Y)_* \nu_\na$.

Let $C = C(Y')$ be the constant given by Proposition~\ref{prop:nordsud}, and define $A = A(C)= \set{\gamma\in \SL(2, \Laurent), \ \norm{\gamma}\geq C}$. We define two measurable maps on $A_n$ with values in $Y'_0$, namely $\Att_{Y'}(\gamma):= \mathrm{att}(\gamma_{Y'})$ and $\Rep_{Y'}(\gamma):= \mathrm{rep}(\gamma_{Y'})$. The image $\mu^n_\na|_{A_n}$ under $\Rep_{Y'}$ (resp. $\Att_{Y'}(\gamma)$) is by definition the distribution of repelling (resp. attracting) points on $Y'$ for the random walk at time $n$.

\begin{lemm}\label{lem:distrib rep-att}
The sequence of atomic measures $(\Rep_{Y'})_*(\mu^n_\na|_{A}) $ (resp.\linebreak $(\Att_{Y'})_*(\mu^n_\na|_{A}) $) converges in mass to $(\res_{Y'})_*\check\nu_\na$ (resp. to $(\res_{Y'})_*\nu_\na$).
\end{lemm}

Taking this lemma for granted for the moment, let us complete the proof of the theorem. Observe first that for any $\gamma \in A$, from the description of the action of $\gamma$ in Proposition~\ref{prop:nordsud}, we see that any cluster value of $(\gamma_{t_j})_*\nu_{Y'_{t_j}}$ is of the form
\[
\big(1- \nu_{Y'_0}\lrpar{\set{\rep(\gamma_{Y'})}}\big) \delta_{\att(\gamma_{Y'})} + \mathrm{error}
=
\delta_{\att(\gamma_{Y'})} + \mathrm{error}
\]
where the error on the right hand side is a signed measure of total mass bounded by $ 2 \nu_{Y'_0}\lrpar{\set{\rep(\gamma_{Y'})}}$. Using the identity $\nu_{Y'_t} = P_t^n\nu_{Y'_t}$ and letting $t=t_j\cv0$ we infer that
\begin{equation}\label{eq:limit stationary}
\nu_{Y'_0} = \int_{A} \delta_{\att(\gamma_{Y'})} \dmu_\na^n(\gamma) + \mathrm{error}
\end{equation}
where the mass of the error is
\begin{align}
\notag \m(\mathrm{error}) & \leq 2 \int_{A} \nu_{Y'_0}\lrpar{\set{\rep(\gamma_{Y'})}} \dmu_\na^n(\gamma) +
\lrpar{1 -\mu_\na^n(A)} \\\label{eq:error}
&= 2 \int_{x\in Y'_0} \nu_{Y'_0}(\{ x \}) \dd \!\left((\Rep_{Y'})_*(\mu_\na^n|_{A})\right)+ \lrpar{1 -\mu_\na^n(A)}.
\end{align}
By Lemma~\ref{lem:distrib rep-att}, $(\Rep_{Y'})_*(\mu_\na^n|_{A})$ converges in mass towards $(\res_{Y'})_*\check\nu_\na$. Since every atom of the latter measure has mass $\le \varepsilon$ by construction, we get that every atom of $(\Rep_{Y'})_*(\mu^n|_{A})$ has mass $\le 2\varepsilon$ for $n$ large enough. It follows that the integral in~\eqref{eq:error} is bounded by $\le 2 \varepsilon$. Since in addition $\mu^n(A) \to 1$ by Lemma~\ref{lem:contraction}, we conclude that for $n$ large enough
\begin{equation}
\nu_{Y'_0} - \int_{A} \delta_{\att(\gamma_{Y'})} \dmu_\na^n(\gamma)
\end{equation}
has total mass $\le 5 \varepsilon$.

To conclude, we observe that applying Lemma~\ref{lem:distrib rep-att} again, the sequence of measures
\[
\int_{A} \delta_{\att(\gamma_{Y'})} \dmu_\na^n(\gamma) = \int_{x\in Y'_0} \delta_x \dd\!\left((\Att_{Y'})_*(\mu_\na^n|_{A})\right)
\]
converges in mass to $(\res_{Y'})_*\nu$. We thus infer that $ \nu_{Y'_0} - (\res_{Y'})_*\nu_\na$ has total mass $\le 5 \varepsilon$. Pushing down this information to $Y$ we get the same bound for $\nu_{Y_0} - (\res_{Y})_*\nu_\na$ on $Y$. Since $Y$ does not depend on $\e$ and $\e$ can be made arbitrarily small, we conclude that $\nu_{Y_0} = (\res_{Y})_*\nu_\na$, as required.
%\qed

\begin{proof}[Proof of Lemma~\ref{lem:distrib rep-att}]
Let $(\Omega, \prob) = (\SL(2,\Laurent)^{\N^*}, \mu_\na^{\N^*})$, and consider the set
\[
\Omega_n = \{ \omega \in \Omega, \, r_n(\omega)\in \SL(2, \Laurent) \text{ is
hyperbolic}\}~.
\]
If $\Rep^n_\na : \Omega_n \to \PP^{1,\an}_{\Laurent}$ denotes the measurable map sending $\omega$ to the repelling fixed point of $r_n(\omega)$, then by Theorem~\ref{thm:loxodromic} $\prob(\Omega_n) \to 1$ and $(\Rep^n_\na)_* \prob \to \check{\nu}_\na$. Pushing forward this convergence by the residue map $\red_{Y'}$ and applying Lemma~\ref{lem:cvg mass}, we thus get that $(\red_{Y'})_* (\Rep^n_\na)_* \prob$ converges in mass to $(\red_{Y'})_*\check{\nu}_\na$.

Now define $\Omega'_n = \Omega_n \cap r_n^{-1}(A)$, which also satisfies $\prob(\Omega'_n) \to 1$ by Lemma~\ref{lem:contraction}. It follows from Remark~\ref{rem:more precise} that for $\omega \in \Omega'_n$, $\rep_{Y'}(r_n(\omega)) = \red_{Y'}(\Rep_\na^n (\omega))$, in other words, the repelling fixed point of $r_n(\omega)$ is mapped under $\red_{Y'}$ to $\rep_{Y'}(r_n(\omega))$. We thus obtain that
\[
(\red_{Y'})_*(\Rep^n_\na)_* \prob = (\Rep_{Y'})_*(\mu^n_\na|_{A}) + \mathrm{error}
\]
where the mass of the error tends to $0$ as $n\to\infty$. This completes the proof.
\end{proof}

\subsection{Proof of Theorem~\ref{Thm:lyapunov na} in the non-elementary case} \label{subs:asymptotics}

As in the previous section we work with a representation $\rho : G \to \SL(2, \mm)$ such that the induced representation $\rho_\na: G \to \SL(2, \Laurent)$ is non-elementary, and further assume that~\eqref{A1b} and~\eqref{A2b} hold. For $t\in \dD^*$ we set $\mu_t = (\rho_t)_* \sm$ and we also put $\mu_\na = (\rho_\na)_* \sm$. Recall that for $t\neq 0$ the Lyapunov exponent
\[
\chi(t):= \lim_{n\cv\infty} \unsur{n} \int \log\norm{\gamma} \dmu_t^n (\gamma)
\]
is a well defined positive number, and that the non-Archimedean Lyapunov exponent
\[
\chi_\na := \lim_{n\cv\infty} \unsur{n} \int \log\norm{\gamma}_\na \dmu_\na^n(\gamma)
\]
is also well-defined and positive by Theorem~\ref{thm:positive}.

We have to show that
\begin{equation}\label{eq:lyapunov na2}
\unsur{ {\log\abs{t}\inv}}\chi(t) \longrightarrow \chi_{\sf na} \text{ as } t\cv 0.
\end{equation}

If $(K, \abs{{\,\cdot\,}})$ is any metrized field, and $z =[z_1:z_2]\in \PP^1(K)$, recall the notation
\[
\sigma(\gamma, z) = \log \frac{\norm{\gamma Z}}{\norm{Z}}~,
\]
where $Z=(z_1, z_2)\in K^2\setminus\set{0}$ and $\norm{Z} = \max(\abs{z_1}, \abs{z_2})$. To establish~\eqref{eq:lyapunov na2}, we first need to relate the classical and non-Archimedean expansion rates for a single group element. We start with the following consequence of Lemma~\ref{lem:uniform disk}.

\begin{lemm}\label{lem:sup2}
There exists a constant $C>0$, such that for any model $Y$, any $g\in G$, and any point $y_{t}\in Y_{t}$ with $|t|\le 1/2$, we have
\begin{equation}\label{eq:sup2}
\abs{\frac{1}{\log|t|^{-1}} \sigma(\gamma_t, y)}
\leq
\log \norm{\gamma}_\na + \frac{C \length (g)}{\log\abs{t}\inv}~,
\end{equation}
where $\gamma = \rho(g) \in \SL(2,\mm)$.
\end{lemm}

\begin{proof}
As before for $t\neq 0$ we can naturally identify the fibers $Y_t$ and $\set{t}\times \pu$, and we write $y_t = [z_{1t}:z_{2t}]$, with $Z_t = (z_{1t}, z_{2t})$ and $\max\{ |z_{1t}|, |z_{2t}|\} =1$. The upper bound
\[
\sigma(\gamma_t, y_{t}) = \log\norm{\gamma_t Z_t} \leq \log\norm{\gamma_t} \le
(\log\abs{t}\inv) \log\norm{\gamma}_{\sf na} + C \length (g)
\]
follows directly from the definitions and Lemma~\ref{lem:uniform disk}. To get the lower bound it is enough to write $ \norm{\gamma_t Z_t} \geq \norm{\gamma_t\inv}\inv \norm{Z_t}$ and remind that $\norm{\gamma_t} = \norm{\gamma_t\inv}$.
\end{proof}

The main step of the proof is the following proposition.

\begin{prop}\label{prop:hybrid}
For every model $Y$, there exists a constant $c(Y)>0$ satisfying the following property. For every $\gamma \in \SL(2, \mm)$, there exists a point $\alpha(\gamma)\in Y$ such that if $t_j\cv0$ and $(y_{t_j})\in Y_{t_j}$ is any sequence not accumulating $\alpha(\gamma)$, then
\begin{equation}\label{eq:norm}
\liminf_{j\cv\infty} \unsur{\log\abs{t_j}\inv} \sigma(\gamma_{t_j}, y_{t_j}) \geq \log\norm{\gamma}_{\sf na} - c(Y).
\end{equation}
\end{prop}

This says that the (positive) upper bound that was obtained in~\eqref{eq:sup2} is almost achieved everywhere on $Y$ when $t\to0$, except at one point, up to an error that is uniform in $\gamma$ (compare Lemma~\ref{lem:expand vector 1}).
%\medskip

Let us postpone the proof of the proposition to the end of the section, and first complete the proof of~\eqref{eq:lyapunov na2}.

Fix $\e>0$. Apply Corollary~\ref{cor:cut to small} to get a model $Y$ in which all the atoms of the residual measure $\nu_{Y_0}= (\red_Y)_* \nu_\na$ are smaller than $\e$. By Theorem~\ref{Thm:atomic}, we have that $\nu_{Y_t}\cv \nu_{Y_0}$ as $t\cv 0$ where $\nu_{Y_t}$ is the pull-back of $\nu_t$ to $Y$.

We first work with a fixed $g\in G$, and as usual we write $\gamma = \rho(g) \in \SL(2,\mm)$. By Lemma~\ref{lem:sup2}, we have that
\begin{equation}\label{eq:666}
\limsup_{t\cv0} \unsur{\log\abs{t }\inv} \int \sigma(\gamma_t, y) \dnu_t(y) \leq \log \norm{\gamma}_{\sf na}.
\end{equation}
To obtain a lower bound, we fix a small neighborhood $U$ of $\alpha(\gamma)$ in $Y$ such that $\nu_t(U) \le 2 \e$ for any $t$. This is possible because $\nu_t \to \nu_{Y_0}$ and $\nu_{Y_0}(\alpha(\gamma)) \le \e$. Then Proposition~\ref{prop:hybrid} shows that if $\eta \ll 1$ is fixed, then for every small enough $t$, for $y\in Y_t\setminus U$ we get
\[
\frac{1}{\log|t|^{-1}} \sigma(\gamma_t, y)
\ge \log \norm{\gamma}_\na - c(Y) - \eta,
\]
whereas for $y\in Y_t \cap U$, Lemma~\ref{lem:sup2} implies that
\[
\frac{1}{\log|t|^{-1}} \sigma(\gamma_t, y)
\ge - \log \norm{\gamma}_\na -\eta.
\]
Combining these two estimates we infer that
\[
\unsur{\log\abs{t }\inv} \int_{Y_t} \sigma(\gamma_t, y) \dnu_t(y) \geq
(1-2\e) (\log \norm{\gamma}_\na - c(Y) - \eta) - 2\e (\log \norm{\gamma}_\na +\eta)
,
\]
therefore since $\eta$ is arbitrary,
\[
\liminf_{t\cv0}\unsur{\log\abs{t }\inv} \int_{Y_t} \sigma(\gamma_t, y) \dnu_t(y) \geq \log \norm{\gamma}_\na - (1-2\e)c(Y) - 2\e \log \norm{\gamma}_\na.
\]
Using this inequality and~\eqref{eq:666} we finally obtain
\begin{equation}\label{eq:limsupinf}
\limsup_{t\cv0} \abs{ \unsur{\log\abs{t }\inv} \int_{Y_t} \sigma(\gamma_t, y) \dnu_t(y) - \log\norm{\gamma}_{\sf na}} \leq c(Y) + 2\e \log \norm{\gamma}_{\sf na}.
\end{equation}

To conclude the argument we integrate this estimate with respect to $g$. Fix an integer $n$ so large that
\[
\frac{c(Y)}{n} <\e \quad\text{ and}\quad \abs{\unsur{n} \int \log \norm{\gamma}_{\sf na} \dmu_\na^n(\gamma) - \chi_{\sf na}} <\e.
\]
We observe that the Furstenberg formula for the Lyapunov exponent iterated $n$ times and read in the model $Y$ is expressed as
\[
\chi(t) = \unsur{n} \int_{Y_t \times \Gamma} {\sigma(\gamma_t, y)} \dnu_t(y) \dmu_t^n(\gamma),
\]
so we can write
\begin{align*}
\abs{ \unsur{\log\abs{t }\inv} \chi(t) - \chi_{\sf na} } &= \abs{\unsur{n} \int_{Y_t \times \Gamma_t} \frac{\sigma(\gamma_t, y)}{{\log\abs{t }\inv}} \dnu_t(y) \dmu_t^n(\gamma) - \chi_{\sf na} }\\
&\leq \abs{\unsur{n} \int_{Y_t \times \Gamma_t} \frac{\sigma(\gamma_t, y)}{{\log\abs{t }\inv}} \dnu_t(y) \dmu_t^n(\gamma) - \unsur{n} \int \log \norm{\gamma}_{\sf na} \dmu_\na^n(\gamma) }+ \e \\
&=: \Delta(t) + \e~.
\end{align*}
By the moment assumption (A2), there exists a finite subset $G' \subset G$ such that
\begin{equation}\label{eq:G'}
\int_{G\setminus G'} \length(g) \dsm^n (g) \le \e \quad\text{ and}\quad
\frac1n \int_{\Gamma_\na \setminus \Gamma'_\na} \log\norm{\gamma}_\na \dmu_\na^n(\gamma) \le \e
\end{equation}
where $\Gamma'_\na := \rho_\na(G')$. To bound the quantity $\Delta(t)$ we split the integrals according to the decomposition $G = G'\cup G\setminus G'$. If $\Delta'(t)$ denotes the contribution coming from $G'$, using~\eqref{eq:G'} and Lemma~\ref{lem:sup2}, we get
\begin{align*}
\Delta(t) & \le \Delta'(t) + 2 \int_{\Gamma_\na \setminus \Gamma'_\na}\frac1n \log\norm{\gamma}_\na \dmu_\na^n(\gamma) + \frac{C}{\log \abs{ t}\inv }
\int_{G\setminus G'} \length(g) \dsm^n (g) \\
& \le \Delta'(t) + (C+2) \e~,
\end{align*}
From~\eqref{eq:limsupinf} we have that
\begin{align*}
\limsup_{t\to 0}\Delta'(t) &=
\limsup_{t\to 0}\abs{\unsur{n} \int_{Y_t \times \Gamma'_t} \frac{\sigma(\gamma_t, y)}{{\log\abs{t }\inv}} \dnu_t(y) \dmu_t^n(\gamma) -\int_{\Gamma' _\na} \unsur{n} \log\norm{\gamma}_\na \dmu_\na^n(\gamma) } \\
&\leq
\unsur{n} \int_{G'} \limsup_{t\to 0} \abs{ \int_{Y_t} \frac{\sigma(\rho_t(g), y)}{{\log\abs{t }\inv}} \dnu_t(y) - \log\norm{\rho_\na(g)} } \dsm^n(g)
\\
& \leq
\frac{c(Y)}n + \frac{ 2\e}{n} \int_{\Gamma'} \log \norm{\gamma}_\na \dmu_\na^n(\gamma)\leq \e + 2 \e (\chi_\na + \e)~,
\end{align*}
where in the second line we use $\mu_t = (\rho_t)_*\sm$ and $\mu_\na = (\rho_\na)_*\sm$. Finally we conclude that
\[
\limsup_{t\cv0} \abs{ \unsur{\log\abs{t }\inv} \chi(t) - \chi_{\sf na} } \leq \e + (C+2)\e + \e + 2 \e (\chi_\na + \e)
\]
Since this estimate makes no reference to the model $Y$ and $\e$ is arbitrary, the theorem follows.


\begin{proof}[Proof of Proposition~\ref{prop:hybrid}]
We first need a version of Proposition~\ref{prop:KAK} for $\SL(2, \mm)$. Indeed since $\mm$ is neither a field nor complete, this proposition cannot be applied directly. Let us explain how to adapt the argument to this concrete setting. We first introduce some notation: for $r>0$ denote by $\mathcal{O}_r$ (resp. $\mm_r$) the ring of holomorphic functions in $D(0, r)$ (resp. of holomorphic functions in $D(0, r)\setminus {0}$ admitting a meromorphic extension at the origin).

\begin{lemm}\label{lem:KAK holomorphe}
For every $\gamma\in \SL(2, \mm)$ there exist $r>0$, $m,n \in \SL(2,\mathcal{O}_r)$ and $a \in \SL(2, \mm_r)$ diagonal such that $\gamma= m\cdot a \cdot n$ and $\norm{a}_\na = \norm{\gamma}_\na$.
\end{lemm}

\begin{proof}
Observe first that a meromorphic family of matrices $(\gamma_t)$ in $\SL(2, \mm)$ extends holomorphically at the origin if and only if for any triple $\set{a,b,c}$ of distinct points in $\pu(\cc)$, then as $t\to 0$, there exist distinct $a',b',c'$ such that $\gamma_t(a)\to a'$, $\gamma_t(b)\to b'$, and $\gamma_t(c)\to c'$.

Let now $\gamma\in \SL(2, \mm)$ and assume that $\gamma\notin\SL(2, \mathcal{O})$. Then on $X = \dD\times \PP^1_\cc$, $\gamma$ contracts $X_0\setminus\set{\rep(\gamma_X)}$ to $\set{\att(\gamma_X)}$. Pick $m\in \SL(2, \mathcal{O})$ such that $m\inv(\att(\gamma_X)) = \infty$ and $n\in \SL(2, \mathcal{O})$ such that $n(\rep(\gamma_X)) = \infty$. Then $\gamma' = m\inv\gamma n\inv$ maps $X_0\setminus\set{0}$ to $0$, which implies that for small $t$, $\gamma_t'$ is loxodromic with an attracting fixed point $\att(\gamma_t)$ close to $\infty$ and a repelling fixed point $\rep(\gamma_t)$ close to $0$. Thus there exist $r>0$ and $h\in \SL(2, \mathcal{O}_r)$ such that $h_t(\att(\gamma_t)) =\infty$, $h_t(\rep(\gamma_t)) = 0$ and $h_t(1) =1$. Then $a= h\gamma'h^{-1}$ fixes 0 and $\infty$, so it is diagonal.

By the first observation, $t\mapsto h_t$ extends holomorphically at the origin, that is $h\in \SL(2, \mathcal{O}_r)$. So the desired decomposition is $\gamma = (h\inv m) a (n h)$. The equality $\norm{a}_\na = \norm{\gamma}_\na$ follows easily.
\end{proof}

We are now ready to prove Proposition~\ref{prop:hybrid}. We start by working on $X$. Pick a sequence of points $(x_{t_j})$ converging to the central fiber, and consider the quantities $\sigma(\gamma_{t_j}, x_{t_j}) = \norm{\gamma_{t_j}X_{t_j}}$. Extract so that $(x_{t_j})$ converges and drop the index $j$ for notational simplicity. If $m\in \SL(2, \mathcal{O})$ then for every $Z\in \mathcal O^2$, $\norm{ m_t Z_t}\asymp \norm{Z_t}$ so by the previous lemma we can assume that $\gamma$ is diagonal, $\gamma_t =\diag(\lambda_t, \lambda_t\inv)$. If $(x_{t})$ does not converge to $[0:1]$, then $\norm{\gamma_{t}X_{t}} \asymp \norm{\gamma_{t}}$ so the desired estimate holds, therefore the interesting case is when $(x_{t})$ converges to $[0:1]$. In this case a lift of norm 1 of $x_t$ will be of the form $X_{t} = (\xi_t, \eta_t)$ with $\abs{\eta_t} = 1$, so $\sigma (\gamma_t, x_t) = \max(\abs{\la_t\xi_t}, \abs{\la_t}\inv)$. From this formula we infer that if for some $l>0$, $\abs{\xi_t} \geq \abs{t}^l$ when $t\cv 0$ then
\[
\liminf_{t\to 0} \unsur{\log\abs{t}\inv} \sigma (\gamma_t, x_t) \geq \abs{\la}_\na -l = \norm{\gamma}_\na -l.
\]

We rely on the following elementary geometric fact.

\begin{lemm}\label{lem:geometric}
Let $\pi:M\cv \dD^2$ be a composition of $N$ blow-ups above the vertical fiber $\set{0}\times \dD$ in the unit bidisk, and denote $M_0 = \pi\inv(\set{0}\times \dD)$.
Then if $\ell>N$, the open set
\[
\pi\inv\lrpar{\set{(t,x)\in \dD^*\times \dD, \ \abs{x}<\abs{t}^\ell}}
\]
clusters at a unique point of $M_0$.
\end{lemm}
\begin{proof}
The open set $\set{(t,x)\in \dD^*\times \dD, \ \abs{x}<\abs{t}^\ell}$ is the union of the curves $\set{ x= c t^\ell}$ in $\dD^*\times \dD$, where $c$ ranges over $\abs{c}<1$. These curves get separated after exactly $\ell$ blow-ups.
\end{proof}

In order to conclude the proof, pick a model $\pi: Y\to X$ and a sequence $(y_{t_j})$ as in the statement of the proposition. Extract a subsequence so that $(y_{t_j})$ converges. Let $N$ be the number of blow-ups required to obtain $Y$. We put $x_{t_j} = \pi(y_{t_j})$ and do the analysis of the first part of the proof. Then, Lemma~\ref{lem:geometric} applied to $l = N+1$ provides a point $\alpha = \alpha(\gamma)$ in the central fiber $Y_0$ such that if $(y_{t_j})$ does not converge to $\alpha$, then
\[
\liminf_{j\cv\infty} \unsur{\log\abs{t_j}\inv} \sigma (\gamma_{t_j}, y_{t_j}) \geq \norm{\gamma}_\na -(N+1).
\]
The result follows.
\end{proof}


\section{Degenerations: elementary representations}\label{sec:elem}

In this section we complete the proof of Theorem~\ref{Thm:lyapunov na} by addressing the case of elementary representations. Let as before $G$ be a finitely generated group endowed with some probability measure $\sm$ satisfying
\begin{enumerate}\renewcommand*{\theenumi}{A\arabic{enumi}${}^+$}\setcounter{enumi}{1}
\item \label{A2+b} there exists $\delta>0$ such that $\int (\mathrm{length}(g))^{1+\delta} \dsm(g)<\infty$,
\end{enumerate}
and let $\rho :G\to \SL(2, \mm)$ be any meromorphic family of representations.

With notation as in Section~\ref{subs:atomic}, Viewing $\mm$ as a subring of $\Laurent$ we denote by $\rho_\na$ the corresponding non-Archimedean representation $G\to \SL(2, \Laurent)$ and $\mu_\na = (\rho_\na)_* \sm$, which satisfies the moment condition
\begin{enumerate}\renewcommand*{\theenumi}{B\arabic{enumi}${}^+$}\setcounter{enumi}{1}
\item \label{B2+b} there exists $\delta>0$ such that $\int \log \|\gamma\|^{1+\delta} \dmu(\gamma_\na)<\infty$
\end{enumerate}
in $\SL(2, \Laurent)$. In particular the non-Archimedean Lyapunov exponent $\chi_\na = \chi(\mu_\na)$ is well-defined. Likewise for $t\in \D^*$ we let $\mu_t = (\rho_t)_* \sm$ and $\chi(t) = \chi(\rho_t(G), \mu_t)$.

\begin{theo} \label{thm:degeneration elementary}
Let $(G, \sm)$ be a finitely generated group endowed with a probability measure satisfying~\eqref{A2+b}, and let $\rho: G\to \SL(2, \mm)$ be such that $ \rho_\na(G) \subset \SL(2,\Laurent)$ is elementary. Then
\begin{equation}\label{eq:lyapunov elementary}
\chi(t) = \big({\log\abs{t}\inv} \big) \chi_{\sf na} + O(1) \text{ as } t\cv 0.
\end{equation}
If in addition $\sm$ is symmetric, then $\chi_{\sf na} = 0$.
\end{theo}

Under mild assumptions, the error term can be understood more precisely, see Section~\ref{subs:continuity elementary}.

Set $\Gamma_\na = \rho_\na(G)$. According to the discussion in Section~\ref{subs:elementary}, if $\Gamma_\na \leq \SL(2, \Laurent)$ is elementary then it is either non-proximal or non-strongly irreducible, so there are three possibilities:
\begin{enumerate}
\item \label{case1} $\Gamma_\na $ has potential good reduction;
\item \label{case2} $\Gamma_\na $ is conjugate to a subgroup of the affine group $\set{z\mapsto az+b, a\in \Laurent^\times, b\in \Laurent}$;
\item \label{case3} $\Gamma_\na $ is conjugate to a subgroup of the group of transformations fixing $\set{0, \infty}$, that is, $\set{z\mapsto \lambda z^{\pm1}, \la\in \Laurent^\times}$.
\end{enumerate}

Note that if we are not in case~\eqref{case1}, then the projection of $\Gamma_\na$ in $\PGL(2, \Laurent)$ is not purely elliptic, and it follows from the analysis of Section~\ref{subs:elementary} that the conjugacy in~\eqref{case2} and~\eqref{case3} lies in $\SL(2, \Laurent)$ (i.e. no field extension is required).

In the remaining part of this section we split the proof of Theorem~\ref{thm:degeneration elementary} according to these three cases.

\subsection{Potential good reduction}

In case~\eqref{case1}, $\Gamma_\na $ is conjugate in $\SL(2, \cc(\!(t^{1/2})\!))$ to a representation fixing the Gau{\ss} point. Lifting to a branched 2-cover (which amounts to making the change of variables $t= u^2$), we can assume that the conjugacy lies in $\SL(2, \Laurent)$, that is there exists $\alpha\in \SL(2, \Laurent)$ such that for every $\gamma\in \Gamma_\na$, $\norm{\alpha\inv\gamma\alpha}\leq 1$.

Observe first that $\mm$ is dense in $\Laurent$ so that there exists a sequence $\alpha_n \in \SL(2,\mm)$ such that $\norm{\alpha - \alpha_n} \to0$. From the continuity of the matrix product, and the ultrametric property, for any sufficiently large integer $n$ we get $\norm{\alpha_n\inv\rho_{\na}(s)\alpha_n}\leq 1$ for all $s$ in a fixed finite set of generators of $G$. In particular, we have $\norm{\alpha_n\inv\rho_{\na}(g)\alpha_n}\leq 1$ for all $g \in G$ so that we may suppose that our original conjugacy $\alpha$ belongs to $\SL(2,\mm)$.

Since the Lyapunov exponent is insensitive to conjugacy, by replacing $\rho$ by\linebreak $\alpha\inv\rho({\,\cdot\,})\alpha$ we can assume that $\rho$ extends holomorphically at the origin. For every $t\neq 0$, by sub-additivity we have the bound $0\leq \chi(t) \leq \int \log \norm{\rho_t(g)} \dsm(g)$. Therefore applying Lemma~\ref{lem:uniform disk} and the moment condition we infer that $\chi(t) = O(1)$ as $t\to 0$. On the other hand, since $\Gamma_\na$ has good reduction, $\chi_\na$ vanishes, and we are done.

\subsection{Affine representations}

Let $({k},\abs{{\,\cdot\,}})$ be any complete valued field and consider any subgroup $\Gamma$ of $\SLk$ endowed with a measure $\mu$, such that the projection of $\Gamma$ in $\PGLk$ lies in the affine group $\Aff(k)$. An element $\gamma\in \Gamma$ can be written in matrix form as
\begin{equation}\label{eq:matrix}
\gamma = \lrpar{
\begin{matrix} \alpha & \beta \\ 0 &\alpha\inv
\end{matrix}},
\end{equation}
corresponding to the M\"obius transformation $\gamma(z) = az+b$, with $a= \alpha^2$ and $b = \beta\alpha$. Thus its norm is
\begin{equation}\label{eq:norm affine}
\norm{\gamma} =
\max\lrpar{ \abs{\alpha}, \abs{\beta}, \abs{\alpha\inv} }= \max\lrpar{\abs{a}^{1/2}, \abs{a}^{ -1/2}, \abs{ba^{-1/2}} }.
\end{equation}

\begin{prop}\label{prop:formula affine}
Let $(k,\abs{{\,\cdot\,}})$ be a complete valued field and let $\mu$ be a measure with countable support in $\SLk$, contained in the affine group, and satisfying~\eqref{B2+b}. Then with notation as above we have
\begin{equation}\label{eq:formula affine}
\chi(\mu) =
\abs{\int \log |\alpha (\gamma)| \, \dmu(\gamma)} = \frac12 \abs{\int \log |a(\gamma)|\, \dmu(\gamma)}\,.
\end{equation}
In particular if $\mu$ is symmetric, $\chi(\mu) = 0$.
\end{prop}

\begin{proof}
For any $\omega = (\gamma_n) \in \Omega$, we write $\gamma_n(z) = a_n(\omega) z+b_n(\omega)$ so that
\[
\ell_n(\omega) = \gamma_n \cdots \gamma_1 (z) = A_n(\omega) z + B_n(\omega) = a_n\cdots a_1 z + \sum_{j=1}^{n-1} a_n\cdots a_{j+1} b_j + b_n
~.
\]
By the law of large numbers (or equivalently the Birkhoff ergodic theorem) we have that
\begin{equation}\label{eq:large numbers}
\frac1n \log| A_n| \to \lambda :=
\int \log |a| \dmu~\text{ a.s.}
\end{equation}
Fix $\varepsilon >0$. For a.e. $\omega$, $ e^{n(\lambda - \varepsilon)} \le |A_n(\omega) |\le e^{n(\lambda + \varepsilon)}$ for large $n$. The moment condition (B2${}^+$) and Chebyshev's inequality yield $\mu \{ \abs{b}> e^{\e j} \} \le C j^{-1-\delta}$, so that by the Borel--Cantelli lemma we get that $b_n(\omega) \le e^{\varepsilon n}$ a.s. for large $n$.

At this point we split the proof into two cases according to the sign of $\lambda= \int \log |a| \dmu$.
Write
\[
\gamma_n\circ \cdots \circ \gamma_1 (z) = a_1\cdots a_n \lrpar{z+ \sum_{j=1}^{n} \frac{b_j}{a_1\cdots a_j}}
= A_n\lrpar{z+ \sum_{j=1}^n b_jA_j\inv}.
\]
If $\lambda>0$ we infer from~\eqref{eq:large numbers} that
a.s. the partial sums of the series $\sum_{j\geq 1} b_j A_j^{-1}$ are bounded, from which it follows that $\abs{B_n} = O(\abs{A_n})$. Therefore
\[
\lim_{n\to \infty} \unsur{n} \log \abs{A_n} = \lambda \quad\text{ and}\quad \limsup_{n\to\infty}
\lrpar{\frac1{n} \log |B_n| - \frac1{2n} \log |A_n| }\leq \frac{\lambda}{2}~.
\]
By~\eqref{eq:norm affine}, we have
\begin{equation}\label{eq:norm affine2}
\frac1n \log \| \ell_n(\omega)\| = \max \left\{ \frac1{2n} \log |A_n|, - \frac1{2n} \log |A_n|, \frac1{n} \log |B_n| - \frac1{2n} \log |A_n|\right\},
\end{equation}
so we conclude that $\chi(\mu) = {\lambda}/{2}$.

On the other hand, if $\lambda\leq 0$, then since almost surely for large $j$, $\abs{b_j}\leq e^{\e j}$ and $e^{(\lambda-\e) j}\leq \abs{A_j}\leq e^{(\lambda+\e) j}$ we deduce that
\[
\abs{ \sum_{j=1}^n b_jA_j\inv} = O \lrpar{e^{(-\lambda +2\e)n}} \text{ hence } \abs{B_n} = O\lrpar{ e^{3\e n}}.
\]
Thus from~\eqref{eq:norm affine2} we get that for large $n$,
\[
-\frac{\lambda}{2} -\e \leq \frac1n \log \| \ell_n(\omega)\| \leq -\frac{\lambda}{2} + 4\e~,
\]
and
$\chi(\mu) = - {\lambda}/{2}$, as required.
\end{proof}

\begin{proof}[Proof of Theorem~\ref{thm:degeneration elementary} in the affine case]
Under the assumptions of the theorem, assume that the projection of $\Gamma_\na$ into $\PGL(2,\Laurent)$ lies in the affine group. For any $g\in G$, we use the same notation as above, writing $\alpha(\rho(g))\in \mm$ for the upper diagonal term of $\rho(g)$, and $\beta(\rho(g))\in \mm$ for its upper right term. We get corresponding coefficients $\alpha(\rho_t(g))\in \cc^*$ (for fixed $t\neq 0$) and $\alpha(\rho_\na(g))\in \Laurent^*$.


By Lemma~\ref{lem:uniform disk} we have that $\alpha(\rho_t(g)) = t^{-\log \abs{\alpha(\rho_\na(g))}} \widetilde\alpha(\rho_t(g))$, where $ t\mapsto \widetilde\alpha(\rho_t(g))$ is holomorphic in $\dD$, and $\abs{\log\norm{\widetilde \alpha }_{L^\infty(\overline{\D}(0, 1/2))} }\leq C(\rho) \length(g)$. Hence
\begin{equation}\label{eq:affine error}
\int \log|\alpha(\rho_t(g))| \, \dsm(g) = (\log|t|^{-1}) \, \int \log \abs{\alpha(\rho_\na(g))} \, \dsm + \mathcal{E}(t)
\end{equation}
where
\begin{equation}\label{eq:affine error2}
|\mathcal{E}(t)| = \left| \int \log \abs{\widetilde\alpha(\rho_t(g))} \, \dsm \right|
\le C \int \length(g) \dsm <+\infty~.
\end{equation}
Therefore applying the formula of Proposition~\ref{prop:formula affine} to $k = \cc$ and $k= \Laurent$ we infer the desired estimate~\eqref{eq:lyapunov elementary} in the affine case.
\end{proof}


\subsection{Representations fixing \texorpdfstring{$\set{0, \infty}$}{set0, infty}}

Back to the general setting, consider now a subgroup $\Gamma\leq \SLk$ whose projection in $\PGL(2,k)$ fixes $\set{0, \infty}$. Then every matrix in $\Gamma$ is of the form
\[
\text{either }\gamma =
\begin{pmatrix} \alpha & 0 \\ 0 &\alpha\inv
\end{pmatrix} \text{ or }
\gamma =
\begin{pmatrix} 0 & -\alpha \\ \alpha \inv &0
\end{pmatrix}~,
\]
and $\| \gamma \| = \max \{|\alpha|, |\alpha|^{-1}\}$. As M\"obius transformations, we have $\gamma(z) = a z$ or $-a/z$ with $a = \alpha^2$, and $\| \gamma\| = \max \{|a|^{1/2}, |a|^{-1/2}\}$.

\begin{prop}\label{prop:formula 0infty}
Let $(k,\abs{{\,\cdot\,}})$ be a complete valued field and let $\mu$ be a measure with countable support in $\SLk$, satisfying the moment condition (B2). Suppose that any element in the support of $\mu$ leaves the pair $\{0,\infty\}$ invariant and at least one element permutes $0$ and $\infty$.

Then $\chi(\mu) = 0$.
\end{prop}

The last case of Theorem~\ref{thm:degeneration elementary} immediately follows, since in this case we have that $\chi(t) \equiv 0 = \chi_\na $.

\begin{proof}
It is more convenient here to use probabilistic language. We denote by $\sE({\,\cdot\,})$ the expectation of a random variable.
%\smallskip

In terms of M\"obius transformations, we are considering a random composition of maps of the form $\gamma_j(z) = \la_j z^{\e_j}$ where $\lambda_j\in k^\times$ and $\e_j \in \{\pm1\}$ are iid random variables.

Write $\ell_n(\omega) = (\gamma_n \circ \cdots \circ \gamma_1) = \Lambda_n z^{\mathcal{E}_n}$, and let $X_n = \log |\Lambda_n|$ and $x_n = \log|\lambda_n|$. A simple computation shows that $\mathcal{E}_n = \prod_{i=1}^{n} \e_i$ and
\begin{equation}\label{eq:Xn}
X_n = x_n + \e_n x_{n-1} + \e_n \e_{n-1} x_{n-2} + \ldots + \e_n\cdots \e_1 x_1
\end{equation}
Note that $(x_n)$ is a sequence of iid real random variables with $\sE(\abs{x_1})<\infty$. Kingman's theorem implies that the sequence $(X_n/n)$ converges a.s. We have to show that its limit is 0.

Let $(n_l)_{l\ge0}$ be the increasing sequence of random times where $\e_{n_l} = -1$, that is, $(n_l)$ is defined by $n_0=0$ and $n_{l+1} = \min\set{j> n_l, \ \e_j = -1}$. Since the $\e_n$ are iid and $\mu$ does not give full mass to the affine group, $(n_{l+1} - n_l)_{l\geq 0}$ is a sequence of iid random variables with a geometric distribution of non-zero parameter $p>0$ (which is the probability that $\gamma$ is not affine).

For $q\geq 1$ put $Y_q = \sum_{j=n_{q-1}}^{n_{q}-1} x_j$ with the convention that $x_0=0$. Observe that $(Y_q)$ forms a sequence of iid random variables with finite first moment, and such that
\[
\sE({Y_1}) = \sE({x_1})+ \sum_{j=1}^{\infty} j(1-p)^j p \sE({x_1}) = \unsur{p} \sE({x_1})~.
\]
It follows from~\eqref{eq:Xn} that for every $l\geq 0$, $X_{n_l-1} = \sum_{j=1}^l (-1)^{l-j} Y_j$.

Finally, let $(Z_{l})_{l\geq 1} =(Y_{2l-1}-Y_{2l})_{l\geq 1} $ which is a sequence of iid random variables with $\sE(\abs{Z_1})<\infty$ and $\sE({Z_1})=0$. Up to sign we have that
\[
X_{n_{2l}-1} = \pm \sum_{j=1}^l Z_j,
\]
thus from the strong law of large numbers we infer that $\unsur{l} X_{n_{2l}-1} \to 0$ a.s. as $l\to \infty$, hence since $l\le n_l$ the same holds for $\unsur{n_{2l}-1} X_{n_{2l}-1}$. The proof is complete.
\end{proof}

\subsection{Continuity of the error term}\label{subs:continuity elementary}

If $\sm$ is finitely supported, the proof of Theorem~\ref{thm:degeneration elementary} actually yields a finer estimate in~\eqref{eq:lyapunov elementary} of the form
\[
\chi(t) = \big({\log\abs{t}\inv} \big) \chi_{\sf na} + C + o(1) \text{ as } t\cv 0.
\]
Indeed:
\begin{itemize}
\item if $\Gamma_\na$ has potential good reduction, the proof reduces the situation to that of a holomorphic family of representations, in which case the result follows from the Furstenberg theory when $\rho_0$ is non elementary in $\SLC$ and from Bocker--Viana~\cite{bocker-viana} when $\rho_0$ is elementary (the finiteness assumption on $\sm$ is used here);
\item if $\Gamma_\na$ is affine we have to show that the error $\mathcal{E}(t)$ in~\eqref{eq:affine error} admits a limit when $t\to0$, which by virtue of~\eqref{eq:affine error2} and Lemma~\ref{lem:uniform disk} follows from the dominated convergence theorem;
\item finally in the case of representations fixing $\set{0, \infty}$ there is nothing to prove because $\chi(t)\equiv 0$.
\end{itemize}


\section{Degenerations: the hybrid approach}\label{sec:hybrid}

We propose an alternative approach to the analysis of the blow-up of the Lyapunov exponent, which is based on the hybrid space constructed by Berkovich and used by Boucksom and Jonsson in~\cite{boucksom-jonsson} and by the first author in~\cite{degeneration}. The introduction of this space allows us to make sense of the convergence of measures $\nu_t \to \nu_{\na}$ and leads to a proof of Theorem~\ref{Thm:hybrid}.

\subsection{The hybrid space}\label{sec:hybrid detail}

We start by briefly recalling the definition the hybrid space, referring to~\cite{boucksom-jonsson,degeneration} for more details.

Let $\cA$ be the subring of $\Laurent$ consisting of those series $f(t) = \sum_{n \in \Z} a_n t^n$ such that $\|f\|_{\hyb} < + \infty$, where
\[
\|f\|_{\hyb}:= \sum_{n= - \infty}^{+\infty} |a_n|_{\hyb}\, e^{-n}~, \quad\text{ and}\quad
\begin{cases}
|a|_{\hyb} = \max\{ |a|, 1\} \text{ if }a\in \cc^* \\ |0|_{\hyb} =0\,.
\end{cases}
\]
Observe that for any $f\in \cA$, the sum has only finitely many negative terms and the series defining $f$ converges in $\overline{\dD}^*_{1/e}$. Endowed with the hybrid norm $\|{\,\cdot\,}\|_{\hyb}$, $\cA$ is a Banach ring, and its Berkovich spectrum $\cD:= M_{\mathrm{ber}}(\cA)$ is defined as usual to be the space of multiplicative semi-norms $|{\,\cdot\,}|$ on $\cA$ such that $|{\,\cdot\,}| \le \|{\,\cdot\,}\|_{\hyb}$, endowed with the topology of pointwise convergence.

It turns out that $\cD$ is naturally a closed disk. To see this, introduce the map $\tau$ from the closed disk of radius $1/e$ to $\cD$ by the formula:
\begin{equation}\label{eqn:def hyb}
\begin{cases}
|f(\tau(0))| = e^{-\ord_{t=0}(f)};
\\
|f(\tau(t))| = |f(t)|^{\frac{-1}{\log\mathopen|t\mathclose|}} \text{ if } 0 < |t| \le 1/e
\end{cases}
\end{equation}
for any $f\in \cA$. One can show that this map is a homeomorphism, see e.g.~\cite[Prop.~1.1]{degeneration}. A note on terminology: an element $x \in \cD$ is a non-negative real valued function on $\cA$, nevertheless as already said it is customary to write $f \mapsto |f(x)| = |f|_x \in \R_+$ for the evaluation map.

%\smallskip
The hybrid affine line $\A^1_{\hyb}:= M_{\mathrm{ber}}(\cA[Z])$ is by definition the set of multiplicative semi-norms $|{\,\cdot\,}|$ on $\cA[Z]$ such that $|{\,\cdot\,}| \le \norm{\cdot}_{\hyb}$ on $\cA$. We endow it with the topology of the pointwise convergence which makes it locally compact. The restriction to $\cA$ of any semi-norm $x$ is a point $p_{\hyb}(x) \in \cD$, and the projection $p_{\hyb} : \A^1_{\hyb}\to \cD$ is a continuous surjective map. Given $x\in \A^1_{\hyb}$, according to the value (zero or non-zero) of $ \tau\inv \circ p_{\hyb}(x)$, the semi-norm $x$ will carry non-Archimedean or Archimedean information.

It follows from the Gelfand--Mazur theorem (see e.g. the proof of~\cite[Prop.~1.1]{degeneration}) that $M_{\mathrm{ber}}(\cc[Z]) \simeq \cc$, so if $t\neq 0$, the fiber $p_{\hyb}\inv(\tau(t))$ is homeomorphic to $\cc$. Furthermore, there exists a unique homeomorphism $\widetilde \psi: \overline{\D}^*_{1/e} \times \cc \to p_{\hyb}^{-1}(\tau (\overline{\D}^*_{1/e})) $ satisfying
\begin{equation}\label{eq:identify}
|g(\widetilde\psi(t,z))| = |g(t,z)|^{\frac{-1}{\log \mathopen| t \mathclose|}}
\end{equation}
for any $(t,z) \in \overline{\D}^*_{1/e} \times \cc $ and any $g\in \cA[Z]$. By construction, we have $ p_{\hyb} \circ \widetilde\psi=\tau \circ \pi$, where $\pi: \overline{\D}^*_{1/e} \times \cc \to\overline{\D}^*_{1/e} $ is the first projection. On the other hand, any semi-norm on $\Laurent[Z]$ can be restricted to $\cA[Z]$ which yields a canonical map $\widetilde\psi_{\na}: \A^{1,\an}_{\Laurent} \to p_{\hyb}\inv(\tau(0))$. This map is a homeomorphism since the completion of $\cA$ with respect to the $t$-adic norm is the field $\Laurent$.

%\smallskip

The hybrid space $(\A^1)^*_{\hyb}$ associated to the punctured affine line $(\A^1)^*$ is the Berkovich spectrum of $\cA[Z,Z^{-1}]$ and can be identified with an open subset of $\A^1_{\hyb}$ whose complement is the set $ g \mapsto |g(0)|$ with $|{\,\cdot\,}| \in \cD$. The latter set is the closure in $\A^1_{\hyb}$ of $\widetilde{\psi}\lrpar{\overline{\D}^*_{1/e} \times \{0\} }$.

%\smallskip

The hybrid projective line is constructed as the union of two copies of $\A^1_{\hyb}$ patched in the usual way. Specifically, the natural inclusions $ \cA[z_1] \to \cA[Z,Z^{-1}]$ and $\cA[z_2] \to \cA[Z,Z^{-1}]$ sending $z_1$ to $Z$, and $z_2$ to $Z^{-1}$ yield two open embeddings $\imath_1, \imath_2\colon(\A^1)^*_{\hyb} \to \A^1_{\hyb}$, and $\phyb$ is defined to be the union of $U_1:=M_{\mathrm{ber}}(\cA[z_1])$ and $U_2 :=M_{\mathrm{ber}}(\cA[z_2])$ glued together using the identification $\imath_1(x) = \imath_2(x)$ for any $x\in (\A^1)^*_{\hyb}$.

The inclusion $U_1 \subset \PP^1_{\hyb}$ yields an open and dense embedding of $\A^1_{\hyb}$ into $\PP^1_{\hyb}$, and the following proposition holds.
\begin{prop}[\cite{degeneration}]
The hybrid space $\PP^1_{\hyb}$ is compact, and there exists a homeomorphism $\psi: \overline{\D}^*_{1/e} \times \PP^1_\cc \to p_{\hyb}^{-1}(\tau (\overline{\D}^*_{1/e})) $ whose restriction to $\overline{\D}^*_{1/e} \times \cc $ is equal to $\widetilde \psi $. Likewise, there is a \emph{canonical} homeomorphism $\psi_\na: \PP^{1,\an}_{\Laurent} \to p_{\hyb}\inv(\tau(0))$ whose restriction to $\A^{1,\an}_{\Laurent}$ is equal to $\widetilde \psi_{\na}$.
\end{prop}

The homeomorphism $\psi_\na$ is canonical so that we will identify in the sequel the fiber $p_{\hyb}\inv(\tau(0))$ with $\PP^{1,\an}_{\Laurent}$.

\begin{rema}
In other words, the proposition asserts that there exists a topology on the disjoint union $\lrpar{\overline{\D}^*_{1/e} \times \PP^1_\cc } \bigsqcup \PP^{1,\an}_{\Laurent}$ such that the map defined by $\psi$ on $\overline{\D}^*_{1/e} \times \PP^1 $ and $\psi_\na$ on $\PP^{1,\an}_{\Laurent}$ is a homeomorphism onto $\phyb$.
\end{rema}

The group $\SL(2,\mm)$ is contained in $\SL(2,\Laurent)$ so it admits a natural action on $\PP^{1,\an}_{\Laurent}$ preserving the analytic structure on this space. It also acts by biholomorphisms on $ \D^*\times \PP^1_\cc $ commuting with the second projection. The next proposition shows that these two actions fit together nicely in the hybrid space. To ease notation we write $\psi_t(z) = \psi(z,t)$.

\begin{prop}\label{prop:action hybrid}
The group $\SL(2,\mm)$ admits a unique action by homeomorphisms on the hybrid space $\PP^1_{\hyb}$ which is compatible with its natural action on $\PP^{1,\an}_{\Laurent}$, and such that
\begin{equation}\label{eq:321}
\psi_t (\gamma_t \cdot z) = \gamma_t \cdot \psi_t(z)
\end{equation}
for any $\gamma\in \SL(2,\mm)$, and any $(t,z)\in \overline{\D}^*_{1/e} \times\PP^1_\cc $. In particular, for all $\gamma\in \SL(2,\mm)$ and $x\in \phyb$, we have $p_{\hyb} (\gamma \cdot x) = p_{\hyb} (x)$.
\end{prop}

\begin{proof}
We define an action of $\SL(2,\mm)$ on $\PP^1_{\hyb}$ by setting $\gamma \cdot x = \gamma_{\na} \cdot x$ when $x\in \PP^{1,\an}_{\Laurent}$, and such that~\eqref{eq:321} holds true. It is only necessary to check that this action is continuous which will follow from the very definition of the hybrid space.

Recall that $\PP^1_{\hyb}$ is the union of two copies $U_1$ and $U_2$ of $\A^1_{\hyb}$. We pick $\gamma = \lrpar{
\begin{smallmatrix} a & b \\ c & d
\end{smallmatrix}} \in \SL(2,\mm)$
and look at the action of $\gamma$ in the first chart in $U_1= M_{\mathrm{ber}}(\cA[z])$. Observe that for any $f\in \cA[z]$, $f(\frac{az+b}{cz+d})$ is the quotient of some $\widetilde{f}\in \cA[z]$ by a polynomial of the form $(cz+d)^N\in \mm[z]$ for some integer $N$. It follows that
\[
|f(\gamma \cdot x)| = \left|f\left(\frac{az+b}{cz+d}\right) (x)\right|
\]
depends continuously on $x$ on the open set $U:= \{x\in U_1, \, |(cz+d)|_x\neq 0\}$, so that $\gamma$ defines a continuous map from $U$ to $U_1$.

If now $f\in \cA[z^{-1}]$, then $f(\frac{az+b}{cz+d})$ is the quotient of an element $g\in \cA[z]$ by a polynomial of the form $(az+b)^N\in \mm[z]$ for some $N$, and we conclude similarly that $\gamma$ defines a continuous map from $U':= \{x\in U_1, \, |(az+b)|_x\neq 0\}$ to $U_2$.

Since $ad-bc =1$, the two open sets $U$ and $U'$ cover $U_1$, which completes the proof.
\end{proof}

Recall from~\eqref{eq:sigma1} the definition of the cocycle $\sigma$.

\begin{prop}\label{prop:hybrid sigma}
For any $\gamma\in \SL(2,\mm)$, the function defined by
\[
\sigma_{\hyb}(\gamma,x) :=
\begin{cases}
\sigma(\gamma_\na,\psi_\na\inv(x)) & \text{ if } x \in p_{\hyb}^{-1}(\tau((0)) = \PP^{1,\an}_{\Laurent}
\\
\frac{\sigma(\gamma_t,z)}{\log\mathopen|t\mathclose|^{-1}}& \text{ if } p_{\hyb}(x) \neq 0 \text{ and } \psi^{-1}(x) = (z,t)
\end{cases}
\]
is continuous on $\PP^1_{\hyb}$.
\end{prop}

\begin{proof}
It is enough to check that the restriction of $\sigma_{\hyb}(\gamma,{\cdot\,})$ to one of the two defining charts of the hybrid projective line is continuous, say on $U = M_{\mathrm{ber}}(\cA[Z])$. Since $\sigma_{\hyb}(g,{\cdot\,})$ is continuous on $\PP^{1,\an}_{\Laurent}$ and on $\psi (\overline{\D}^*_{1/e} \times \PP^1_\cc)$ separately, we only need to check that $\sigma_{\hyb}(g,x_i) \to \sigma_{\hyb}(g,x)$ for any net of points $x_i = \psi(t_i, z_i)$, $(t_i, z_i)\in {\D}^*_{1/e} \times \cc $ indexed by some inductive set $I$ and such that $x_i \to x \in \PP^{1,\an}_{\Laurent}$ along $I$. Note that this implies $t_i \to 0$ and $|g(x_i)| \to |g(x)|$ for any $g\in \cA[Z]$, i.e.
\[
|g_{t_i}(z_i)|^{\frac{-1}{\log\mathopen| t_i\mathclose|}} \to |g(x)|~.
\]
By switching the charts we are working in and extracting a subfamily if necessary, we may also suppose that $|Z(x)| \le 1$ and $|z_i| \le 1$ for all $I\in I$. Let $a,b,c$ and $d\in \mm$ be the coefficients of $g$. By definition, we have
\[
\sigma_{\hyb}(g,x_i) =
\frac{\log \max\{ |a(t_i) z_i +b(t_i)|, |c(t_i) z_i + d(t_i)|\}}{\log\mathopen|t_i\mathclose|^{-1}}~.
\]
On the other hand for any $h\in \cA[Z]$ we have $\lim_i |h(x_i)| = |h(x)|$. Since $\mm \subset \cA$,~\eqref{eq:identify} yields
\[
\lim_i |(a Z + b)(x_i)| =\lim_i |a(t_i)(z_i) + b(t_i)|^{\frac{-1}{\log\mathopen| t_i\mathclose|}} = |(aZ+b)(x)|
\]
which implies $\sigma_{\hyb}(g,x_i) \to \sigma_{\hyb}(g,x)$ as required.
\end{proof}


\subsection{Convergence of measures in the hybrid space: proof of Theorem~\ref{Thm:hybrid}}

Consider a representation $\rho: G \to \SL(2,\mm)$ such that $\rho_{\na}$ is non-elementary, and $\sm$ a measure on $G$ satisfying (A1).

By Theorem~\ref{thm:positive}, we may consider the unique stationary measure $\nu_{\na}$ associated to $\rho_{\na}$: this is a probability measure supported in $\PP^{1,\an}_{\Laurent} = p_{\hyb}^{-1}(\tau(0)) \subset \phyb$. By Lemma~\ref{lem:specialization NE}, for any small enough $t\neq0$ the representation $\rho_t$ is non-elementary, and we denote by $\nu_t$ the image in $ \set{t} \times \PP^1_\cc $ of the unique stationary measure associated to $\rho_t$ under the natural inclusion $\PP^1_\cc \subset \set{t} \times \PP^1_\cc $. We shall see that any limit point of $\nu_t$ is a stationary measure on $\PP^{1,\an}_{\Laurent}$ hence equal to $\nu_{\na}$ so that $\nu_t \to \nu_{\na}$ (here for simplicity we drop the mention to the embeddings $\psi$ and $\psi_\na$). Since the hybrid space is not metrizable, some care needs to be taken when arguing in this way, and we thus proceed as follows.

Consider the set of probability measures $M = \{ \psi_*(\nu_t), \, 0 < |t| \le 1/e\}$ in $\phyb$. Let $\overline{M}$ be the closure of $M$ in the space of probability measures, for the weak-$\star$ topology associated to the hybrid topology. Since $\PP^{1}_{\hyb}$ is compact, so does $\overline{M}$.

Let us prove that $\overline{M} \setminus M = \{(\psi_\na)_* \nu_{\na}\}$.

We claim that $\overline{M} \setminus M\neq \emptyset$. Indeed for any $\delta \in (0, 1/e)$, define $M_\delta =\{ \psi_*(\nu_t), \, 0 < |t| \le \delta\}$. This forms an increasing family of subsets of $M$. Observe that any measure in $M_\delta$ has its support included in $p_{\hyb}^{-1}(\tau(\overline{\D}^*_\delta))$. The intersection $\bigcap_{\delta>0} \overline{M_\delta}$ is non-empty as an intersection of compact sets, and it is included in $\overline{M} \setminus M$ since any measure in this intersection has its support included in $p_{\hyb}^{-1}(\tau(0))$. This proves our claim. This also proves that any measure in $\overline{M} \setminus M$ is supported on $p_{\hyb}^{-1}(\tau(0)) = \PP^{1,\an}_{\Laurent}$.

Pick now any measure $\nu \in \overline{M} \setminus M$. Let $\varphi_{\na} : \PP^{1,\an}_{\Laurent} \to \R$ be an arbitrary continuous function. Since $\phyb$ is compact, it is a normal topological space, so the Tietze--Urysohn extension lemma applies. We can thus find a continuous function $\phi: \phyb \to \R$ whose restriction to $\PP^{1,\an}_{\Laurent}$ is equal to $\varphi_{\na}\circ \psi_\na\inv$. Let us introduce the convolution operator acting on continuous functions of $\phyb$ by setting
\[
\sm \ast \phi (x) = \int_G \phi(\rho(g)^{-1} \cdot x) \dsm(g)~.
\]
Observe that for any $0 < |t| \le 1/e$, we have
\begin{align*}
\int_{\PP^1(\cc)} (\sm \ast \phi) \dd(\psi_*(\nu_t)) &=
\int_{\PP^1} \int_G \phi(\rho(g)^{-1} \cdot x) \dd(\psi_*(\nu_t)) \, \dsm(g)
\\ & = \int_{\PP^1} \int_G \phi(\rho(g) ^{-1} \cdot \psi(z,t)) \dnu_t \, \dsm(g)
\\
& \overset{\mathclap{\eqref{eq:321}}}{=}\, \int_{\PP^1} \int_G (\phi \circ \psi_t) (\rho_t(g)^{-1} \cdot z) \dnu_t
\\ &= \int_{\PP^1} (\phi \circ \psi_t) \dd((\rho_t)_*\sm) * \nu_t) = \int_{\PP^1} \phi \, \dd(\psi_{*}(\nu_t))
\end{align*}
so by definition of the weak-$\star$ topology we get that $\int (\sm \ast \phi) \dnu= \int \phi \dnu$. This implies that $(\rho_{\na*}(\sm)) \ast (\psi_\na^*\nu) = \psi_\na^*\nu$, hence $\nu = (\psi_\na)_*\nu_{\na}$ since $\rho_{\na}$ admits a unique stationary measure.

%\smallskip

Finally let us show that $\psi_*(\nu_t) \to(\psi_\na)_* \nu_{\na}$. We argue by contradiction, and pick $\e>0$, a continuous function $\phi$ on $\phyb$, a sequence $t_n\to0$ such that
\[
\int \phi \, \dd(\psi_*(\nu_{t_n})) \ge \int \phi\, \dd((\psi_\na)_*\nu_{\na}) + \e.
\]
Since $\nu_{\na}$ belongs to the accumulation set $\bigcap_m \overline{\bigcup_{n\ge m} \{\psi_*(\nu_{t_n})\}}$ of the sequence $\psi_*(\nu_{t_n})$, it follows that the open set $\{ \nu, \, \int \phi \, \dnu < \int \phi\, \dd((\psi_\na)_*\nu_{\na}) + \e\}$ contains infinitely many measures of the form $\psi_*(\nu_{t_n})$, which is contradictory, thereby finishing the proof.
%\qed

\subsection{The hybrid approach to Theorem~\ref{Thm:lyapunov na} for non-elementary representations}

Let $\e>0$ be any positive small real number. By condition (A2) there exists a finite subset $G'$ of $G$ such that $\int_{G\setminus G'} \length(g) \dsm(g) \le \e$. By Lemma~\ref{lem:sup2} we have
\begin{align*}
\left|
\int_{G\setminus G'} \! \int_{\PP^1_\cc} \frac{\sigma(\rho_t(g),v)}{\log|t|^{-1}} \dsm(g)\dnu_t(v)
\right| &\mathop{\le}\limits^{\eqref{eq:sup2}}\!
\int_{G\setminus G'}\!
\lrpar{ \log \|\rho(g)\|_{\na} + \frac{C}{\log\mathopen|t\mathclose|^{-1}}
\length(g) }\dsm(g)
\end{align*}
which is bounded by $ 2C\e$ for small enough $t$. Likewise from Lemma~\ref{lem:expand vector 1} we infer that
\[
\left|\int_{G\setminus G'} \! \int_{\PP^1_{\Laurent}} \sigma(\rho_\na(g), v) \dsm(g) \right|
\le
\left|\int_{G\setminus G'} \log \|\rho(g)\|_{\na}\right|
\le C\e~.
\]
So using the above and Furstenberg's formula for the Lyapunov exponent we get
\begin{multline*}
\left| \frac{\chi(t)}{\log\abs{t}\inv} - \chi_{\na} \right|\\
\begin{aligned}
&= \left|
\int_{G} \int_{\PP^1_\cc} \frac{\sigma(\rho_t(g),v)}{\log|t|^{-1}} \dsm(g) \dnu_t(v) -
\int_{G}\int_{\PP^{1,\an}_{\Laurent}} \sigma(\rho_\na(g), v) \dsm(g) \dnu_\na(v)
\right|
\\
& \le 2 C \e +
\left|
\int_{G'} \! \left(
\int_{\PP^1_\cc} \frac{\sigma(\rho_t(g),v)}{\log|t|^{-1}} \dnu_t(v) - \int_{\PP^{1,\an}_{\Laurent}} \sigma(\rho_\na(g), v) \dnu_\na(v)
\right) \dsm(g)
\right|.
\end{aligned}
\end{multline*}
Viewed in the hybrid space, the difference of integrals in the last line rewrites as
\[
\int_\phyb \frac{\sigma_{\hyb}(\rho(g),x)}{\log|t|^{-1}} \,\dd\!\lrpar{(\psi_t)_*\nu_t }(x) - \int_\phyb {\sigma_{\hyb}(\rho(g),x)} \,\dd\! \lrpar{(\psi_\na)_*\nu_\na}(x)
\,,
\]
so using the finiteness of $G'$, Proposition~\ref{prop:hybrid sigma} and Theorem~\ref{Thm:hybrid} we deduce that
\[
\limsup_{t\to0} \abs{\log\abs{t}\inv \chi(t) - \chi_{\na} }\le 2C \e,
\]
and we conclude
by letting $\e\to0$.
%\qed

\begin{thebibliography}{MMMMMM}
\bibitem[ACS83]{avron-craig-simon} Avron, Joseph; Craig, Walter; Simon, Barry Large coupling behaviour of the Lyapunov exponent for tight binding one-dimensional random systems, J. Phys. A, Math. Gen., Volume 16 (1983) no. 7, pp. 209-211 \href{https://ahl.centre-mersenne.org/item/AHL_2019__2__515_0/#avron-craig-simon}{Publisher reference}.
\bibitem[Alp87]{alperin} Alperin, Roger C. An elementary account of Selberg’s lemma, Enseign. Math., Volume 33 (1987) no. 3-4, pp. 269-273 \href{https://ahl.centre-mersenne.org/item/AHL_2019__2__515_0/#alperin}{Publisher reference}.
\bibitem[Bea95]{beardon} Beardon, Alan F. The geometry of discrete groups, Graduate Texts in Mathematics, 91, Springer, 1995, xii+337 pages (Corrected reprint of the 1983 original) \href{https://ahl.centre-mersenne.org/item/AHL_2019__2__515_0/#beardon}{Publisher reference}.
\bibitem[Ber90]{Ber} Berkovich, Vladimir G. Spectral theory and analytic geometry over non-Archimedean fields, Mathematical Surveys and Monographs, 33, American Mathematical Society, 1990, x+169 pages \href{https://ahl.centre-mersenne.org/item/AHL_2019__2__515_0/#Ber}{Publisher reference}.
\bibitem[Ber09]{berko} Berkovich, Vladimir G. A non-Archimedean interpretation of the weight zero subspaces of limit mixed Hodge structures, Algebra, arithmetic, and geometry: in honor of Yu. I. Manin. Vol. I (Progress in Mathematics), Volume 269, Birkhäuser, 2009, pp. 49-67 \href{https://ahl.centre-mersenne.org/item/AHL_2019__2__515_0/#berko}{Publisher reference}.
\bibitem[BJ17]{boucksom-jonsson} Boucksom, Sébastien; Jonsson, Mattias Tropical and non-Archimedean limits of degenerating families of volume forms, J. Éc. Polytech., Math., Volume 4 (2017), pp. 87-139 \href{https://ahl.centre-mersenne.org/item/AHL_2019__2__515_0/#boucksom-jonsson}{Publisher reference}.
\bibitem[BL85]{bougerol-lacroix} Bougerol, Philippe; Lacroix, Jean Products of random matrices with applications to Schrödinger operators, Progress in Probability and Statistics, 8, Birkhäuser, 1985, xii+283 pages \href{https://ahl.centre-mersenne.org/item/AHL_2019__2__515_0/#bougerol-lacroix}{Publisher reference}.
\bibitem[BNV17]{bocker-viana} Bocker-Neto, Carlos; Viana, Marcelo Continuity of Lyapunov exponents for random two-dimensional matrices, Ergodic Theory Dyn. Syst., Volume 37 (2017) no. 5, pp. 1413-1442 \href{https://ahl.centre-mersenne.org/item/AHL_2019__2__515_0/#bocker-viana}{Publisher reference}.
\bibitem[BR10]{baker-rumely} Baker, Matthew; Rumely, Robert Potential theory and dynamics on the Berkovich projective line, Mathematical Surveys and Monographs, 159, American Mathematical Society, 2010, xxxiv+428 pages \href{https://ahl.centre-mersenne.org/item/AHL_2019__2__515_0/#baker-rumely}{Publisher reference}.
\bibitem[CM87]{culler-morgan} Culler, Marc; Morgan, John W. Group actions on R-trees, Proc. Lond. Math. Soc., Volume 55 (1987) no. 3, pp. 571-604 \href{https://ahl.centre-mersenne.org/item/AHL_2019__2__515_0/#culler-morgan}{Publisher reference}.
\bibitem[CS83]{culler-shalen} Culler, Marc; Shalen, Peter B. Varieties of group representations and splittings of 3-manifolds, Ann. Math., Volume 117 (1983) no. 1, pp. 109-146 \href{https://ahl.centre-mersenne.org/item/AHL_2019__2__515_0/#culler-shalen}{Publisher reference}.
\bibitem[DD12]{kleinbif} Deroin, Bertrand; Dujardin, Romain Random walks, Kleinian groups, and bifurcation currents, Invent. Math., Volume 190 (2012) no. 1, pp. 57-118 \href{https://ahl.centre-mersenne.org/item/AHL_2019__2__515_0/#kleinbif}{Publisher reference}.
\bibitem[DD15]{bers1} Deroin, Bertrand; Dujardin, Romain Lyapunov exponents for surface group representations, Commun. Math. Phys., Volume 340 (2015) no. 2, pp. 433-469 \href{https://ahl.centre-mersenne.org/item/AHL_2019__2__515_0/#bers1}{Publisher reference}.
\bibitem[DeM03]{demarco03} DeMarco, Laura Dynamics of rational maps: Lyapunov exponents, bifurcations, and capacity, Math. Ann., Volume 326 (2003) no. 1, pp. 43-73 \href{https://ahl.centre-mersenne.org/item/AHL_2019__2__515_0/#demarco03}{Publisher reference}.
\bibitem[DeM16]{demarco16} DeMarco, Laura Bifurcations, intersections, and heights, Algebra Number Theory, Volume 10 (2016) no. 5, pp. 1031-1056 \href{https://ahl.centre-mersenne.org/item/AHL_2019__2__515_0/#demarco16}{Publisher reference}.
\bibitem[DF16]{demarco-faber-2} DeMarco, Laura; Faber, Xander Degenerations of complex dynamical systems II: analytic and algebraic stability, Math. Ann., Volume 365 (2016) no. 3-4, pp. 1669-1699 (With an appendix by Jan Kiwi) \href{https://ahl.centre-mersenne.org/item/AHL_2019__2__515_0/#demarco-faber-2}{Publisher reference}.
\bibitem[DMF14]{demarco-faber} De Marco, Laura; Faber, Xander Degenerations of complex dynamical systems, Forum Math. Sigma, Volume 2 (2014), e6, 36 pages \href{https://ahl.centre-mersenne.org/item/AHL_2019__2__515_0/#demarco-faber}{Publisher reference}.
\bibitem[Fan18]{NL} Fantini, Lorenzo Normalized Berkovich spaces and surface singularities, Trans. Am. Math. Soc., Volume 370 (2018) no. 11, pp. 7815-7859 \href{https://ahl.centre-mersenne.org/item/AHL_2019__2__515_0/#NL}{Publisher reference}.
\bibitem[Fav19]{degeneration} Favre, Charles Degeneration of endomorphisms of the complex projective space in the hybrid space. (2019) (to appear in J. Inst. Math. Jussieu) \href{https://ahl.centre-mersenne.org/item/AHL_2019__2__515_0/#degeneration}{Publisher reference}.
\bibitem[FJ04]{valtree} Favre, Charles; Jonsson, Mattias The valuative tree, Lecture Notes in Mathematics, 1853, Springer, 2004, xiv+234 pages \href{https://ahl.centre-mersenne.org/item/AHL_2019__2__515_0/#valtree}{Publisher reference}.
\bibitem[FK83]{furstenberg-kifer} Furstenberg, Harry; Kifer, Yuri Random matrix products and measures on projective spaces, Isr. J. Math., Volume 46 (1983) no. 1-2, pp. 12-32 \href{https://ahl.centre-mersenne.org/item/AHL_2019__2__515_0/#furstenberg-kifer}{Publisher reference}.
\bibitem[Fur63]{furstenberg} Furstenberg, Harry Noncommuting random products, Trans. Am. Math. Soc., Volume 108 (1963), pp. 377-428 \href{https://ahl.centre-mersenne.org/item/AHL_2019__2__515_0/#furstenberg}{Publisher reference}.
\bibitem[Fur02]{furman} Furman, Alex Random walks on groups and random transformations, Handbook of dynamical systems, Vol. 1A, North-Holland, 2002, pp. 931-1014 \href{https://ahl.centre-mersenne.org/item/AHL_2019__2__515_0/#furman}{Publisher reference}.
\bibitem[GR85]{guivarch-raugi} Guivarc’h, Yves; Raugi, Albert Frontière de Furstenberg, propriétés de contraction et théorèmes de convergence, Z. Wahrscheinlichkeitstheor. Verw. Geb., Volume 69 (1985) no. 2, pp. 187-242 \href{https://ahl.centre-mersenne.org/item/AHL_2019__2__515_0/#guivarch-raugi}{Publisher reference}.
\bibitem[Gui90]{guivarch} Guivarc’h, Yves Produits de matrices aléatoires et applications aux propriétés géométriques des sous-groupes du groupe linéaire, Ergodic Theory Dyn. Syst., Volume 10 (1990) no. 3, pp. 483-512 \href{https://ahl.centre-mersenne.org/item/AHL_2019__2__515_0/#guivarch}{Publisher reference}.
\bibitem[Jon15]{jonsson} Jonsson, Mattias Dynamics of Berkovich spaces in low dimensions, Berkovich spaces and applications (Lecture Notes in Mathematics), Volume 2119, Springer, 2015, pp. 205-366 \href{https://ahl.centre-mersenne.org/item/AHL_2019__2__515_0/#jonsson}{Publisher reference}.
\bibitem[Kap09]{kapovich} Kapovich, Michael Hyperbolic manifolds and discrete groups, Modern Birkhäuser Classics, Birkhäuser, 2009, xxviii+467 pages (Reprint of the 2001 edition) \href{https://ahl.centre-mersenne.org/item/AHL_2019__2__515_0/#kapovich}{Publisher reference}.
\bibitem[LP89]{lepage-holder} Le Page, Émile Régularité du plus grand exposant caractéristique des produits de matrices aléatoires indépendantes et applications, Ann. Inst. Henri Poincaré, Probab. Stat., Volume 25 (1989) no. 2, pp. 109-142 \href{https://ahl.centre-mersenne.org/item/AHL_2019__2__515_0/#lepage-holder}{Publisher reference}.
\bibitem[MT18]{maher-tiozzo} Maher, Joseph; Tiozzo, Giulio Random walks on weakly hyperbolic groups, J. Reine Angew. Math., Volume 742 (2018), pp. 187-239 \href{https://ahl.centre-mersenne.org/item/AHL_2019__2__515_0/#maher-tiozzo}{Publisher reference}.
\bibitem[Ota15]{otal} Otal, Jean-Pierre Compactification of spaces of representations after Culler, Morgan and Shalen, Berkovich spaces and applications (Lecture Notes in Mathematics), Volume 2119, Springer, 2015, pp. 367-413 \href{https://ahl.centre-mersenne.org/item/AHL_2019__2__515_0/#otal}{Publisher reference}.
\bibitem[Poi13]{poineau} Poineau, Jérôme Les espaces de Berkovich sont angéliques, Bull. Soc. Math. Fr., Volume 141 (2013) no. 2, pp. 267-297 \href{https://ahl.centre-mersenne.org/item/AHL_2019__2__515_0/#poineau}{Publisher reference}.
\bibitem[Tem15]{temkin} Temkin, Michael Introduction to Berkovich analytic spaces, Berkovich spaces and applications (Lecture Notes in Mathematics), Volume 2119, Springer, 2015, pp. 3-66 \href{https://ahl.centre-mersenne.org/item/AHL_2019__2__515_0/#temkin}{Publisher reference}.
\bibitem[YW15]{yang-wang} Yang, Jinghua; Wang, Yuefei On p-adic Möbius maps (2015) (\url{https://arxiv.org/abs/1512.01305}) \href{https://ahl.centre-mersenne.org/item/AHL_2019__2__515_0/#yang-wang}{Publisher reference}.
\end{thebibliography}
\end{document}
